\documentclass[12pt, a4paper, final]{article}

% ============================================================
% Base packages (mirrors template main.tex)
% ============================================================
\usepackage[utf8]{inputenc}
\usepackage[T1]{fontenc}
\usepackage[english]{babel}
\usepackage{amsmath, amssymb, amsthm}

% ============================================================
% Load custom configuration (template schema)
% ============================================================
\input{./configuration/packages}
\input{./configuration/functions}
\input{./configuration/code-style}
\input{./configuration/style}
\newenvironment{comentario}{%
    \color{blue}
}{%
    \par\medskip
}

% ============================================================
% QuTu project — physics macros
% ============================================================
\newcommand{\hb}{\hbar}
\newcommand{\alp}{\alpha}
\newcommand{\eps}{\varepsilon}
\newcommand{\Vb}{V_{\!b}}
\newcommand{\xe}{x_e}
\newcommand{\half}{\tfrac{1}{2}}
\newcommand{\fourth}{\tfrac{1}{4}}
\newcommand{\minn}{\min(n,m)}
\newcommand{\Phik}{\Phi_k}
\newcommand{\trec}{t_{\text{rec}}}
\newcommand{\ttun}{T_{\text{tun}}}
\newcommand{\Delt}{\Delta}
\newcommand{\Omg}{\Omega}
\newcommand{\adag}{\hat{a}^{\dagger}}
\newcommand{\xL}{x_{\!L}}
\newcommand{\xR}{x_{\!R}}
\newcommand{\Aij}{A_{ij}}
\newcommand{\Xij}{X_{ij}}
\newcommand{\Qij}{Q_{ij}}

% Theorem environments
\theoremstyle{definition}
\newtheorem{result}{Result}[section]
\newtheorem{remark}{Remark}[section]

\newcommand{\ttitle}{Theoretical Framework for Quantum Tunneling in Double-Well Potentials}
\newcommand{\stitle}{\large{Basis Set Methods, Wavepacket Dynamics, and Physical Observables}}
\newcommand{\finaldate}{\today}
\newcommand{\authorname}{José Antonio \textsc{Quiñonero Gris}}
\newcommand{\univname}{Universidad de Murcia}
\newcommand{\worktype}{Theory Reference}
\newcommand{\supname}{José \textsc{Zúñiga Román}}
\newcommand{\degreename}{Química}
\newcommand{\deptname}{Departamento de Química Física}
\newcommand{\facultyname}{Facultad de Química}
\newcommand{\institutename}{QuTu Project}
\newcommand{\centername}{Universidad de Murcia}
\newcommand{\kkeywords}{quantum tunneling, double-well potential, harmonic oscillator basis, variational method, wavepacket dynamics, QuTu}
\newcommand{\HRule}{\rule{.9\linewidth}{.6pt}}

% \input{./configuration/embed_video}   % Uncomment to enable PDF video embedding
\usepackage{./configuration/mybraket}

% ============================================================
% Document-level settings
% ============================================================
\graphicspath{{./figures/}}

\AtBeginDocument{%
  \hypersetup{
    pdfauthor   = {\authorname},
    pdftitle    = {\ttitle},
    pdfsubject  = {Quantum Mechanics, Variational Method, Tunneling},
    pdfkeywords = {\kkeywords},
  }%
}

\author{\authorname}
\title{
  \vspace{-1cm}
  \textbf{\ttitle}\\[0.8em]
  \large\stitle\\[0.4em]
  \normalsize A Graduate-Level Theory Reference for the QuTu Project
}
\date{\today}

% ============================================================
\begin{document}
% ============================================================

%
% --- Title ---
%
% Option 1: simple title block (default)
%
\maketitle
%
% Option 2: minimal info block — same page
%
% \thispagestyle{empty}
% \input{front-page/minimal}
%
% Option 3: minimal separate title page
%
% \thispagestyle{empty}
% \input{front-page/minimal-titlepage}\clearpage
%
% Option 4: elaborated title page
%
% \input{front-page/elaborated-titlepage}\afterpage{\blankpage}

% ============================================================
%                       Abstract
% ============================================================
\begin{abstract}
This document presents the complete analytical framework underlying the
\textbf{QuTu} quantum tunneling simulation code. We treat the
one-dimensional double-well problem motivated by the ammonia inversion
mode. Five basis-set representations are developed in full detail:
(I)~harmonic oscillator functions centered at the origin,
(II)~shifted harmonic oscillator functions forming a two-center basis,
(III)~Gaussian basis functions, (IV)~the Fourier Grid Hamiltonian
(DVR), and (V)~the particle-in-a-box finite basis representation (PIB-FBR).
For each basis, all Hamiltonian matrix elements are derived
analytically with complete intermediate algebra. The theory of
wavepacket construction (two-state, four-state, and Gaussian initial
states), time evolution, and the computation of physical observables
(survival probability, position expectation value, tunneling splitting,
turning points) is presented in a self-contained manner.
For the asymmetric double well we analyse in detail the construction of
localized states via the two-level effective Hamiltonian and the
mixing-angle formalism, the sign convention for selecting the initial
well, the contrast between the parity-based localization mechanism of
the symmetric case and the two-dimensional eigenstate rotation of the
asymmetric case, and the fundamental complementarity between
localization quality and tunneling efficiency~$\eta = \sin^2(2\theta)$.
An analytical comparison of basis sets for the asymmetric case
establishes the practical advantages of the DVR over the shifted-HO
basis when parity symmetry is broken: the DVR eliminates all integral
computation through the Fourier Grid Hamiltonian representation,
provides a standard (rather than generalized) eigenvalue problem, and
is insensitive to changes in the asymmetry parameter.
The general case of arbitrary polynomial potentials and custom
wavepackets is also treated. Throughout, atomic units
($\hb = m_e = e = a_0 = 1$) are employed unless stated otherwise.
\end{abstract}

\tableofcontents
\newpage

% ============================================================
%                       Sections
% ============================================================
\section{The Double-Well Potential}\label{sec:potential}
% ============================================================
% ============================================================

% ------------------------------------------------------------
\subsection{Symmetric double-well}\label{sec:V_sym}
% ------------------------------------------------------------

The symmetric quartic double-well potential is written in the standard form
\begin{equation}\label{eq:V_ab}
  V(x) = a\,x^4 - b\,x^2, \qquad a > 0,\; b > 0.
\end{equation}
This potential is symmetric under parity: $V(-x) = V(x)$.
It has a local maximum (barrier) at $x = 0$ with $V(0) = 0$, and two degenerate minima.

\subsubsection{Location of the minima}

Setting the first derivative to zero:
\begin{equation}
  V'(x) = 4a\,x^3 - 2b\,x = 2x(2a\,x^2 - b) = 0.
\end{equation}
The solutions are $x = 0$ (local maximum) and
\begin{equation}
  2a\,x^2 = b
  \quad\Longrightarrow\quad
  \boxed{x_{\min} = \pm\sqrt{\frac{b}{2a}}\,.}
\end{equation}
We verify that these are minima by computing the second derivative:
\begin{equation}
  V''(x) = 12a\,x^2 - 2b.
\end{equation}
At $x = 0$: $V''(0) = -2b < 0$ (maximum, confirmed).
At $x_{\min}$: $V''(x_{\min}) = 12a\cdot\frac{b}{2a} - 2b = 6b - 2b = 4b > 0$ (minimum, confirmed).

\subsubsection{Barrier height}

The barrier height is the depth of the minima relative to the maximum at the origin:
\begin{equation}
  \Vb = V(0) - V(x_{\min})
  = 0 - \left[a\left(\frac{b}{2a}\right)^2 - b\left(\frac{b}{2a}\right)\right]
  = -\frac{b^2}{4a} + \frac{b^2}{2a}
  = \frac{b^2}{4a}\,.
\end{equation}
Therefore,
\begin{equation}
  \boxed{\Vb = \frac{b^2}{4a}\,.}
\end{equation}

\subsubsection{Local harmonic frequency at the minima}

Near each minimum, the potential is approximately harmonic with angular frequency
\begin{equation}
  \omega_{\text{local}} = \sqrt{\frac{V''(x_{\min})}{m}} = \sqrt{\frac{4b}{m}}\,,
\end{equation}
where $m$ is the mass of the particle. In atomic units ($\hb = 1$), the harmonic oscillator width parameter associated with this local frequency is
\begin{equation}\label{eq:alpha_local}
  \alpha_{\text{local}} = \frac{m\,\omega_{\text{local}}}{\hb} = m\,\omega_{\text{local}} = \sqrt{4\,m\,b} = 2\sqrt{m\,b}\,.
\end{equation}
For the reduced mass $\mu$ of the NH$_3$ inversion mode, $m \to \mu$, and
\begin{equation}
  \boxed{\omega_{\text{local}} = \sqrt{\frac{4b}{\mu}}\,,\qquad
  \alpha_{\text{local}} = 2\sqrt{\mu\,b}\,.}
\end{equation}

\subsubsection{Alternative parameterization in terms of $\xe$ and $\Vb$}

It is often convenient to parameterize the potential directly in terms of the
equilibrium distance $\xe \equiv |x_{\min}|$ and the barrier height $\Vb$.
From the results above:
\begin{equation}
  \xe = \sqrt{\frac{b}{2a}}\,,\qquad \Vb = \frac{b^2}{4a}\,.
\end{equation}
Solving for $a$ and $b$:
\begin{align}
  \text{From }\xe^2 &= \frac{b}{2a} \implies b = 2a\,\xe^2\,,\\
  \text{Substitute into }\Vb &= \frac{(2a\,\xe^2)^2}{4a} = a\,\xe^4\,,
\end{align}
yielding
\begin{equation}\label{eq:ab_from_xeVb}
  \boxed{a = \frac{\Vb}{\xe^4}\,,\qquad b = \frac{2\Vb}{\xe^2}\,.}
\end{equation}
The potential becomes
\begin{equation}\label{eq:V_xeVb}
  V(x) = \frac{\Vb}{\xe^4}\,x^4 - \frac{2\Vb}{\xe^2}\,x^2
  = \frac{\Vb}{\xe^4}(x^2 - \xe^2)^2 - \Vb\,.
\end{equation}
In this form, $V(\pm\xe) = -\Vb$ (the minima) and $V(0) = 0$ (the barrier top).
Alternatively, one may shift the energy zero so that $V(\pm\xe) = 0$:
\begin{equation}\label{eq:V_xeVb_shifted}
  V(x) = \frac{\Vb}{\xe^4}(x^2 - \xe^2)^2\,.
\end{equation}

\subsubsection{NH$_3$-specific parameterization}

For the ammonia inversion mode, the physical input consists of:
\begin{itemize}
  \item Equilibrium distance $\xe$ (in \AA{}ngstr\"om), to be converted to atomic units: $\xe^{\text{(a.u.)}} = \xe^{(\text{\AA})}/a_0$ with $a_0 = 0.529177\,$\AA.
  \item Barrier height $\Vb$ (in cm$^{-1}$), to be converted to atomic units: $\Vb^{\text{(a.u.)}} = \Vb^{(\text{cm}^{-1})} / 219474.63\;E_h/\text{cm}^{-1}$.
\end{itemize}
Then $a$ and $b$ are computed from \cref{eq:ab_from_xeVb} in atomic units.

The local harmonic frequency at the minima in cm$^{-1}$ is
\begin{equation}
  \nu_{\text{local}} = \frac{\omega_{\text{local}}}{2\pi c}
  = \frac{1}{2\pi c}\sqrt{\frac{4b}{\mu}}\,,
\end{equation}
where $c$ is the speed of light.

\subsubsection{Reduced mass for NH$_3$ inversion}

The umbrella inversion mode of NH$_3$ involves the nitrogen atom moving along the
$C_3$ symmetry axis while the three hydrogen atoms move in the opposite direction.
In the one-dimensional model, the inversion coordinate $x$ measures the displacement
of the nitrogen from the plane of the three hydrogens. The reduced mass for this motion is
\begin{equation}
  \boxed{\mu = \frac{3\,m_H\,m_N}{3\,m_H + m_N}\,.}
\end{equation}
\textbf{Derivation.} Let $M_H = 3m_H$ be the total mass of the hydrogen subsystem.
In the center-of-mass frame, the reduced mass for the relative motion of the nitrogen
atom against the hydrogen plane is
\begin{equation}
  \mu = \frac{M_H \cdot m_N}{M_H + m_N} = \frac{3\,m_H\,m_N}{3\,m_H + m_N}\,.
\end{equation}
Using $m_H = 1.00782503\,\text{u}$ and $m_N = 14.0030740\,\text{u}$
(where $1\,\text{u} = 1822.888\,m_e$):
\begin{equation}
  \mu = \frac{3(1.00782503)(14.0030740)}{3(1.00782503) + 14.0030740}\,\text{u}
  = \frac{42.3488}{17.0265}\,\text{u}
  = 2.4871\,\text{u}
  = 4534.2\,m_e\,.
\end{equation}

% ------------------------------------------------------------
\subsection{Asymmetric double-well}\label{sec:V_asym}
% ------------------------------------------------------------

The asymmetric double-well potential adds a linear bias term:
\begin{equation}\label{eq:V_asym}
  V(x) = a\,x^4 - b\,x^2 + c\,x, \qquad c \neq 0.
\end{equation}
The linear term $cx$ breaks the parity symmetry $V(-x) \neq V(x)$.

\subsubsection{Shift of minima: perturbative treatment}

For $|c| \ll b^2/\sqrt{b/(2a)}$ (weak asymmetry), we treat $cx$ as a perturbation.
Let the unperturbed minima be at $x_0 = \pm\sqrt{b/(2a)} \equiv \pm\xe$.
Expanding around $x_0$, the shifted minimum is at $x_0 + \delta x$ where
\begin{equation}
  V'(x_0 + \delta x) = 0 \implies
  \underbrace{V'(x_0)}_{=0} + V''(x_0)\,\delta x + c = 0
  \implies \delta x = -\frac{c}{V''(x_0)}\,.
\end{equation}
Since $V'(x) = 4ax^3 - 2bx + c$ and $V''(x_0) = 12a\,x_0^2 - 2b = 4b$:
\begin{equation}
  \delta x = -\frac{c}{4b}\,.
\end{equation}
The perturbed minima are thus
\begin{equation}\label{eq:xmin_pert}
  x_{\pm} \approx \pm\xe - \frac{c}{4b} + \mathcal{O}(c^2)\,.
\end{equation}
Both minima shift by the same amount $-c/(4b)$ in the $x$-direction.
The energy difference between the two wells to first order in $c$ is
\begin{equation}
  V(x_-) - V(x_+) \approx c\,x_- - c\,x_+
  = c(-\xe - \xe) = -2c\,\xe\,,
\end{equation}
so $\Delta V \approx -2c\,\xe$. The well at $x = +\xe$ is deeper when $c < 0$.

\subsubsection{Shift of minima: exact treatment}

The exact condition for the minima is
\begin{equation}\label{eq:cubic_roots}
  V'(x) = 4a\,x^3 - 2b\,x + c = 0\,.
\end{equation}
This is a depressed cubic equation in $x$. Dividing by $4a$:
\begin{equation}
  x^3 - \frac{b}{2a}\,x + \frac{c}{4a} = 0\,.
\end{equation}
This has the form $x^3 + px + q = 0$ with $p = -b/(2a)$ and $q = c/(4a)$.
The discriminant is
\begin{equation}
  \Delta_3 = -4p^3 - 27q^2 = 4\left(\frac{b}{2a}\right)^3 - 27\left(\frac{c}{4a}\right)^2
  = \frac{b^3}{2a^3} - \frac{27c^2}{16a^2}\,.
\end{equation}
When $\Delta_3 > 0$, i.e., $c^2 < 8b^3/(27a)$, the cubic has three distinct real roots
(two minima and one maximum). For $\Delta_3 = 0$, two roots coalesce (the barrier
disappears for one well). For $\Delta_3 < 0$, there is only one real root (one minimum).

The three real roots (when $\Delta_3 > 0$) are given by Cardano--Vieta's trigonometric method:
\begin{equation}
  x_k = 2\sqrt{\frac{b}{6a}}\cos\left[\frac{1}{3}\arccos\left(-\frac{c}{4a}\cdot\frac{1}{2}\left(\frac{6a}{b}\right)^{3/2}\right) - \frac{2\pi k}{3}\right],\quad k = 0,1,2.
\end{equation}
Simplifying the argument of the arccosine:
\begin{equation}
  x_k = 2\sqrt{\frac{b}{6a}}\cos\left[\frac{1}{3}\arccos\left(-\frac{3c}{4b}\sqrt{\frac{6a}{b}}\right) - \frac{2\pi k}{3}\right].
\end{equation}

\subsubsection{Change in barrier heights}

With the asymmetric potential, there are now two distinct barrier heights. Let $x_-$, $x_{\max}$, $x_+$ be the three critical points ordered as $x_- < x_{\max} < x_+$. The left barrier is
\begin{equation}
  V_{b,L} = V(x_{\max}) - V(x_-)\,,
\end{equation}
and the right barrier is
\begin{equation}
  V_{b,R} = V(x_{\max}) - V(x_+)\,.
\end{equation}
The difference is
\begin{equation}
  V_{b,L} - V_{b,R} = V(x_+) - V(x_-) \approx 2c\,\xe\,.
\end{equation}
The barrier is higher on the side of the shallower well.

\subsubsection{General polynomial potential}

For maximum generality, the code supports an arbitrary polynomial potential:
\begin{equation}\label{eq:V_poly}
  V(x) = \sum_{k=0}^{K} v_k\,x^k\,.
\end{equation}
The Hamiltonian matrix elements require the integrals $\mel{n}{x^k}{m}$ for each power $k$.
These are computed recursively using the ladder-operator formalism (see \cref{sec:HO_matrix_elements}).

% ============================================================
% ============================================================

\section{Variational Method Framework}\label{sec:variational}
% ============================================================
% ============================================================

% ------------------------------------------------------------
\subsection{Rayleigh--Ritz variational principle}
% ------------------------------------------------------------

The time-independent Schr\"odinger equation is
\begin{equation}\label{eq:TISE}
  \hat{H}\ket{\Psi} = E\ket{\Psi}\,.
\end{equation}
The variational principle states that for any normalized trial state $\ket{\tilde{\Psi}}$,
\begin{equation}\label{eq:variational_bound}
  \mel{\tilde{\Psi}}{\hat{H}}{\tilde{\Psi}} \geq E_0\,,
\end{equation}
with equality if and only if $\ket{\tilde{\Psi}}$ is the ground state.
More generally, the Hylleraas--Undheim--MacDonald (HUM) theorem guarantees that
for a variational calculation in an $N$-dimensional subspace, the $k$-th eigenvalue
$E_k^{(N)}$ is an upper bound to the $k$-th exact eigenvalue $E_k^{\text{exact}}$:
\begin{equation}
  E_k^{(N)} \geq E_k^{\text{exact}},\quad k = 0, 1, \ldots, N-1\,.
\end{equation}
Furthermore, enlarging the basis can only improve (lower) each variational eigenvalue:
$E_k^{(N+1)} \leq E_k^{(N)}$.

% ------------------------------------------------------------
\subsection{Expansion in a finite basis}
% ------------------------------------------------------------

We expand the trial wavefunction in a set of $N$ basis functions $\{\phi_i\}_{i=1}^{N}$:
\begin{equation}\label{eq:expansion}
  \ket{\tilde{\Psi}} = \sum_{i=1}^{N} c_i\,\ket{\phi_i}\,.
\end{equation}
The energy functional is
\begin{equation}
  E[\{c_i\}] = \frac{\mel{\tilde{\Psi}}{\hat{H}}{\tilde{\Psi}}}{\braket{\tilde{\Psi}}}
  = \frac{\sum_{i,j} c_i^* c_j\,H_{ij}}{\sum_{i,j} c_i^* c_j\,S_{ij}}\,,
\end{equation}
where $H_{ij} = \mel{\phi_i}{\hat{H}}{\phi_j}$ and $S_{ij} = \braket{\phi_i}{\phi_j}$
are the Hamiltonian and overlap matrix elements, respectively.

% ------------------------------------------------------------
\subsection{Derivation of the generalized eigenvalue problem}
% ------------------------------------------------------------

Minimizing $E$ with respect to the coefficients $c_k^*$ (treating $c_k$ and $c_k^*$ as independent):
\begin{align}
  \frac{\partial}{\partial c_k^*}\left[\sum_{i,j}c_i^*c_j\,H_{ij} - E\sum_{i,j}c_i^*c_j\,S_{ij}\right] &= 0\notag\\
  \sum_j H_{kj}\,c_j - E\sum_j S_{kj}\,c_j &= 0\,,\quad k = 1,\ldots, N\,.
\end{align}
In matrix form, this is the \emph{generalized eigenvalue problem}:
\begin{equation}\label{eq:gen_eig}
  \boxed{\mathbf{H}\,\mathbf{c} = E\,\mathbf{S}\,\mathbf{c}\,,}
\end{equation}
where $\mathbf{H}$ is the $N\times N$ Hamiltonian matrix, $\mathbf{S}$ is the $N\times N$
overlap matrix, and $\mathbf{c}$ is the coefficient vector.

\subsubsection{Standard eigenvalue problem (orthonormal basis)}

When the basis is orthonormal, $S_{ij} = \delta_{ij}$, and the generalized problem
reduces to the \emph{standard} eigenvalue problem:
\begin{equation}\label{eq:std_eig}
  \boxed{\mathbf{H}\,\mathbf{c} = E\,\mathbf{c}\,.}
\end{equation}
This is the case for the harmonic oscillator basis (Basis~I) and the DVR basis (Basis~IV).
The shifted HO basis (Basis~II) and Gaussian basis (Basis~III) have non-trivial overlap
and require either the generalized eigenvalue solver or a preliminary orthogonalization step.

% ------------------------------------------------------------
\subsection{L\"owdin canonical orthogonalization}\label{sec:Lowdin}
% ------------------------------------------------------------

When $\mathbf{S} \neq \mathbf{I}$, one can transform the generalized problem into
a standard one by the \emph{L\"owdin symmetric orthogonalization} (also called
canonical orthogonalization).

\subsubsection{Step 1: Diagonalize the overlap matrix}

Since $\mathbf{S}$ is real, symmetric, and positive semi-definite, it can be diagonalized:
\begin{equation}
  \mathbf{S} = \mathbf{U}\,\bm{\sigma}\,\mathbf{U}^T\,,
\end{equation}
where $\mathbf{U}$ is an orthogonal matrix of eigenvectors and
$\bm{\sigma} = \text{diag}(\sigma_1, \sigma_2, \ldots, \sigma_N)$
contains the eigenvalues $\sigma_i \geq 0$.

\subsubsection{Step 2: Remove near-linear dependencies}

Eigenvalues $\sigma_i < \eps_{\text{thresh}}$ (typically $\eps_{\text{thresh}} = 10^{-8}$)
indicate near-linear dependence in the basis. The corresponding eigenvectors are
discarded. Let $N' \leq N$ be the number of retained eigenvalues with
$\sigma_i \geq \eps_{\text{thresh}}$, and let $\mathbf{U}'$ and $\bm{\sigma}'$
denote the truncated matrices.

\subsubsection{Step 3: Construct the transformation matrix}

Define the matrix
\begin{equation}
  \mathbf{X} = \mathbf{U}'\,(\bm{\sigma}')^{-1/2}\,,
\end{equation}
where $(\bm{\sigma}')^{-1/2} = \text{diag}(\sigma_1^{-1/2}, \ldots, \sigma_{N'}^{-1/2})$.
This is an $N \times N'$ matrix.

\subsubsection{Step 4: Verify orthonormality}

The transformed overlap matrix is
\begin{equation}
  \mathbf{X}^T\,\mathbf{S}\,\mathbf{X}
  = (\bm{\sigma}')^{-1/2}\,(\mathbf{U}')^T\,\mathbf{U}'\,\bm{\sigma}'\,(\mathbf{U}')^T\,\mathbf{U}'\,(\bm{\sigma}')^{-1/2}
  = (\bm{\sigma}')^{-1/2}\,\bm{\sigma}'\,(\bm{\sigma}')^{-1/2}
  = \mathbf{I}_{N'}\,.
\end{equation}
The transformed basis is orthonormal.

\subsubsection{Step 5: Solve the standard eigenvalue problem}

Define the transformed Hamiltonian:
\begin{equation}
  \tilde{\mathbf{H}} = \mathbf{X}^T\,\mathbf{H}\,\mathbf{X}\,.
\end{equation}
This is an $N' \times N'$ real symmetric matrix. Diagonalize it:
\begin{equation}
  \tilde{\mathbf{H}}\,\tilde{\mathbf{c}}_k = E_k\,\tilde{\mathbf{c}}_k\,.
\end{equation}
The eigenvectors in the original (non-orthogonal) basis are recovered as
\begin{equation}\label{eq:Lowdin_backtransform}
  \mathbf{c}_k = \mathbf{X}\,\tilde{\mathbf{c}}_k\,.
\end{equation}

\begin{result}[L\"owdin procedure summary]
\begin{equation}
  \mathbf{S} \xrightarrow{\text{diag.}} \mathbf{U}\,\bm{\sigma}\,\mathbf{U}^T
  \xrightarrow{\text{truncate}} \mathbf{X} = \mathbf{U}'(\bm{\sigma}')^{-1/2}
  \xrightarrow{\text{transform}} \tilde{\mathbf{H}} = \mathbf{X}^T\mathbf{H}\mathbf{X}
  \xrightarrow{\text{diag.}} E_k,\;\mathbf{c}_k = \mathbf{X}\tilde{\mathbf{c}}_k\,.
\end{equation}
\end{result}

% ------------------------------------------------------------
\subsection{Convergence and the variational upper bound}
% ------------------------------------------------------------

By the HUM theorem, the variational eigenvalues $\{E_k^{(N)}\}$ provide upper bounds to
the exact eigenvalues and converge monotonically from above as $N \to \infty$.
A practical convergence criterion for the $k$-th level is
\begin{equation}
  |E_k^{(N)} - E_k^{(N-\delta N)}| < \delta_{\text{conv}}\,,
\end{equation}
where $\delta N$ is the basis increment and $\delta_{\text{conv}}$ is a prescribed tolerance
(e.g., $10^{-10}\,E_h$ for spectroscopic accuracy).

For the quartic double-well, the lowest levels typically converge with $N \sim 20$--$50$
harmonic oscillator functions. Higher excited states require larger bases due to the
polynomial growth of the potential.

% ============================================================
% ============================================================

\section{Basis Set~I: Harmonic Oscillator Basis (Centered at Origin)}\label{sec:HO}
% ============================================================
% ============================================================

This is the current implementation in QuTu. The basis functions are the eigenstates of the
one-dimensional harmonic oscillator centered at the origin:
\begin{equation}\label{eq:HO_basis}
  \varphi_n(x) \equiv \varphi_n(x;\alpha)
  = \mathcal{N}_n\,H_n\!\bigl(\sqrt{\alpha}\,x\bigr)\,e^{-\alpha x^2/2}\,,
  \qquad n = 0, 1, 2, \ldots,
\end{equation}
with normalization constant
\begin{equation}\label{eq:Nn}
  \mathcal{N}_n = \left(\frac{\alpha}{\pi}\right)^{1/4}\frac{1}{\sqrt{2^n\,n!}}\,.
\end{equation}
Here $H_n(\xi)$ are the (``physicist's'') Hermite polynomials satisfying the recurrence
\begin{equation}
  H_{n+1}(\xi) = 2\xi\,H_n(\xi) - 2n\,H_{n-1}(\xi)\,,\quad
  H_0(\xi) = 1,\;H_1(\xi) = 2\xi\,.
\end{equation}
The basis is orthonormal: $\braket{\varphi_m}{\varphi_n} = \delta_{mn}$.

% ------------------------------------------------------------
\subsection{Ladder operators}\label{sec:ladder}
% ------------------------------------------------------------

We define the lowering and raising operators
\begin{equation}\label{eq:ladder_def}
  \hat{a} = \sqrt{\frac{\alpha}{2}}\,x + \frac{1}{\sqrt{2\alpha}}\frac{d}{dx}\,,
  \qquad
  \adag = \sqrt{\frac{\alpha}{2}}\,x - \frac{1}{\sqrt{2\alpha}}\frac{d}{dx}\,,
\end{equation}
satisfying $[\hat{a}, \adag] = 1$ and
\begin{equation}
  \hat{a}\ket{n} = \sqrt{n}\,\ket{n-1}\,,\qquad
  \adag\ket{n} = \sqrt{n+1}\,\ket{n+1}\,.
\end{equation}

\subsubsection{Position and derivative in terms of ladder operators}

Solving \cref{eq:ladder_def} for $x$ and $d/dx$:
\begin{equation}\label{eq:x_and_ddx}
  \boxed{x = \frac{1}{\sqrt{2\alpha}}\bigl(\hat{a} + \adag\bigr)\,,\qquad
  \frac{d}{dx} = \sqrt{\frac{\alpha}{2}}\bigl(\hat{a} - \adag\bigr)\,.}
\end{equation}
One verifies:
\begin{equation}
  \hat{a} + \adag = 2\sqrt{\frac{\alpha}{2}}\,x = \sqrt{2\alpha}\,x \implies x = \frac{\hat{a}+\adag}{\sqrt{2\alpha}}\,.
\end{equation}

% ------------------------------------------------------------
\subsection{Matrix elements of $x^k$}\label{sec:HO_matrix_elements}
% ------------------------------------------------------------

\subsubsection{$\mel{m}{x}{n}$}

Using \cref{eq:x_and_ddx}:
\begin{align}
  \mel{m}{x}{n}
  &= \frac{1}{\sqrt{2\alpha}}\mel{m}{\hat{a}+\adag}{n}\notag\\
  &= \frac{1}{\sqrt{2\alpha}}\bigl(\sqrt{n}\,\delta_{m,n-1} + \sqrt{n+1}\,\delta_{m,n+1}\bigr)\,.
\end{align}
Defining $m_* = \min(m,n)$:
\begin{equation}\label{eq:x1_HO}
  \boxed{\mel{m}{x}{n} = \frac{\sqrt{m_*+1}}{\sqrt{2\alpha}}\,\delta_{|m-n|,1}\,.}
\end{equation}
\textbf{Selection rule:} nonzero only for $|m-n| = 1$ (tridiagonal matrix).

\subsubsection{$\mel{m}{x^2}{n}$}

We compute $x^2 = (\hat{a}+\adag)^2/(2\alpha)$ and expand:
\begin{equation}
  (\hat{a}+\adag)^2 = \hat{a}^2 + \hat{a}\adag + \adag\hat{a} + (\adag)^2
  = \hat{a}^2 + (\adag)^2 + 2\hat{N} + 1\,,
\end{equation}
where $\hat{N} = \adag\hat{a}$ is the number operator and we used $\hat{a}\adag = \hat{N} + 1$.

Acting on $\ket{n}$:
\begin{align}
  (\hat{a}+\adag)^2\ket{n}
  &= \sqrt{n(n-1)}\,\ket{n-2} + (2n+1)\ket{n} + \sqrt{(n+1)(n+2)}\,\ket{n+2}\,.
\end{align}
Therefore:
\begin{equation}\label{eq:x2_HO}
  \boxed{\mel{m}{x^2}{n} = \frac{1}{2\alpha}
  \begin{cases}
    \sqrt{(m_*+1)(m_*+2)} & |m-n| = 2\,,\\
    2n+1 & m = n\,,\\
    0 & \text{otherwise}\,.
  \end{cases}}
\end{equation}
Equivalently, the diagonal element is $(n+\half)/\alpha$.

\textbf{Selection rule:} nonzero for $|m-n| = 0$ or $2$ (pentadiagonal matrix).

% ------------------------------------------------------------
\subsubsection{$\mel{m}{x^3}{n}$}
% ------------------------------------------------------------

We compute $(\hat{a}+\adag)^3\ket{n}$ by applying $(\hat{a}+\adag)$ to
$(\hat{a}+\adag)^2\ket{n} = \sqrt{n(n-1)}\ket{n-2}+(2n+1)\ket{n}+\sqrt{(n+1)(n+2)}\ket{n+2}$
and collecting terms:
\begin{align}
  (\hat{a}+\adag)^3\ket{n}
  &= \sqrt{n(n-1)(n-2)}\,\ket{n-3} + 3n^{3/2}\,\ket{n-1}\notag\\
  &\quad + 3(n+1)^{3/2}\,\ket{n+1} + \sqrt{(n+1)(n+2)(n+3)}\,\ket{n+3}\,.
\end{align}
Since $\hat{x}^3 = (\hat{a}+\adag)^3/(2\alpha)^{3/2}$:
\begin{equation}\label{eq:x3_HO}
  \boxed{\mel{m}{x^3}{n} = \frac{1}{(2\alpha)^{3/2}}
  \begin{cases}
    3(m_*+1)^{3/2} & |m-n| = 1\,,\\[4pt]
    \sqrt{(m_*+1)(m_*+2)(m_*+3)} & |m-n| = 3\,,\\[4pt]
    0 & \text{otherwise}\,.
  \end{cases}}
\end{equation}

\textbf{Selection rule:} nonzero for $|m-n| = 1$ or $3$ (the odd-power operator
connects states of opposite parity). This is the term that appears in the
asymmetric potential $V(x) = ax^4 - bx^2 + \varepsilon x$ (see \cref{sec:V_asym})
and breaks the parity-block structure of the Hamiltonian matrix.

\subsubsection{$\mel{m}{x^4}{n}$: complete derivation}

We need $x^4 = (\hat{a}+\adag)^4/(2\alpha)^2$. We compute $(\hat{a}+\adag)^4\ket{n}$
by applying $(\hat{a}+\adag)$ to $(\hat{a}+\adag)^2\ket{n}$ twice, or equivalently by
applying $(\hat{a}+\adag)$ to $(\hat{a}+\adag)^3\ket{n}$.

\textbf{Step~1:} Compute $(\hat{a}+\adag)^3\ket{n}$.

Apply $(\hat{a}+\adag)$ to $(\hat{a}+\adag)^2\ket{n}$:
\begin{align}
  &(\hat{a}+\adag)\bigl[\sqrt{n(n-1)}\ket{n-2} + (2n+1)\ket{n} + \sqrt{(n+1)(n+2)}\ket{n+2}\bigr]\notag\\
  &= \sqrt{n(n-1)}\bigl[\sqrt{n-2}\ket{n-3}+\sqrt{n-1}\ket{n-1}\bigr]\notag\\
  &\quad + (2n+1)\bigl[\sqrt{n}\ket{n-1}+\sqrt{n+1}\ket{n+1}\bigr]\notag\\
  &\quad + \sqrt{(n+1)(n+2)}\bigl[\sqrt{n+2}\ket{n+1}+\sqrt{n+3}\ket{n+3}\bigr]\,.
\end{align}
Collecting coefficients for each ket:
\begin{align}
  \ket{n-3}&:\quad \sqrt{n(n-1)(n-2)}\,,\notag\\
  \ket{n-1}&:\quad (n-1)\sqrt{n} + (2n+1)\sqrt{n} = 3n\sqrt{n} = 3n^{3/2}\,,\notag\\
  \ket{n+1}&:\quad (2n+1)\sqrt{n+1} + (n+2)\sqrt{n+1} = 3(n+1)^{3/2}\,,\notag\\
  \ket{n+3}&:\quad \sqrt{(n+1)(n+2)(n+3)}\,.\notag
\end{align}
So:
\begin{align}\label{eq:aad3}
  (\hat{a}+\adag)^3\ket{n}
  &= \sqrt{n(n-1)(n-2)}\ket{n-3} + 3n\sqrt{n}\ket{n-1}\notag\\
  &\quad + 3(n+1)\sqrt{n+1}\ket{n+1} + \sqrt{(n+1)(n+2)(n+3)}\ket{n+3}\,.
\end{align}

\textbf{Step~2:} Compute $(\hat{a}+\adag)^4\ket{n} = (\hat{a}+\adag)\cdot[(\hat{a}+\adag)^3\ket{n}]$.

Applying $(\hat{a}+\adag)$ to each term in \cref{eq:aad3}:
\begin{align}
  &(\hat{a}+\adag)\sqrt{n(n-1)(n-2)}\ket{n-3}\notag\\
  &\qquad= \sqrt{n(n-1)(n-2)}\bigl[\sqrt{n-3}\ket{n-4}+\sqrt{n-2}\ket{n-2}\bigr]\,,\\[4pt]
  &(\hat{a}+\adag)\cdot 3n\sqrt{n}\ket{n-1}\notag\\
  &\qquad= 3n\sqrt{n}\bigl[\sqrt{n-1}\ket{n-2}+\sqrt{n}\ket{n}\bigr]\,,\\[4pt]
  &(\hat{a}+\adag)\cdot 3(n+1)\sqrt{n+1}\ket{n+1}\notag\\
  &\qquad= 3(n+1)\sqrt{n+1}\bigl[\sqrt{n+1}\ket{n}+\sqrt{n+2}\ket{n+2}\bigr]\,,\\[4pt]
  &(\hat{a}+\adag)\sqrt{(n+1)(n+2)(n+3)}\ket{n+3}\notag\\
  &\qquad= \sqrt{(n+1)(n+2)(n+3)}\bigl[\sqrt{n+3}\ket{n+2}+\sqrt{n+4}\ket{n+4}\bigr]\,.
\end{align}

Collecting coefficients:
\begin{align}
  \ket{n-4}&:\quad \sqrt{n(n-1)(n-2)(n-3)}\,,\notag\\[3pt]
  \ket{n-2}&:\quad (n-2)\sqrt{n(n-1)} + 3n\sqrt{n(n-1)}\notag\\
            &\phantom{:}\quad = (n-2+3n)\sqrt{n(n-1)} = (4n-2)\sqrt{n(n-1)}\notag\\
            &\phantom{:}\quad = 2(2n-1)\sqrt{n(n-1)}\,,\notag\\[3pt]
  \ket{n}&:\quad 3n^2 + 3(n+1)^2 = 3(2n^2+2n+1)\,,\notag\\[3pt]
  \ket{n+2}&:\quad 3(n+1)\sqrt{(n+1)(n+2)} + (n+3)\sqrt{(n+1)(n+2)}\notag\\
            &\phantom{:}\quad = (3n+3+n+3)\sqrt{(n+1)(n+2)} = (4n+6)\sqrt{(n+1)(n+2)}\notag\\
            &\phantom{:}\quad = 2(2n+3)\sqrt{(n+1)(n+2)}\,,\notag\\[3pt]
  \ket{n+4}&:\quad \sqrt{(n+1)(n+2)(n+3)(n+4)}\,.\notag
\end{align}

Therefore, since $x^4 = (\hat{a}+\adag)^4/(4\alpha^2)$:

\begin{equation}\label{eq:x4_HO}
  \boxed{\mel{m}{x^4}{n} = \frac{1}{4\alpha^2}
  \begin{cases}
    \sqrt{(m_*\!+\!1)(m_*\!+\!2)(m_*\!+\!3)(m_*\!+\!4)} & |m-n| = 4\,,\\[4pt]
    2(2m_*+3)\sqrt{(m_*\!+\!1)(m_*\!+\!2)} & |m-n| = 2\,,\\[4pt]
    3(2n^2+2n+1) & m = n\,,\\[4pt]
    0 & \text{otherwise}\,.
  \end{cases}}
\end{equation}

\textbf{Selection rule:} nonzero for $|m-n| = 0, 2, 4$.

% ------------------------------------------------------------
\subsection{Kinetic energy matrix elements}\label{sec:HO_kinetic}
% ------------------------------------------------------------

The kinetic energy operator in atomic units ($\hb = 1$) is
\begin{equation}
  \hat{T} = -\frac{1}{2\mu}\frac{d^2}{dx^2}\,.
\end{equation}
Using $d/dx = \sqrt{\alpha/2}(\hat{a} - \adag)$:
\begin{align}
  \frac{d^2}{dx^2} &= \frac{\alpha}{2}(\hat{a}-\adag)^2
  = \frac{\alpha}{2}\bigl(\hat{a}^2 + (\adag)^2 - 2\hat{N} - 1\bigr)\,.
\end{align}
Therefore:
\begin{align}
  \hat{T} &= -\frac{1}{2\mu}\cdot\frac{\alpha}{2}\bigl(\hat{a}^2 + (\adag)^2 - 2\hat{N} - 1\bigr)\notag\\
  &= \frac{\alpha}{4\mu}\bigl(2\hat{N}+1 - \hat{a}^2 - (\adag)^2\bigr)\,.
\end{align}

Matrix elements:
\begin{align}
  T_{mn} &= \mel{m}{\hat{T}}{n}
  = \frac{\alpha}{4\mu}\mel{m}{2\hat{N}+1-\hat{a}^2-(\adag)^2}{n}\notag\\
  &= \frac{\alpha}{4\mu}\Bigl[(2n+1)\delta_{mn} - \sqrt{n(n-1)}\,\delta_{m,n-2} - \sqrt{(n+1)(n+2)}\,\delta_{m,n+2}\Bigr]\,.
\end{align}
Using $m_* = \min(m,n)$:
\begin{equation}\label{eq:T_HO}
  \boxed{T_{mn} = \frac{\alpha}{4\mu}
  \begin{cases}
    2n + 1 & m = n\,,\\[4pt]
    -\sqrt{(m_*+1)(m_*+2)} & |m-n| = 2\,,\\[4pt]
    0 & \text{otherwise}\,.
  \end{cases}}
\end{equation}
Equivalently, $T_{nn} = \frac{\alpha}{2\mu}(n+\half)$.

\textbf{Verification:} For the harmonic oscillator with $V = \half\mu\omega^2 x^2$
and $\alpha = \mu\omega$, the total energy should be $E_n = \omega(n+\half)$.
The kinetic energy contribution to the diagonal is $T_{nn} = \frac{\mu\omega}{2\mu}(n+\half)
= \frac{\omega}{2}(n+\half)$, which equals exactly half the total energy, consistent
with the virial theorem for a harmonic potential.

% ------------------------------------------------------------
\subsection{Symmetry factorization: even and odd subspaces}\label{sec:parity}
% ------------------------------------------------------------

For the symmetric double-well $V(x) = ax^4 - bx^2$, the potential satisfies $V(-x) = V(x)$.
The parity operator $\hat{\Pi}$ acts as $\hat{\Pi}\varphi_n(x) = \varphi_n(-x) = (-1)^n\varphi_n(x)$
(since $H_n(-\xi) = (-1)^n H_n(\xi)$).

\textbf{Claim:} $H_{mn} = 0$ whenever $m + n$ is odd.

\textbf{Proof:} The Hamiltonian matrix element consists of:
\begin{enumerate}
  \item $T_{mn}$: nonzero only for $|m-n| = 0$ or $2$, both of which require $m+n$ even.
  \item $\mel{m}{x^2}{n}$: nonzero only for $|m-n| = 0$ or $2$, requiring $m+n$ even.
  \item $\mel{m}{x^4}{n}$: nonzero only for $|m-n| = 0$, $2$, or $4$, all requiring $m+n$ even.
\end{enumerate}
Since all terms vanish when $m+n$ is odd, $H_{mn} = 0$ for $m+n$ odd. $\square$

Therefore, the Hamiltonian matrix decouples into two blocks:
\begin{equation}
  \mathbf{H} = \begin{pmatrix} \mathbf{H}_{\text{even}} & \mathbf{0} \\ \mathbf{0} & \mathbf{H}_{\text{odd}} \end{pmatrix},
\end{equation}
where $\mathbf{H}_{\text{even}}$ is built from $\{|0\rangle, |2\rangle, |4\rangle, \ldots\}$
and $\mathbf{H}_{\text{odd}}$ from $\{|1\rangle, |3\rangle, |5\rangle, \ldots\}$.

\textbf{Dimensions.} For a total of $N$ basis functions $\{|0\rangle, |1\rangle, \ldots, |N-1\rangle\}$:
\begin{equation}
  N_{\text{even}} = \left\lceil\frac{N}{2}\right\rceil\,,\qquad
  N_{\text{odd}} = \left\lfloor\frac{N}{2}\right\rfloor\,.
\end{equation}
The even-parity block yields the even-parity eigenstates $\Phi_0, \Phi_2, \Phi_4, \ldots$
(including the ground state), and the odd-parity block yields $\Phi_1, \Phi_3, \Phi_5, \ldots$.
The combined sorted spectrum is $E_0 < E_1 < E_2 < E_3 < \cdots$.

The computational cost is $2 \times (N/2)^3 = N^3/4$, a factor of 4 savings over full
$N\times N$ diagonalization.

For the \emph{asymmetric} potential $V(x) = ax^4 - bx^2 + cx$, the linear term couples
even and odd states ($\mel{m}{x}{n}$ is nonzero for $|m-n|=1$, i.e., $m+n$ odd), and the
full $N\times N$ matrix must be diagonalized.

\subsubsection*{Asymmetric case: full $N\times N$ matrix structure}

The asymmetric Hamiltonian matrix element combines the symmetric double-well terms with
the linear perturbation $\varepsilon\hat{x}$:
\begin{equation}\label{eq:H_asym}
  H_{mn}^{\text{asym}} = H_{mn}^{\text{sym}} + \varepsilon\mel{m}{\hat{x}}{n}\,.
\end{equation}
The $\mel{m}{\hat{x}}{n}$ term (see \cref{eq:x1_HO}) is nonzero for $|m-n|=1$, which
connects even and odd basis functions. The result is a real symmetric \emph{pentadiagonal}
matrix of bandwidth~4 (nonzero off-diagonals at $|m-n|=1,2,4$):
\begin{equation}
  \mathbf{H}^{\text{asym}} =
  \begin{pmatrix}
    D_0    & L_{01} & S_{02} & 0      & Q_{04} & \cdots \\
    L_{01} & D_1    & L_{12} & S_{13} & 0      & \cdots \\
    S_{02} & L_{12} & D_2    & L_{23} & S_{24} & \cdots \\
    0      & S_{13} & L_{23} & D_3    & L_{34} & \cdots \\
    Q_{04} & 0      & S_{24} & L_{34} & D_4    & \cdots \\
    \vdots & \vdots & \vdots & \vdots & \vdots & \ddots
  \end{pmatrix},
\end{equation}
where $D_n$ are diagonal elements; $L_{n,n\pm1} = \varepsilon\sqrt{\minn+1}/\sqrt{2\alpha}$
are the linear couplings ($|m-n|=1$) from $\mel{m}{\hat{x}}{n}$; $S_{mn}$ are the
quadratic couplings ($|m-n|=2$) from $\mel{m}{x^2}{n}$ (see \cref{eq:x2_HO}); and
$Q_{mn}$ are the quartic couplings ($|m-n|=4$) from $\mel{m}{x^4}{n}$ (see
\cref{eq:x4_HO}). Note that $|m-n|=3$ remains zero because there is no cubic term
in the symmetric part, and $\mel{m}{\hat{x}}{n}$ does not contribute at $|m-n|=3$.
The full matrix is diagonalized by a real symmetric eigensolver (e.g., LAPACK
\texttt{DSYEV}).

% ------------------------------------------------------------
\subsection{Optimal $\alpha$ parameter}\label{sec:alpha_opt}
% ------------------------------------------------------------

The width parameter $\alpha$ determines the spatial extent of the basis functions.
We seek $\alpha_{\text{opt}}$ that minimizes the ground-state energy.

\subsubsection{Variational optimization}

Consider the ground-state expectation value of the Hamiltonian $\hat{H} = \hat{T} + ax^4 - bx^2$
in the state $\ket{0}$:
\begin{align}
  E_0(\alpha) &= T_{00} + a\mel{0}{x^4}{0} - b\mel{0}{x^2}{0}\notag\\
  &= \frac{\alpha}{4\mu} + a\cdot\frac{3}{4\alpha^2} - b\cdot\frac{1}{2\alpha}\,.
\end{align}
Setting $dE_0/d\alpha = 0$:
\begin{align}
  \frac{dE_0}{d\alpha} &= \frac{1}{4\mu} - \frac{3a}{2\alpha^3} + \frac{b}{2\alpha^2} = 0\,.
\end{align}
This is a cubic equation in $\alpha$ that does not have a simple closed-form solution
in general. However, in two limiting regimes, simple expressions emerge.

\subsubsection{Deep-well limit: harmonic approximation}

When the barrier is much higher than the zero-point energy, the particle is well-localized
in one of the wells. Near the minimum at $x_{\min}$, the potential is approximately
harmonic: $V(x) \approx V(x_{\min}) + \half V''(x_{\min})(x-x_{\min})^2$ with
$V''(x_{\min}) = 4b$ (see \cref{sec:V_sym}). The harmonic oscillator with frequency
$\omega_{\text{local}} = \sqrt{4b/\mu}$ has the optimal width parameter
\begin{equation}\label{eq:alpha_local_opt}
  \boxed{\alpha_{\text{opt}} = \mu\,\omega_{\text{local}} = 2\sqrt{\mu\,b}\,.}
\end{equation}
This follows from the standard HO result $\alpha = m\omega/\hb$, with $m = \mu$,
$\omega = \omega_{\text{local}}$, and $\hb = 1$ in atomic units.

\subsubsection{Pure quartic limit}

When $b = 0$ (no quadratic term, pure quartic oscillator $V = ax^4$), the optimization
condition becomes
\begin{equation}
  \frac{1}{4\mu} = \frac{3a}{2\alpha^3}
  \implies \alpha^3 = 6a\mu
  \implies \alpha_{\text{opt}} = (6a\mu)^{1/3}\,.
\end{equation}
Substituting $a = \Vb/\xe^4$:
\begin{equation}\label{eq:alpha_quartic}
  \boxed{\alpha_{\text{opt}} = \left(\frac{6\mu\Vb}{\xe^4}\right)^{1/3}\,.}
\end{equation}
This is the expression used when the quartic confinement dominates.

\begin{remark}
In practice, a numerical optimization of $\alpha$ with respect to the full
ground-state energy (not just the $\ket{0}$ expectation value) yields the best results.
This involves repeatedly diagonalizing the Hamiltonian for different $\alpha$ values.
The analytical estimates \eqref{eq:alpha_local_opt} and \eqref{eq:alpha_quartic}
serve as excellent starting points for such an optimization.
\end{remark}

% ============================================================
% ============================================================

\section{Basis Set~II: Shifted Harmonic Oscillator (Two-Center Basis)}\label{sec:shifted_HO}
% ============================================================
% ============================================================

The harmonic oscillator basis centered at the origin converges slowly for deep double wells
because none of the basis functions are concentrated near the potential minima.
A natural improvement is to use HO functions centered at the two minima.

% ------------------------------------------------------------
\subsection{Definition}
% ------------------------------------------------------------

Define left-centered and right-centered HO functions:
\begin{align}\label{eq:shifted_HO_def}
  \varphi_n^{(L)}(x) &= \varphi_n(x - \xL;\alpha) = \mathcal{N}_n\,H_n\!\bigl(\sqrt{\alpha}(x-\xL)\bigr)\,e^{-\alpha(x-\xL)^2/2}\,,\\
  \varphi_n^{(R)}(x) &= \varphi_n(x - \xR;\alpha) = \mathcal{N}_n\,H_n\!\bigl(\sqrt{\alpha}(x-\xR)\bigr)\,e^{-\alpha(x-\xR)^2/2}\,,
\end{align}
where $\xL = -\xe$ and $\xR = +\xe$ are the locations of the two potential minima,
and $\alpha$ is the common width parameter (typically $\alpha = \alpha_{\text{local}}$
from \cref{eq:alpha_local_opt}).

The full basis consists of $N_L$ left-centered and $N_R$ right-centered functions:
\begin{equation}
  \bigl\{\varphi_0^{(L)}, \varphi_1^{(L)}, \ldots, \varphi_{N_L-1}^{(L)},\;
         \varphi_0^{(R)}, \varphi_1^{(R)}, \ldots, \varphi_{N_R-1}^{(R)}\bigr\}\,.
\end{equation}
The total number of basis functions is $N = N_L + N_R$.

% ------------------------------------------------------------
\subsection{Overlap matrix}\label{sec:shifted_overlap}
% ------------------------------------------------------------

\subsubsection{Same-center overlaps (trivial)}

Since the HO eigenstates centered at the same point form an orthonormal set:
\begin{equation}
  S_{mn}^{LL} = \braket{\varphi_m^{(L)}}{\varphi_n^{(L)}} = \delta_{mn}\,,\qquad
  S_{mn}^{RR} = \braket{\varphi_m^{(R)}}{\varphi_n^{(R)}} = \delta_{mn}\,.
\end{equation}

\subsubsection{Cross-center overlaps: $S_{mn}^{LR}$}

The nontrivial overlaps are
\begin{equation}
  S_{mn}^{LR} = \braket{\varphi_m^{(L)}}{\varphi_n^{(R)}}
  = \int_{-\infty}^{\infty}\varphi_m(x-\xL;\alpha)\,\varphi_n(x-\xR;\alpha)\,dx\,.
\end{equation}
Define the displacement $d = \xR - \xL = 2\xe$.

\textbf{Derivation via the displacement operator.}
A displaced harmonic oscillator eigenstate can be expanded in terms of the undisplaced eigenstates.
The coherent-state displacement operator is
\begin{equation}
  \hat{D}(\beta) = e^{\beta\adag - \beta^*\hat{a}}\,,\quad
  \text{with}\;\beta = \sqrt{\frac{\alpha}{2}}\,x_0\;\text{(real for real displacement $x_0$)}\,.
\end{equation}
The displaced state is
\begin{equation}
  \ket{n; x_0} \equiv \varphi_n(x - x_0) = \hat{D}(\beta)\ket{n}\,.
\end{equation}
The overlap between two displaced states centered at $\xL$ and $\xR$ is
\begin{equation}
  S_{mn}^{LR} = \mel{m}{\hat{D}^\dagger(\beta_L)\hat{D}(\beta_R)}{n}
  = \mel{m}{\hat{D}(\beta_R - \beta_L)\,e^{i\phi}}{n}\,,
\end{equation}
where $\beta_L = \sqrt{\alpha/2}\,\xL$, $\beta_R = \sqrt{\alpha/2}\,\xR$, and
$\phi = \text{Im}(\beta_L^*\beta_R) = 0$ since both displacements are real.
Thus,
\begin{equation}
  S_{mn}^{LR} = \mel{m}{\hat{D}(\gamma)}{n}\,,\qquad \gamma = \beta_R - \beta_L = \sqrt{\frac{\alpha}{2}}\,d = \sqrt{2\alpha}\,\xe\,.
\end{equation}

The matrix elements of the displacement operator in the Fock basis are given by the
well-known formula (derivable from the Baker--Campbell--Hausdorff decomposition):
\begin{equation}\label{eq:displacement_matrix}
  \mel{m}{\hat{D}(\gamma)}{n} = e^{-\gamma^2/2}\sqrt{\frac{n!}{m!}}\,\gamma^{m-n}\,L_n^{(m-n)}(\gamma^2)\,,
  \quad m \geq n\,,
\end{equation}
where $L_n^{(k)}$ is the associated Laguerre polynomial. For $m < n$, use
$\mel{m}{\hat{D}(\gamma)}{n} = (-1)^{n-m}\mel{n}{\hat{D}(\gamma)}{m}|_{\gamma\to -\gamma}$.
For real $\gamma$, this simplifies to:
\begin{equation}
  \mel{m}{\hat{D}(\gamma)}{n} = (-1)^{m-n}\mel{n}{\hat{D}(\gamma)}{m}\quad\text{when }m < n\text{ and }\gamma \text{ real}.
\end{equation}

More explicitly, for $m \geq n$ and real $\gamma$:
\begin{equation}\label{eq:S_LR}
  \boxed{S_{mn}^{LR} = e^{-\gamma^2/2}\sqrt{\frac{n!}{m!}}\,\gamma^{m-n}\,L_n^{(m-n)}(\gamma^2)\,,
  \quad \gamma = \sqrt{2\alpha}\,\xe\,.}
\end{equation}
For $m < n$: $S_{mn}^{LR} = (-1)^{n-m}S_{nm}^{LR}$.

Note that $S_{mn}^{RL} = (S_{nm}^{LR})^*= S_{nm}^{LR}$ since the basis functions are real.

\textbf{Special cases:}
\begin{align}
  S_{00}^{LR} &= e^{-\gamma^2/2}\,,\\
  S_{10}^{LR} &= \gamma\,e^{-\gamma^2/2}\,,\\
  S_{01}^{LR} &= -\gamma\,e^{-\gamma^2/2}\,,\\
  S_{11}^{LR} &= (1-\gamma^2)\,e^{-\gamma^2/2}\,.
\end{align}

\textbf{Derivation of \cref{eq:displacement_matrix}.}
Write $\hat{D}(\gamma) = e^{-\gamma^2/2}e^{\gamma\adag}e^{-\gamma\hat{a}}$ (normal ordering via BCH).
Then:
\begin{align}
  \mel{m}{\hat{D}(\gamma)}{n}
  &= e^{-\gamma^2/2}\mel{m}{e^{\gamma\adag}e^{-\gamma\hat{a}}}{n}\notag\\
  &= e^{-\gamma^2/2}\sum_{k=0}^{n}\frac{(-\gamma)^k}{k!}\sqrt{\frac{n!}{(n-k)!}}
     \,\mel{m}{e^{\gamma\adag}}{n-k}\notag\\
  &= e^{-\gamma^2/2}\sum_{k=0}^{n}\frac{(-\gamma)^k}{k!}\sqrt{\frac{n!}{(n-k)!}}
     \sum_{j=0}^{\infty}\frac{\gamma^j}{j!}\sqrt{\frac{(n-k+j)!}{(n-k)!}}\,\delta_{m,n-k+j}\,.
\end{align}
Setting $j = m - n + k$ (which requires $j \geq 0$, i.e., $k \geq n - m$):
\begin{align}
  \mel{m}{\hat{D}(\gamma)}{n}
  &= e^{-\gamma^2/2}\sum_{k=\max(0,n-m)}^{n}
     \frac{(-1)^k\gamma^{m-n+2k}}{k!(m-n+k)!}\sqrt{\frac{n!\,m!}{[(n-k)!]^2}}\notag\\
  &= e^{-\gamma^2/2}\sqrt{\frac{n!}{m!}}\,\gamma^{m-n}
     \sum_{k=0}^{n}\frac{(-1)^k\gamma^{2k}\,n!}{k!(m-n+k)!(n-k)!}\quad (m\geq n)\,.
\end{align}
The sum is recognized as the associated Laguerre polynomial:
$L_n^{(m-n)}(\gamma^2) = \sum_{k=0}^{n}\binom{m}{n-k}\frac{(-\gamma^2)^k}{k!}$, confirming \cref{eq:displacement_matrix}.

% ------------------------------------------------------------
\subsection{Kinetic energy matrix elements}\label{sec:shifted_kinetic}
% ------------------------------------------------------------

\subsubsection{Same-center blocks: $T_{mn}^{LL}$, $T_{mn}^{RR}$}

Since the HO eigenstates centered at the same point are simply translated copies of
the origin-centered HO, and the kinetic energy operator $\hat{T} = -\frac{1}{2\mu}\frac{d^2}{dx^2}$
is translation-invariant, the matrix elements are identical to the origin-centered case:
\begin{equation}
  T_{mn}^{LL} = T_{mn}^{RR} = \frac{\alpha}{4\mu}
  \begin{cases}
    2n+1 & m = n\,,\\
    -\sqrt{(m_*+1)(m_*+2)} & |m-n|=2\,,\\
    0 & \text{otherwise}\,.
  \end{cases}
\end{equation}

\subsubsection{Cross-center blocks: $T_{mn}^{LR}$}

The cross-center kinetic energy matrix element is
\begin{equation}
  T_{mn}^{LR} = -\frac{1}{2\mu}\int_{-\infty}^{\infty}\varphi_m^{(L)}(x)\,\frac{d^2}{dx^2}\varphi_n^{(R)}(x)\,dx\,.
\end{equation}
Using the expansion of displaced HO states in terms of origin-centered states:
\begin{equation}
  \varphi_n^{(R)}(x) = \sum_k d_{nk}\,\varphi_k(x)\,,\qquad
  \varphi_m^{(L)}(x) = \sum_l d'_{ml}\,\varphi_l(x)\,,
\end{equation}
where $d_{nk} = \mel{k}{\hat{D}(\beta_R)}{n}$ and $d'_{ml} = \mel{l}{\hat{D}(\beta_L)}{m}$.

Then:
\begin{equation}
  T_{mn}^{LR} = \sum_{k,l} (d'_{ml})^* d_{nk}\,T_{lk}^{(0)}\,,
\end{equation}
where $T_{lk}^{(0)}$ are the origin-centered kinetic energy matrix elements.

In practice, it is more efficient to compute $T_{mn}^{LR}$ using the identity
\begin{equation}
  T_{mn}^{LR} = -\frac{1}{2\mu}\mel{\varphi_m^{(L)}}{\frac{d^2}{dx^2}}{\varphi_n^{(R)}}
  = \frac{1}{2\mu}\int_{-\infty}^{\infty}\frac{d\varphi_m^{(L)}}{dx}\cdot\frac{d\varphi_n^{(R)}}{dx}\,dx\,,
\end{equation}
where the second equality follows from integration by parts (boundary terms vanish).
The derivative of a shifted HO function is
\begin{equation}
  \frac{d}{dx}\varphi_n(x-x_0;\alpha) = \sqrt{\frac{\alpha}{2}}\Bigl[\sqrt{n}\,\varphi_{n-1}(x-x_0;\alpha) - \sqrt{n+1}\,\varphi_{n+1}(x-x_0;\alpha)\Bigr]\,.
\end{equation}
Therefore:
\begin{align}
  T_{mn}^{LR}
  &= \frac{\alpha}{4\mu}\sum_{j,k}(-1)^{p_j}(-1)^{p_k}s_j s_k\,S^{LR}_{m+\delta_j,\,n+\delta_k}\,,
\end{align}
where the sum runs over the terms produced by the derivative:
\begin{align}
  T_{mn}^{LR}
  &= \frac{\alpha}{4\mu}\Bigl[
    \sqrt{mn}\,S_{m-1,n-1}^{LR}
    - \sqrt{m(n+1)}\,S_{m-1,n+1}^{LR}\notag\\
  &\qquad\qquad
    - \sqrt{(m+1)n}\,S_{m+1,n-1}^{LR}
    + \sqrt{(m+1)(n+1)}\,S_{m+1,n+1}^{LR}
  \Bigr]\,.
\end{align}

\begin{equation}\label{eq:T_LR}
  \boxed{T_{mn}^{LR} = \frac{\alpha}{4\mu}\Bigl[
    \sqrt{mn}\,S_{m-1,n-1}^{LR}
    - \sqrt{m(n\!+\!1)}\,S_{m-1,n+1}^{LR}
    - \sqrt{(m\!+\!1)n}\,S_{m+1,n-1}^{LR}
    + \sqrt{(m\!+\!1)(n\!+\!1)}\,S_{m+1,n+1}^{LR}
  \Bigr]}
\end{equation}

% ------------------------------------------------------------
\subsection{Potential energy matrix elements}\label{sec:shifted_potential}
% ------------------------------------------------------------

For $V(x) = ax^4 - bx^2$, we need the matrix elements $\mel{\varphi_m^{(P)}}{x^k}{\varphi_n^{(Q)}}$
where $P, Q \in \{L, R\}$.

\subsubsection{Binomial expansion technique}

The key idea is to express $x^k$ in terms of displacements from one of the centers.
For the $(L,R)$ block, write $x = (x - \xR) + \xR$ and expand:
\begin{equation}
  x^k = \sum_{j=0}^{k}\binom{k}{j}(x-\xR)^j\,\xR^{k-j}\,.
\end{equation}
Then:
\begin{equation}
  \mel{\varphi_m^{(L)}}{x^k}{\varphi_n^{(R)}}
  = \sum_{j=0}^{k}\binom{k}{j}\xR^{k-j}\mel{\varphi_m^{(L)}}{(x-\xR)^j}{\varphi_n^{(R)}}\,.
\end{equation}
The matrix element $\mel{\varphi_m^{(L)}}{(x-\xR)^j}{\varphi_n^{(R)}}$ involves the
$j$-th power of the position operator measured from $\xR$, which is diagonal in the
$R$-centered ladder operators. Specifically, letting $\hat{a}_R$ and $\hat{a}_R^\dagger$
be the ladder operators for the $R$-centered HO:
\begin{equation}
  x - \xR = \frac{\hat{a}_R + \hat{a}_R^\dagger}{\sqrt{2\alpha}}\,.
\end{equation}
The matrix elements of $(x-\xR)^j$ in the $R$-centered basis are given by
\cref{eq:x1_HO,eq:x2_HO,eq:x4_HO} with the replacement $x \to x - \xR$.
Therefore:
\begin{equation}
  \mel{\varphi_m^{(L)}}{(x-\xR)^j}{\varphi_n^{(R)}}
  = \sum_p \mel{m}{\hat{D}^\dagger(\beta_L-\beta_R)}{p}\,\mel{p}{(x-\xR)^j}{n}_R\,,
\end{equation}
where $\mel{p}{(x-\xR)^j}{n}_R$ is the standard HO matrix element of $x^j$ in the $R$-centered basis.

Alternatively, for the $(L,L)$ block, $x = (x-\xL) + \xL$ and the matrix elements of
$(x-\xL)^j$ in the $L$-centered basis are simply the origin-centered HO matrix elements
$\mel{m}{x^j}{n}$ computed in \cref{sec:HO_matrix_elements}:
\begin{equation}
  \mel{\varphi_m^{(L)}}{x^k}{\varphi_n^{(L)}}
  = \sum_{j=0}^{k}\binom{k}{j}\xL^{k-j}\mel{m}{(x-\xL)^j}{n}_L
  = \sum_{j=0}^{k}\binom{k}{j}\xL^{k-j}\mel{m}{x^j}{n}_{\text{HO}}\,.
\end{equation}

\subsubsection{Explicit formulas for $k = 2$}

For the same-center $(L,L)$ block:
\begin{align}
  \mel{\varphi_m^{(L)}}{x^2}{\varphi_n^{(L)}}
  &= \mel{m}{(x-\xL)^2 + 2\xL(x-\xL) + \xL^2}{n}_L\notag\\
  &= \mel{m}{x^2}{n}_{\text{HO}} + 2\xL\mel{m}{x}{n}_{\text{HO}} + \xL^2\,\delta_{mn}\,.
\end{align}

For the cross-center $(L,R)$ block:
\begin{equation}
  \mel{\varphi_m^{(L)}}{x^2}{\varphi_n^{(R)}}
  = \mel{m}{(x-\xR)^2}{n}^{LR} + 2\xR\mel{m}{(x-\xR)}{n}^{LR} + \xR^2\,S_{mn}^{LR}\,,
\end{equation}
where $\mel{m}{(x-\xR)^j}{n}^{LR} = \sum_p S_{mp}^{LR}\mel{p}{x^j}{n}_{\text{HO}}$ (since
$(x-\xR)^j$ has the standard HO matrix elements in the $R$-centered basis, and the
overlap $S_{mp}^{LR}$ converts from the $L$-basis row index to the $R$-basis).

More explicitly:
\begin{align}
  \mel{m}{(x-\xR)}{n}^{LR} &= \sum_p S_{mp}^{LR}\frac{1}{\sqrt{2\alpha}}\bigl(\sqrt{n}\,\delta_{p,n-1}+\sqrt{n+1}\,\delta_{p,n+1}\bigr)\notag\\
  &= \frac{1}{\sqrt{2\alpha}}\bigl(\sqrt{n}\,S_{m,n-1}^{LR}+\sqrt{n+1}\,S_{m,n+1}^{LR}\bigr)\,.
\end{align}

\subsubsection{Explicit formulas for $k = 4$}

By the same binomial expansion, writing $x = (x-\xR) + \xR$:
\begin{align}
  x^4 &= (x-\xR)^4 + 4\xR(x-\xR)^3 + 6\xR^2(x-\xR)^2 + 4\xR^3(x-\xR) + \xR^4\,.
\end{align}
Each term's matrix element in the $(L,R)$ block is computed by inserting the overlap:
\begin{align}
  \mel{m}{(x-\xR)^j}{n}^{LR} = \sum_p S_{mp}^{LR}\mel{p}{x^j}{n}_{\text{HO}}\,,
\end{align}
using the known HO matrix elements of $x^j$ for $j = 0, 1, 2, 3, 4$.

\subsubsection{Extension to asymmetric potential}

For $V(x) = ax^4 - bx^2 + cx$, the additional linear term adds:
\begin{equation}
  c\mel{\varphi_m^{(P)}}{x}{\varphi_n^{(Q)}} = c\bigl[\mel{m}{x-x_Q}{n}^{PQ} + x_Q\,S_{mn}^{PQ}\bigr]\,.
\end{equation}

% ============================================================
% ============================================================

\section{Basis Set~III: Gaussian Basis Functions}\label{sec:gaussian}
% ============================================================
% ============================================================

An alternative non-orthogonal basis consists of unnormalized Gaussians:
\begin{equation}\label{eq:Gaussian_def}
  G_i(x) = \exp\bigl[-\alpha_i(x - x_i)^2\bigr]\,,\qquad i = 1, 2, \ldots, N\,,
\end{equation}
where $\alpha_i > 0$ is the width parameter and $x_i$ is the center of the $i$-th Gaussian.

% ------------------------------------------------------------
\subsection{Product of two Gaussians}\label{sec:Gauss_product}
% ------------------------------------------------------------

\begin{result}[Gaussian product theorem]
\begin{equation}\label{eq:Gauss_product}
  G_i(x)\,G_j(x) = e^{-\Qij}\,e^{-\Aij(x - \Xij)^2}\,,
\end{equation}
where
\begin{equation}\label{eq:AXQ}
  \Aij = \alpha_i + \alpha_j\,,\qquad
  \Xij = \frac{\alpha_i x_i + \alpha_j x_j}{\Aij}\,,\qquad
  \Qij = \frac{\alpha_i\alpha_j(x_i - x_j)^2}{\Aij}\,.
\end{equation}
\end{result}

\textbf{Complete derivation.} The product is
\begin{align}
  G_i(x)\,G_j(x) &= \exp\bigl[-\alpha_i(x-x_i)^2 - \alpha_j(x-x_j)^2\bigr]\,.
\end{align}
Expand the exponent:
\begin{align}
  f(x) &\equiv \alpha_i(x-x_i)^2 + \alpha_j(x-x_j)^2\notag\\
  &= \alpha_i(x^2 - 2x\,x_i + x_i^2) + \alpha_j(x^2 - 2x\,x_j + x_j^2)\notag\\
  &= (\alpha_i+\alpha_j)x^2 - 2(\alpha_i x_i + \alpha_j x_j)x + \alpha_i x_i^2 + \alpha_j x_j^2\notag\\
  &= \Aij\,x^2 - 2\Aij\Xij\,x + \alpha_i x_i^2 + \alpha_j x_j^2\,.
\end{align}
Complete the square in $x$:
\begin{align}
  f(x) &= \Aij\bigl(x^2 - 2\Xij\,x + \Xij^2\bigr) - \Aij\Xij^2 + \alpha_i x_i^2 + \alpha_j x_j^2\notag\\
  &= \Aij(x - \Xij)^2 + \bigl(\alpha_i x_i^2 + \alpha_j x_j^2 - \Aij\Xij^2\bigr)\,.
\end{align}
Now compute the constant term. Note:
\begin{align}
  \Aij\Xij^2 &= \frac{(\alpha_i x_i + \alpha_j x_j)^2}{\Aij}
  = \frac{\alpha_i^2 x_i^2 + 2\alpha_i\alpha_j x_i x_j + \alpha_j^2 x_j^2}{\alpha_i+\alpha_j}\,.
\end{align}
Therefore:
\begin{align}
  &\alpha_i x_i^2 + \alpha_j x_j^2 - \Aij\Xij^2\notag\\
  &= \frac{(\alpha_i x_i^2 + \alpha_j x_j^2)(\alpha_i+\alpha_j) - (\alpha_i x_i + \alpha_j x_j)^2}{\alpha_i+\alpha_j}\notag\\
  &= \frac{\alpha_i^2 x_i^2 + \alpha_i\alpha_j x_j^2 + \alpha_i\alpha_j x_i^2 + \alpha_j^2 x_j^2
       - \alpha_i^2 x_i^2 - 2\alpha_i\alpha_j x_i x_j - \alpha_j^2 x_j^2}{\alpha_i+\alpha_j}\notag\\
  &= \frac{\alpha_i\alpha_j(x_i^2 + x_j^2 - 2x_i x_j)}{\alpha_i+\alpha_j}
  = \frac{\alpha_i\alpha_j(x_i-x_j)^2}{\Aij} = \Qij\,.
\end{align}
Hence $f(x) = \Aij(x-\Xij)^2 + \Qij$, and $G_i G_j = e^{-f(x)} = e^{-\Qij}e^{-\Aij(x-\Xij)^2}$. $\square$

\textbf{Physical interpretation:} The product of two Gaussians is again a Gaussian (up to a
constant factor $e^{-\Qij}$), centered at the weighted average position $\Xij$ with combined
width $\Aij$. The pre-exponential factor $e^{-\Qij}$ quantifies how much the two Gaussians
overlap: it equals 1 when $x_i = x_j$ and decays exponentially as $|x_i - x_j|$ increases.

% ------------------------------------------------------------
\subsection{Overlap integral $S_{ij}$}\label{sec:Gauss_overlap}
% ------------------------------------------------------------

\begin{equation}
  S_{ij} = \int_{-\infty}^{\infty}G_i(x)\,G_j(x)\,dx
  = e^{-\Qij}\int_{-\infty}^{\infty}e^{-\Aij(x-\Xij)^2}\,dx\,.
\end{equation}
Using the standard Gaussian integral $\int_{-\infty}^{\infty}e^{-Au^2}\,du = \sqrt{\pi/A}$ with
the substitution $u = x - \Xij$:
\begin{equation}\label{eq:Gauss_S}
  \boxed{S_{ij} = \sqrt{\frac{\pi}{\Aij}}\,e^{-\Qij}\,.}
\end{equation}

\textbf{Diagonal case} ($i = j$, $x_i = x_j$, $\alpha_i = \alpha_j = \alpha$):
$\Aij = 2\alpha$, $\Qij = 0$, $S_{ii} = \sqrt{\pi/(2\alpha)}$.

% ------------------------------------------------------------
\subsection{Gaussian moments $I_n$}\label{sec:Gauss_moments}
% ------------------------------------------------------------

Define the moment integrals
\begin{equation}\label{eq:In_def}
  I_n(A, X) = \int_{-\infty}^{\infty}x^n\,e^{-A(x-X)^2}\,dx\,.
\end{equation}
These are the building blocks for all potential energy matrix elements.

\subsubsection{Derivation of the general formula}

Substitute $u = x - X$, so $x = u + X$ and $dx = du$:
\begin{equation}
  I_n = \int_{-\infty}^{\infty}(u+X)^n\,e^{-Au^2}\,du
  = \sum_{k=0}^{n}\binom{n}{k}X^{n-k}\int_{-\infty}^{\infty}u^k\,e^{-Au^2}\,du\,.
\end{equation}
The Gaussian moments of $u$ are:
\begin{equation}\label{eq:u_moments}
  \int_{-\infty}^{\infty}u^k\,e^{-Au^2}\,du =
  \begin{cases}
    0 & k\;\text{odd}\,,\\[4pt]
    \dfrac{(k-1)!!}{(2A)^{k/2}}\sqrt{\dfrac{\pi}{A}} & k\;\text{even}\,,
  \end{cases}
\end{equation}
where $(k-1)!! = 1\cdot 3\cdot 5\cdots(k-1)$ is the double factorial (with $(-1)!! = 1$ by convention).

\textbf{Derivation of \cref{eq:u_moments}:} The odd moments vanish by symmetry.
For even moments, starting from the base case
\begin{equation}
  \int_{-\infty}^{\infty}e^{-Au^2}\,du = \sqrt{\frac{\pi}{A}} = I_0(A,0)\,,
\end{equation}
differentiate $n$ times with respect to $-A$ (or equivalently, use the generating function):
\begin{equation}
  \int_{-\infty}^{\infty}u^{2k}e^{-Au^2}\,du
  = \left(-\frac{\partial}{\partial A}\right)^k\sqrt{\frac{\pi}{A}}
  = \frac{(2k-1)!!}{(2A)^k}\sqrt{\frac{\pi}{A}}\,.
\end{equation}
This can be verified by induction: writing $J_k = \int u^{2k}e^{-Au^2}du$, integration
by parts gives $J_k = \frac{2k-1}{2A}J_{k-1}$, confirming the recurrence.

Therefore, the general moment is:
\begin{equation}\label{eq:In_general}
  \boxed{I_n(A, X) = \sqrt{\frac{\pi}{A}}\sum_{\substack{k=0\\k\;\text{even}}}^{n}\binom{n}{k}\frac{(k-1)!!}{(2A)^{k/2}}\,X^{n-k}\,.}
\end{equation}

\subsubsection{Recursion relation}

A useful recursion can be derived from the substitution $x = (x - X) + X$:
\begin{align}
  I_n &= \int x^n e^{-A(x-X)^2}dx = \int [(x-X)+X]x^{n-1}e^{-A(x-X)^2}dx\notag\\
  &= \int(x-X)x^{n-1}e^{-A(x-X)^2}dx + X\,I_{n-1}\,.
\end{align}
For the first term, integrate by parts with $f = x^{n-1}$ and $dg = (x-X)e^{-A(x-X)^2}dx$,
giving $g = -\frac{1}{2A}e^{-A(x-X)^2}$:
\begin{align}
  \int(x-X)x^{n-1}e^{-A(x-X)^2}dx
  &= \underbrace{\left[-\frac{x^{n-1}}{2A}e^{-A(x-X)^2}\right]_{-\infty}^{\infty}}_{=\,0}
  + \frac{n-1}{2A}\int x^{n-2}e^{-A(x-X)^2}dx\notag\\
  &= \frac{n-1}{2A}\,I_{n-2}\,.
\end{align}
Therefore:
\begin{equation}\label{eq:In_recursion}
  \boxed{I_n = X\,I_{n-1} + \frac{n-1}{2A}\,I_{n-2}\,,\qquad n \geq 2\,.}
\end{equation}
Starting values: $I_0 = \sqrt{\pi/A}$, $I_1 = X\sqrt{\pi/A}$.

\subsubsection{Explicit tabulation of $I_0$ through $I_4$}

Applying the recursion:
\begin{alignat}{2}
  I_0 &= \sqrt{\frac{\pi}{A}}\,, &\qquad& \label{eq:I0}\\[6pt]
  I_1 &= X\,I_0 = X\sqrt{\frac{\pi}{A}}\,, && \label{eq:I1}\\[6pt]
  I_2 &= X\,I_1 + \frac{1}{2A}\,I_0
       = \left(X^2 + \frac{1}{2A}\right)\sqrt{\frac{\pi}{A}}\,, && \label{eq:I2}\\[6pt]
  I_3 &= X\,I_2 + \frac{2}{2A}\,I_1
       = X\left(X^2 + \frac{1}{2A}\right)\sqrt{\frac{\pi}{A}} + \frac{X}{A}\sqrt{\frac{\pi}{A}}\notag\\
      &= \left(X^3 + \frac{3X}{2A}\right)\sqrt{\frac{\pi}{A}}\,, && \label{eq:I3}\\[6pt]
  I_4 &= X\,I_3 + \frac{3}{2A}\,I_2
       = X\left(X^3 + \frac{3X}{2A}\right)\sqrt{\frac{\pi}{A}} + \frac{3}{2A}\left(X^2+\frac{1}{2A}\right)\sqrt{\frac{\pi}{A}}\notag\\
      &= \left(X^4 + \frac{3X^2}{A}\cdot\frac{1}{2}+\frac{3X^2}{2A}+\frac{3}{4A^2}\right)\sqrt{\frac{\pi}{A}}\notag\\
      &= \left(X^4 + \frac{6X^2}{2A} + \frac{3}{4A^2}\right)\sqrt{\frac{\pi}{A}}\notag\\
      &= \left(X^4 + \frac{3X^2}{A} + \frac{3}{4A^2}\right)\sqrt{\frac{\pi}{A}}\,. && \label{eq:I4}
\end{alignat}

Let us verify \cref{eq:I4} using \cref{eq:In_general}. The even values of $k$ in $\sum_{k=0,2,4}\binom{4}{k}\frac{(k-1)!!}{(2A)^{k/2}}X^{4-k}$:
\begin{align}
  k=0&:\quad \binom{4}{0}\cdot 1 \cdot X^4 = X^4\,,\notag\\
  k=2&:\quad \binom{4}{2}\cdot\frac{1}{2A}\cdot X^2 = \frac{6X^2}{2A} = \frac{3X^2}{A}\,,\notag\\
  k=4&:\quad \binom{4}{4}\cdot\frac{3!!}{(2A)^2} = \frac{3}{4A^2}\,.
\end{align}
Sum: $(X^4 + 3X^2/A + 3/(4A^2))\sqrt{\pi/A}$. Consistent. $\checkmark$

\begin{center}
\renewcommand{\arraystretch}{1.8}
\begin{tabular}{cl}
  \toprule
  $n$ & $I_n(A,X) / \sqrt{\pi/A}$ \\
  \midrule
  0 & $1$ \\
  1 & $X$ \\
  2 & $X^2 + \dfrac{1}{2A}$ \\
  3 & $X^3 + \dfrac{3X}{2A}$ \\
  4 & $X^4 + \dfrac{3X^2}{A} + \dfrac{3}{4A^2}$ \\
  \bottomrule
\end{tabular}
\end{center}

% ------------------------------------------------------------
\subsection{Potential energy matrix elements: symmetric case}\label{sec:Gauss_V_sym}
% ------------------------------------------------------------

For $V(x) = ax^4 - bx^2$:
\begin{align}
  V_{ij} &= \int_{-\infty}^{\infty}G_i(x)\,V(x)\,G_j(x)\,dx
  = a\int x^4\,G_i G_j\,dx - b\int x^2\,G_i G_j\,dx\notag\\
  &= e^{-\Qij}\bigl[a\,I_4(\Aij,\Xij) - b\,I_2(\Aij,\Xij)\bigr]\,.
\end{align}
Substituting the explicit forms of $I_4$ and $I_2$:
\begin{align}\label{eq:Vij_sym}
  V_{ij} &= e^{-\Qij}\sqrt{\frac{\pi}{\Aij}}\left\{
    a\left[\Xij^4 + \frac{3\Xij^2}{\Aij} + \frac{3}{4\Aij^2}\right]
    - b\left[\Xij^2 + \frac{1}{2\Aij}\right]
  \right\}\,.
\end{align}
\begin{equation}\label{eq:Vij_sym_boxed}
  \boxed{V_{ij} = e^{-\Qij}\sqrt{\frac{\pi}{\Aij}}\left\{
    a\Xij^4 + \frac{3a\Xij^2}{\Aij} + \frac{3a}{4\Aij^2}
    - b\Xij^2 - \frac{b}{2\Aij}
  \right\}}
\end{equation}

\textbf{Dimensional verification.} In atomic units:
$[a] = E_h/a_0^4$, $[b] = E_h/a_0^2$, $[\alpha] = a_0^{-2}$, $[x] = a_0$.
Then $[\Xij] = a_0$, $[\Aij] = a_0^{-2}$, $[\Qij] = \text{dimensionless}$, and
$[\sqrt{\pi/\Aij}] = a_0$. Check: $[a\Xij^4\sqrt{\pi/\Aij}] = (E_h/a_0^4)(a_0^4)(a_0) = E_h\cdot a_0$.
Wait --- this should have dimensions of energy times length (since $V_{ij}$ is an integral
over $x$ of an energy density times a dimensionless integrand). Indeed, $V_{ij}$ as defined
is not yet an energy: it must be combined with the overlap to form a ratio, or the basis
functions must be normalized. Since we use unnormalized Gaussians, $V_{ij}$ has dimensions
$[V_{ij}] = E_h \cdot a_0$ (energy $\times$ length), consistent with $S_{ij} = \sqrt{\pi/\Aij}\,e^{-\Qij}$
also having dimensions of length. The ratio $V_{ij}/S_{ij}$ has dimensions of energy, as required.

\textbf{Diagonal case} ($i = j$, $x_i = x_j$, $\alpha_i = \alpha_j = \alpha$):
$\Aij = 2\alpha$, $\Xij = x_i$, $\Qij = 0$. Then:
\begin{equation}
  V_{ii} = \sqrt{\frac{\pi}{2\alpha}}\left\{
    a\,x_i^4 + \frac{3a\,x_i^2}{2\alpha} + \frac{3a}{16\alpha^2}
    - b\,x_i^2 - \frac{b}{4\alpha}
  \right\}\,.
\end{equation}

% ------------------------------------------------------------
\subsection{Potential energy matrix elements: asymmetric case}\label{sec:Gauss_V_asym}
% ------------------------------------------------------------

For $V(x) = ax^4 - bx^2 + cx$:
\begin{equation}\label{eq:Vij_asym}
  V_{ij} = e^{-\Qij}\bigl[a\,I_4(\Aij,\Xij) - b\,I_2(\Aij,\Xij) + c\,I_1(\Aij,\Xij)\bigr]\,.
\end{equation}
Using $I_1 = \Xij\sqrt{\pi/\Aij}$:
\begin{equation}
  \boxed{V_{ij} = e^{-\Qij}\sqrt{\frac{\pi}{\Aij}}\left\{
    a\Xij^4 + \frac{3a\Xij^2}{\Aij} + \frac{3a}{4\Aij^2}
    - b\Xij^2 - \frac{b}{2\Aij}
    + c\Xij
  \right\}}
\end{equation}

% ------------------------------------------------------------
\subsection{Kinetic energy matrix elements $T_{ij}$}\label{sec:Gauss_kinetic}
% ------------------------------------------------------------

The kinetic energy matrix element in atomic units ($\hb = 1$) is
\begin{equation}
  T_{ij} = -\frac{1}{2\mu}\int_{-\infty}^{\infty}G_i(x)\,\frac{d^2}{dx^2}G_j(x)\,dx\,.
\end{equation}

\subsubsection{Step 1: Compute the second derivative of $G_j$}

\begin{align}
  \frac{d}{dx}G_j(x) &= -2\alpha_j(x-x_j)\,G_j(x)\,,\\
  \frac{d^2}{dx^2}G_j(x) &= \frac{d}{dx}\bigl[-2\alpha_j(x-x_j)G_j(x)\bigr]\notag\\
  &= -2\alpha_j\,G_j(x) + 4\alpha_j^2(x-x_j)^2\,G_j(x)\notag\\
  &= \bigl[-2\alpha_j + 4\alpha_j^2(x-x_j)^2\bigr]G_j(x)\,.
\end{align}
Therefore:
\begin{equation}\label{eq:d2Gj}
  \frac{d^2}{dx^2}G_j(x) = \bigl[-2\alpha_j + 4\alpha_j^2(x-x_j)^2\bigr]G_j(x)\,.
\end{equation}

\subsubsection{Step 2: Assemble the integral}

\begin{align}
  T_{ij} &= -\frac{1}{2\mu}\int G_i(x)\bigl[-2\alpha_j + 4\alpha_j^2(x-x_j)^2\bigr]G_j(x)\,dx\notag\\
  &= -\frac{1}{2\mu}\left\{-2\alpha_j\int G_i G_j\,dx + 4\alpha_j^2\int(x-x_j)^2\,G_i G_j\,dx\right\}\notag\\
  &= \frac{\alpha_j}{\mu}\,S_{ij} - \frac{2\alpha_j^2}{\mu}\int(x-x_j)^2\,G_i G_j\,dx\,.
\end{align}

\subsubsection{Step 3: Evaluate $\int(x-x_j)^2\,G_i\,G_j\,dx$}

Using the Gaussian product formula, $G_i G_j = e^{-\Qij}e^{-\Aij(x-\Xij)^2}$.
Write $x - x_j = (x - \Xij) + (\Xij - x_j)$. Define
\begin{equation}\label{eq:Delta_j}
  \delta_j \equiv \Xij - x_j = \frac{\alpha_i x_i + \alpha_j x_j}{\Aij} - x_j
  = \frac{\alpha_i(x_i - x_j)}{\Aij}\,.
\end{equation}
Then:
\begin{align}
  (x-x_j)^2 &= (x-\Xij)^2 + 2\delta_j(x-\Xij) + \delta_j^2\,.
\end{align}
Integrating against $e^{-\Aij(x-\Xij)^2}$ (substituting $u = x - \Xij$):
\begin{align}
  &\int(x-x_j)^2\,e^{-\Aij(x-\Xij)^2}\,dx\notag\\
  &= \int u^2\,e^{-\Aij u^2}\,du + 2\delta_j\underbrace{\int u\,e^{-\Aij u^2}\,du}_{=\,0} + \delta_j^2\int e^{-\Aij u^2}\,du\notag\\
  &= \frac{1}{2\Aij}\sqrt{\frac{\pi}{\Aij}} + \delta_j^2\sqrt{\frac{\pi}{\Aij}}\notag\\
  &= \left(\frac{1}{2\Aij} + \delta_j^2\right)\sqrt{\frac{\pi}{\Aij}}\,.
\end{align}
Multiplying by $e^{-\Qij}$:
\begin{equation}
  \int(x-x_j)^2\,G_i G_j\,dx = e^{-\Qij}\sqrt{\frac{\pi}{\Aij}}\left(\frac{1}{2\Aij}+\delta_j^2\right)
  = S_{ij}\left(\frac{1}{2\Aij}+\delta_j^2\right)\,.
\end{equation}

\subsubsection{Step 4: Combine}

\begin{align}
  T_{ij} &= \frac{\alpha_j}{\mu}\,S_{ij} - \frac{2\alpha_j^2}{\mu}\,S_{ij}\left(\frac{1}{2\Aij}+\delta_j^2\right)\notag\\
  &= \frac{\alpha_j}{\mu}\,S_{ij}\left[1 - \frac{\alpha_j}{\Aij} - 2\alpha_j\delta_j^2\right]\,.
\end{align}
Substituting $\delta_j = \alpha_i(x_i-x_j)/\Aij$ gives $\delta_j^2 = \alpha_i^2(x_i-x_j)^2/\Aij^2$.
Therefore:
\begin{equation}\label{eq:Tij_Gauss}
  \boxed{T_{ij} = \frac{\alpha_j}{\mu}\,S_{ij}\left[1 - \frac{\alpha_j}{\Aij} - \frac{2\alpha_j\alpha_i^2(x_i-x_j)^2}{\Aij^2}\right]}
\end{equation}

Equivalently, using $\Qij = \alpha_i\alpha_j(x_i-x_j)^2/\Aij$:
\begin{equation}
  T_{ij} = \frac{\alpha_j}{\mu}\,S_{ij}\left[1 - \frac{\alpha_j}{\Aij} - \frac{2\alpha_i\Qij}{\Aij}\right]\,.
\end{equation}

Or substituting $S_{ij} = \sqrt{\pi/\Aij}\,e^{-\Qij}$:
\begin{equation}\label{eq:Tij_explicit}
  T_{ij} = \frac{\alpha_j}{\mu}\,e^{-\Qij}\sqrt{\frac{\pi}{\Aij}}
  \left[1 - \frac{\alpha_j}{\Aij} - \frac{2\alpha_j\alpha_i^2(x_i-x_j)^2}{\Aij^2}\right]\,.
\end{equation}

\subsubsection{Verification: diagonal case}

For $i = j$ (same center, same width): $\alpha_i = \alpha_j = \alpha$, $x_i = x_j$,
$\Aij = 2\alpha$, $\Qij = 0$, $\delta_j = 0$.
\begin{equation}
  T_{ii} = \frac{\alpha}{\mu}\sqrt{\frac{\pi}{2\alpha}}\left[1 - \frac{\alpha}{2\alpha} - 0\right]
  = \frac{\alpha}{\mu}\sqrt{\frac{\pi}{2\alpha}}\cdot\frac{1}{2}
  = \frac{\alpha}{2\mu}\,S_{ii}\,.
\end{equation}
\textbf{Check:} For a single Gaussian $G(x) = e^{-\alpha x^2}$, the expectation value of kinetic
energy is $\langle T\rangle = T_{ii}/S_{ii} = \alpha/(2\mu)$. This equals $p^2/(2\mu)$ with
$\langle p^2\rangle = \alpha$ (the variance of the momentum distribution for a Gaussian
of width $1/\sqrt{2\alpha}$ in position space is $\Delta p^2 = \alpha$,
and $\langle p\rangle = 0$, so $\langle p^2\rangle = \alpha$).
This is correct for a minimum-uncertainty state with $\Delta x = 1/(2\sqrt{\alpha})$ and
$\Delta p = \sqrt{\alpha}$, satisfying $\Delta x\Delta p = 1/2 = \hb/2$. $\checkmark$

\subsubsection{Alternative derivation via integration by parts}

One may also compute $T_{ij}$ using
\begin{equation}
  T_{ij} = \frac{1}{2\mu}\int\frac{dG_i}{dx}\frac{dG_j}{dx}\,dx\,,
\end{equation}
where $dG_i/dx = -2\alpha_i(x-x_i)G_i(x)$. This gives:
\begin{equation}
  T_{ij} = \frac{2\alpha_i\alpha_j}{\mu}\int(x-x_i)(x-x_j)\,G_i G_j\,dx\,.
\end{equation}
Writing $(x-x_i)(x-x_j) = (x-\Xij+\Xij-x_i)(x-\Xij+\Xij-x_j) = (u+\delta_i)(u+\delta_j)$
where $u = x - \Xij$, $\delta_i = \Xij - x_i$, $\delta_j = \Xij - x_j$:
\begin{align}
  &\int(u+\delta_i)(u+\delta_j)e^{-\Aij u^2}\,du
  = \frac{1}{2\Aij}\sqrt{\frac{\pi}{\Aij}} + \delta_i\delta_j\sqrt{\frac{\pi}{\Aij}}\,.
\end{align}
Therefore:
\begin{equation}
  T_{ij} = \frac{2\alpha_i\alpha_j}{\mu}\,e^{-\Qij}\sqrt{\frac{\pi}{\Aij}}\left(\frac{1}{2\Aij}+\delta_i\delta_j\right)\,.
\end{equation}
Since $\delta_i = \alpha_j(x_j-x_i)/\Aij$ and $\delta_j = \alpha_i(x_i-x_j)/\Aij$,
we have $\delta_i\delta_j = -\alpha_i\alpha_j(x_i-x_j)^2/\Aij^2$:
\begin{equation}
  T_{ij} = \frac{2\alpha_i\alpha_j}{\mu}\,S_{ij}\left(\frac{1}{2\Aij} - \frac{\alpha_i\alpha_j(x_i-x_j)^2}{\Aij^2}\right)\,.
\end{equation}
This is equivalent to \cref{eq:Tij_Gauss}, as can be verified by algebraic simplification.
For equal widths $\alpha_i = \alpha_j = \alpha$: $\delta_i = -\delta_j$, $\delta_i\delta_j = -\delta_j^2$, and
\begin{equation}
  T_{ij}^{(\alpha_i=\alpha_j=\alpha)} = \frac{2\alpha^2}{\mu}\,S_{ij}\left(\frac{1}{4\alpha} - \frac{(x_i-x_j)^2\alpha^2}{4\alpha^2}\right)
  = \frac{\alpha}{2\mu}\,S_{ij}\left(1 - \frac{\alpha(x_i-x_j)^2}{2}\right)\,,
\end{equation}
which from \cref{eq:Tij_Gauss} with $\Aij = 2\alpha$:
$\frac{\alpha}{\mu}\left[1 - \frac{1}{2} - \frac{2\alpha^3(x_i-x_j)^2}{4\alpha^2}\right]
= \frac{\alpha}{\mu}\left[\frac{1}{2} - \frac{\alpha(x_i-x_j)^2}{2}\right]
= \frac{\alpha}{2\mu}\left[1 - \alpha(x_i-x_j)^2\right]$. This matches (note $\Qij = \alpha(x_i-x_j)^2/2$ for equal widths). $\checkmark$

% ------------------------------------------------------------
\subsection{L\"owdin orthogonalization for the Gaussian basis}\label{sec:Gauss_Lowdin}
% ------------------------------------------------------------

The Gaussian basis functions are not orthogonal ($S_{ij} \neq \delta_{ij}$ in general).
The overlap matrix can become nearly singular when Gaussians are placed too close together
or have very different widths. The L\"owdin procedure (\cref{sec:Lowdin}) handles this:

\begin{enumerate}
  \item Diagonalize $\mathbf{S}$: $\mathbf{S} = \mathbf{U}\,\bm{\sigma}\,\mathbf{U}^T$.
  \item Discard eigenvectors with $\sigma_i < \eps_{\text{thresh}}$ (typically $10^{-8}$).
  \item Construct $\mathbf{X} = \mathbf{U}'(\bm{\sigma}')^{-1/2}$.
  \item Transform: $\tilde{\mathbf{H}} = \mathbf{X}^T(\mathbf{T}+\mathbf{V})\mathbf{X}$.
  \item Diagonalize $\tilde{\mathbf{H}}$ to get energies and eigenvectors in the orthogonalized basis.
  \item Back-transform: $\mathbf{c} = \mathbf{X}\tilde{\mathbf{c}}$.
\end{enumerate}

The reduced dimension $N' < N$ implies that some linearly dependent combinations of
Gaussians have been projected out. This is physically benign provided $N'$ is still
large enough for convergence.

% ------------------------------------------------------------
\subsection{Parameter selection for the double well}\label{sec:Gauss_params}
% ------------------------------------------------------------

\subsubsection{Centers: uniform grid}

A simple choice is a uniform grid of $N$ centers spanning the classically accessible region:
\begin{equation}
  x_i = x_{\min}^{\text{grid}} + (i-1)\,\Delta x\,,\qquad i = 1, \ldots, N\,,
\end{equation}
where $\Delta x = (x_{\max}^{\text{grid}} - x_{\min}^{\text{grid}})/(N-1)$.
The grid boundaries $x_{\min,\max}^{\text{grid}}$ should extend well beyond the
classical turning points of the highest energy level of interest, typically
$x_{\max}^{\text{grid}} \approx 2$--$3\,\xe$.

\subsubsection{Common width parameter}

A natural choice is the local harmonic width at the potential minima:
\begin{equation}
  \alpha = \alpha_{\text{local}} = \mu\,\omega_{\text{local}} = 2\sqrt{\mu\,b}\,,
\end{equation}
or alternatively $\alpha = \mu\omega_{\text{local}}/2 = \sqrt{\mu\,b}$.
The optimal width can also be determined variationally by minimizing the ground-state
energy with respect to a common $\alpha$.

\subsubsection{Number of Gaussians}

The required number depends on the grid spacing relative to the Gaussian width:
$\Delta x \lesssim 1/\sqrt{2\alpha}$ ensures sufficient overlap between adjacent Gaussians.
For the double-well with typical parameters, $N \sim 20$--$60$ Gaussians suffice
for spectroscopic accuracy of the lowest 4--8 levels.

% ============================================================
% ============================================================

\section{Basis Set~V: Particle-in-a-Box Basis (PIB-FBR)}\label{sec:PIB}
% ============================================================
% ============================================================

The particle-in-a-box (PIB) basis is a Finite Basis Representation (FBR) built from
the exact eigenfunctions of the kinetic energy operator on a bounded domain.
It possesses three properties that make it particularly well suited to the double-well
tunneling problem.  First, the basis is \emph{complete} on $[-L/2, L/2]$ in the
$L^2$ sense, and eigenvalues converge to the exact values from above as the truncation
size $N\to\infty$ (variational theorem).  Second, the only free parameter is the
box half-length $L/2$, which plays a purely geometric role with no reference to
any equilibrium geometry of the potential.  Third, for any polynomial potential
$V(x) = \sum_k c_k x^k$, all Hamiltonian matrix elements are obtainable in
\emph{closed form} via elementary integration, eliminating quadrature entirely.
The last property is especially significant for the asymmetric double well
$V(x) = ax^4 - bx^2 + cx$ ($c \neq 0$): unlike the shifted harmonic oscillator
basis (Basis~II), the PIB basis requires no recentering or modification when $c$
changes.  The close relationship between the PIB-FBR and the sinc-DVR (Basis~IV,
\cref{sec:DVR}) is made precise in \cref{sec:PIB_DVR}.

% ------------------------------------------------------------
\subsection{PIB eigenfunctions and eigenvalues}\label{sec:PIB_wf}
% ------------------------------------------------------------

\subsubsection{Schrödinger equation and boundary conditions}

Consider a particle of unit mass confined to the segment $[-L/2, L/2]$ by an
infinite square well.  In atomic units ($\hb = 1$) the time-independent
Schrödinger equation inside the box is
\begin{equation}\label{eq:PIB_TISE}
  \hat{T}\,\phi(x) = -\,\frac{1}{2}\,\frac{d^2\phi}{dx^2} = E\,\phi(x)\,,
\end{equation}
subject to Dirichlet boundary conditions
\begin{equation}\label{eq:PIB_BC}
  \phi\!\left(-\tfrac{L}{2}\right) = \phi\!\left(\tfrac{L}{2}\right) = 0\,.
\end{equation}
Setting $k = \sqrt{2E}$ ($E > 0$), the general solution is
$\phi(x) = A\cos(kx) + B\sin(kx)$.
Evaluating \cref{eq:PIB_BC} at $x = \pm L/2$ and adding/subtracting the two equations
gives two independent conditions:
\begin{equation}
  2A\cos\!\left(\tfrac{kL}{2}\right) = 0\,,\qquad
  2B\sin\!\left(\tfrac{kL}{2}\right) = 0\,.
\end{equation}
Non-trivial solutions require either $A = 0$ or $B = 0$.

\textbf{Family~1 (even parity, $B = 0$):}
$\cos(kL/2) = 0 \Rightarrow kL/2 = (2j-1)\pi/2$, so $k_n = n\pi/L$ with $n = 1, 3, 5, \ldots$

\textbf{Family~2 (odd parity, $A = 0$):}
$\sin(kL/2) = 0 \Rightarrow kL/2 = j\pi$, so $k_n = n\pi/L$ with $n = 2, 4, 6, \ldots$

Both families share the same wavenumber formula $k_n = n\pi/L$, differing only in
which trigonometric function appears.

\subsubsection{Unified indexing and normalized eigenfunctions}

Ordering all solutions by increasing $k_n$, we label them with a single quantum
number $n = 1, 2, 3, \ldots$\,:
\begin{equation}\label{eq:PIB_basis}
  \boxed{\phi_n(x) =
  \begin{cases}
    \displaystyle\sqrt{\frac{2}{L}}\cos\!\left(\frac{n\pi x}{L}\right),
    & n = 1, 3, 5, \ldots\;\text{(even parity)}\,,\\[10pt]
    \displaystyle\sqrt{\frac{2}{L}}\sin\!\left(\frac{n\pi x}{L}\right),
    & n = 2, 4, 6, \ldots\;\text{(odd parity)}\,,
  \end{cases}}
\end{equation}
defined on $[-L/2, L/2]$ and zero outside.  The parity of $\phi_n$ under $x \to -x$
is $(-1)^{n+1}$:
\begin{equation}\label{eq:PIB_parity}
  \phi_n(-x) = (-1)^{n+1}\,\phi_n(x)\,.
\end{equation}
The normalization is verified directly.  For odd $n$,
\begin{align}
  \int_{-L/2}^{L/2}\!\bigl|\phi_n(x)\bigr|^2\,dx
  &= \frac{2}{L}\int_{-L/2}^{L/2}\cos^2\!\left(\frac{n\pi x}{L}\right)dx\notag\\
  &= \frac{2}{L}\cdot\frac{1}{2}\int_{-L/2}^{L/2}\!\left[1 + \cos\!\left(\frac{2n\pi x}{L}\right)\right]dx\notag\\
  &= \frac{1}{L}\left[L + \frac{L}{2n\pi}\,\sin\!\left(\frac{2n\pi x}{L}\right)\bigg|_{-L/2}^{L/2}\right]
  = 1\,,\quad\checkmark
\end{align}
since $\sin(n\pi) = 0$ for all integer $n$; the argument for even $n$ is identical.

\subsubsection{Orthonormality and eigenvalues}

The eigenfunctions satisfy
\begin{equation}\label{eq:PIB_ortho}
  \braket{\phi_m}{\phi_n} = \int_{-L/2}^{L/2}\phi_m(x)\,\phi_n(x)\,dx = \delta_{mn}\,.
\end{equation}
Orthogonality within the same family follows from the standard trigonometric
identity $\int_{-a}^{a}\cos(p\pi x/a)\cos(q\pi x/a)\,dx = a\,\delta_{pq}$; for
the mixed case (one cosine, one sine) the integrand is odd over a symmetric interval
and vanishes identically.  The exact zeroth-order energies are:
\begin{equation}\label{eq:PIB_En}
  \boxed{E_n^{(0)} = \frac{k_n^2}{2} = \frac{n^2\pi^2}{2L^2}\,,\quad n = 1, 2, 3, \ldots}
\end{equation}

% ------------------------------------------------------------
\subsection{Kinetic energy matrix elements}\label{sec:PIB_T}
% ------------------------------------------------------------

Because each $\phi_n$ is an eigenfunction of $\hat{T}$ with eigenvalue $E_n^{(0)}$,
the kinetic energy matrix in the PIB basis is exactly diagonal.  To see this, note
that for odd $n$:
\begin{equation}
  -\frac{1}{2}\frac{d^2}{dx^2}\cos\!\left(\frac{n\pi x}{L}\right)
  = \frac{n^2\pi^2}{2L^2}\cos\!\left(\frac{n\pi x}{L}\right)\,,
\end{equation}
and analogously for even $n$.  Therefore:
\begin{align}
  T_{mn} = \mel{\phi_m}{\hat{T}}{\phi_n}
  &= E_n^{(0)}\braket{\phi_m}{\phi_n}\notag\\
  &= E_n^{(0)}\,\delta_{mn}\,.
\end{align}
\begin{equation}\label{eq:PIB_T}
  \boxed{T_{mn} = \frac{n^2\pi^2}{2L^2}\,\delta_{mn}\,.}
\end{equation}
No numerical integration is ever required for the kinetic energy: the diagonal
elements are given by \cref{eq:PIB_En} and all off-diagonal elements are
identically zero.  This contrasts sharply with the DVR (Basis~IV, \cref{sec:DVR}),
where $\hat{T}$ is represented by a dense matrix in the position grid; the two
representations are related by the FBR$\to$DVR transformation discussed in
\cref{sec:PIB_DVR}.

% ------------------------------------------------------------
\subsection{Potential energy matrix elements}\label{sec:PIB_V}
% ------------------------------------------------------------

For any polynomial potential $V(x) = \sum_k c_k\,x^k$, the potential matrix
decomposes as
\begin{equation}
  V_{mn} = \mel{\phi_m}{V(x)}{\phi_n} = \sum_k c_k\,I_{mn}^{(k)}\,,
\end{equation}
where the fundamental integrals are
\begin{equation}\label{eq:PIB_Imn}
  I_{mn}^{(k)} \equiv \int_{-L/2}^{L/2}\phi_m(x)\,x^k\,\phi_n(x)\,dx\,.
\end{equation}
All $I_{mn}^{(k)}$ for polynomial $V(x)$ admit closed-form expressions obtained
by repeated integration by parts.  The following subsections derive these
expressions for $k = 1, 2, 3, 4$; all formulas have been independently verified
by symbolic integration with SymPy.

% ............................................................
\subsubsection{Parity selection rule}\label{sec:PIB_parity_sel}
% ............................................................

The integrand of \cref{eq:PIB_Imn} has parity
$(-1)^{m+1}\cdot(-1)^k\cdot(-1)^{n+1} = (-1)^{m+n+k}$ under $x \to -x$.
Since the integration domain is the symmetric interval $[-L/2, L/2]$, any odd
integrand integrates to zero:
\begin{equation}\label{eq:PIB_sel}
  \boxed{I_{mn}^{(k)} = 0\quad\text{if }m + n + k\text{ is odd.}}
\end{equation}
This \emph{parity selection rule} determines which pairs $(m,n)$ are coupled by
each power of $x$ before any algebra is performed:
\begin{itemize}
  \item $k$ odd (powers $x, x^3, \ldots$): couples states of \emph{opposite} parity ($m+n$ odd).
  \item $k$ even (powers $x^2, x^4, \ldots$): couples states of \emph{same} parity ($m+n$ even).
\end{itemize}

% ............................................................
\subsubsection{$k = 1$: the linear term $cx$}\label{sec:PIB_k1}
% ............................................................

By \cref{eq:PIB_sel}, $I_{mn}^{(1)} \neq 0$ only when $m + n$ is odd, i.e.\ one
index is even and the other odd.  Take $m$ odd and $n$ even (the case $m$ even,
$n$ odd follows by symmetry):
\begin{align}
  I_{mn}^{(1)}
  &= \frac{2}{L}\int_{-L/2}^{L/2}
    \cos\!\left(\frac{m\pi x}{L}\right)x\sin\!\left(\frac{n\pi x}{L}\right)dx\notag\\
  &= \frac{1}{L}\int_{-L/2}^{L/2}x\left[
    \sin\!\left(\frac{(m+n)\pi x}{L}\right) - \sin\!\left(\frac{(m-n)\pi x}{L}\right)
    \right]dx\,,\label{eq:PIB_k1_int}
\end{align}
where we used $\cos\alpha\sin\beta = \frac{1}{2}[\sin(\alpha+\beta)-\sin(\alpha-\beta)]$.
For any nonzero integer $p$, integration by parts gives
\begin{align}
  \int_{-L/2}^{L/2}x\sin\!\left(\frac{p\pi x}{L}\right)dx
  &= \left[-\frac{Lx}{p\pi}\cos\!\left(\frac{p\pi x}{L}\right)\right]_{-L/2}^{L/2}
    + \frac{L}{p\pi}\int_{-L/2}^{L/2}\cos\!\left(\frac{p\pi x}{L}\right)dx\notag\\
  &= -\frac{L^2}{p\pi}\cos\!\left(\frac{p\pi}{2}\right)
    + \frac{2L^2}{p^2\pi^2}\sin\!\left(\frac{p\pi}{2}\right)\,.
  \label{eq:PIB_k1_ibp}
\end{align}
For $m$ odd and $n$ even, both $p = m+n$ and $p = m-n$ are odd, so
$\cos(p\pi/2) = 0$ and $\sin(p\pi/2) = (-1)^{(p-1)/2}$.  Hence
\begin{equation}
  \int_{-L/2}^{L/2}x\sin\!\left(\frac{p\pi x}{L}\right)dx
  = \frac{2L^2}{p^2\pi^2}(-1)^{(p-1)/2}\,,\qquad p = m\pm n\,.
\end{equation}
Moreover, since $n$ is even, $(-1)^{(m-n-1)/2} = (-1)^{(m+n-1)/2}\cdot(-1)^n
= (-1)^{(m+n-1)/2}$.  Substituting into \cref{eq:PIB_k1_int} and using
$(m-n)^2 - (m+n)^2 = -4mn$:
\begin{equation}\label{eq:PIB_I1}
  \boxed{I_{mn}^{(1)} = \frac{-8mnL}{\pi^2(m^2-n^2)^2}\,(-1)^{(m+n-1)/2}\,,
  \quad m+n\text{ odd}\,,\;m\neq n\,.}
\end{equation}
The diagonal elements $I_{nn}^{(1)} = 0$ for all $n$, reflecting $\langle\phi_n|x|\phi_n\rangle = 0$
by parity on a symmetric domain.  The off-diagonal coupling decays as
$\sim mn/(m^2-n^2)^2 \sim 1/m^3$ for large $m$ at fixed $n$, ensuring absolute
convergence of the basis expansion.

% ............................................................
\subsubsection{$k = 2$: the quadratic term $-bx^2$}\label{sec:PIB_k2}
% ............................................................

By \cref{eq:PIB_sel}, $I_{mn}^{(2)} \neq 0$ only when $m+n$ is even.

\textbf{Diagonal elements ($m = n$).}
Using $f_n^2(x) = \frac{1}{2}\bigl[1 \pm \cos(2n\pi x/L)\bigr]$ for cosine/sine families:
\begin{equation}
  I_{nn}^{(2)} = \frac{1}{L}\left[\int_{-L/2}^{L/2}x^2\,dx
    \pm\int_{-L/2}^{L/2}x^2\cos\!\left(\frac{2n\pi x}{L}\right)dx\right]\,.
\end{equation}
The first integral gives $L^3/12$.  For the second, with $p = 2n$ (even):
double integration by parts yields $\int_{-L/2}^{L/2}x^2\cos(p\pi x/L)\,dx
= (-1)^n L^3/(2n^2\pi^2)$.  After inserting and noting that for odd $n$,
$(-1)^n = -1$, and for even $n$, the $\pm$ and $(-1)^n$ signs conspire to give
the same result in both cases:
\begin{equation}\label{eq:PIB_I2_diag}
  \boxed{I_{nn}^{(2)} = L^2\!\left(\frac{1}{12} - \frac{1}{2n^2\pi^2}\right).}
\end{equation}

\textbf{Off-diagonal elements ($m \neq n$, $m+n$ even).}
Both $m,n$ odd (cosine-cosine): $\cos(m\pi x/L)\cos(n\pi x/L)
= \frac{1}{2}[\cos((m-n)\pi x/L) + \cos((m+n)\pi x/L)]$.
Both $m,n$ even (sine-sine): $\sin(m\pi x/L)\sin(n\pi x/L)
= \frac{1}{2}[\cos((m-n)\pi x/L) - \cos((m+n)\pi x/L)]$.
In either case $p = m\pm n$ is even, say $p = 2q$, and
$\int_{-L/2}^{L/2}x^2\cos(p\pi x/L)\,dx = (-1)^q L^3/(2q^2\pi^2)$.
Defining $s \equiv (-1)^{m+1} = +1$ for both odd, $-1$ for both even:
\begin{equation}\label{eq:PIB_I2_offdiag}
  \boxed{I_{mn}^{(2)} = \frac{L^2}{\pi^2}\!\left[
    \frac{(-1)^{(m-n)/2}}{(m-n)^2}
    + s\,\frac{(-1)^{(m+n)/2}}{(m+n)^2}
  \right],\quad m\neq n,\;m+n\text{ even,}}
\end{equation}
where $s = (-1)^{m+1}$.

% ............................................................
\subsubsection{$k = 3$: the cubic term (for generality)}\label{sec:PIB_k3}
% ............................................................

Although the double-well potential $V = ax^4 - bx^2 + cx$ contains no cubic term,
we record the result for a general polynomial.  By \cref{eq:PIB_sel},
$I_{mn}^{(3)} \neq 0$ only for $m+n$ odd.
The derivation follows the same pattern as \cref{sec:PIB_k1}: write
$\cos(m\pi x/L)\sin(n\pi x/L) = \frac{1}{2}[\sin((m+n)\pi x/L)
- \sin((m-n)\pi x/L)]$ and evaluate $K_p = \int_{-L/2}^{L/2}x^3\sin(p\pi x/L)\,dx$
for odd $p$ via three integrations by parts.  The first IBP reduces $K_p$ to
$\frac{3L}{p\pi}J_p$, where $J_p = \int_{-L/2}^{L/2}x^2\cos(p\pi x/L)\,dx$
with odd $p$; two further IBPs give
$J_p = (-1)^{(p-1)/2}L^3(p^2\pi^2-8)/(2p^3\pi^3)$.
Assembling and simplifying using $(-1)^{(m-n-1)/2} = (-1)^{(m+n-1)/2}$ (as in
\cref{sec:PIB_k1}):
\begin{equation}\label{eq:PIB_I3}
  \boxed{I_{mn}^{(3)} = (-1)^{(m+n-1)/2}\,\frac{6mnL^3}{\pi^2(m^2-n^2)^2}
  \!\left[\frac{16(m^2+n^2)}{\pi^2(m^2-n^2)^2} - 1\right],\quad m+n\text{ odd.}}
\end{equation}

% ............................................................
\subsubsection{$k = 4$: the quartic term $ax^4$}\label{sec:PIB_k4}
% ............................................................

By \cref{eq:PIB_sel}, $I_{mn}^{(4)} \neq 0$ only for $m+n$ even.

\textbf{Diagonal elements.}
Using $f_n^2 = \frac{1}{2}[1\pm\cos(2n\pi x/L)]$, we need
$J_p^{(4)} = \int_{-L/2}^{L/2}x^4\cos(p\pi x/L)\,dx$ for $p = 2n$.
Four integrations by parts reduce this to previously evaluated integrals.
The first IBP gives $J_p^{(4)} = -(4L/p\pi)K_p$ (boundary vanishes since
$\sin(n\pi) = 0$), where $K_p = \int x^3\sin(p\pi x/L)\,dx$ with even $p$.
Evaluating $K_{2n}$ via IBP yields $K_{2n} = (-1)^n L^4(12-n^2\pi^2)/(8n^3\pi^3)$,
so $J_{2n}^{(4)} = (-1)^n L^5(n^2\pi^2-12)/(4n^4\pi^4)$.
The sign cancellation between the two parities gives a universal result:
\begin{equation}\label{eq:PIB_I4_diag}
  \boxed{I_{nn}^{(4)} = L^4\!\left(\frac{1}{80} - \frac{1}{4n^2\pi^2}
    + \frac{3}{n^4\pi^4}\right).}
\end{equation}

\textbf{Off-diagonal elements ($m \neq n$, $m+n$ even).}
For even $p$, $J_p^{(4)} = (-1)^{p/2}L^5(p^2\pi^2-48)/p^4\pi^4$.
With the same parity factor $s = (-1)^{m+1}$:
\begin{equation}\label{eq:PIB_I4_offdiag}
  \boxed{I_{mn}^{(4)} = \frac{L^4}{\pi^4}\!\left[
    \frac{(-1)^{(m-n)/2}\bigl[(m-n)^2\pi^2-48\bigr]}{(m-n)^4}
    +s\,\frac{(-1)^{(m+n)/2}\bigl[(m+n)^2\pi^2-48\bigr]}{(m+n)^4}
  \right],\quad m\neq n,\;m+n\text{ even.}}
\end{equation}

% ............................................................
\subsubsection{Summary table of selection rules and matrix elements}
% ............................................................

\begin{center}
\renewcommand{\arraystretch}{1.5}
\begin{tabular}{cclll}
  \toprule
  $k$ & Term in $V$ & Nonzero when & Diagonal & Off-diagonal \\
  \midrule
  1 & $cx$ & $m+n$ odd & $0$ & \cref{eq:PIB_I1} \\
  2 & $-bx^2$ & $m+n$ even & \cref{eq:PIB_I2_diag} & \cref{eq:PIB_I2_offdiag} \\
  3 & $cx^3$ & $m+n$ odd & $0$ & \cref{eq:PIB_I3} \\
  4 & $ax^4$ & $m+n$ even & \cref{eq:PIB_I4_diag} & \cref{eq:PIB_I4_offdiag} \\
  \bottomrule
\end{tabular}
\end{center}

\begin{remark}
All formulas in \crefrange{eq:PIB_I1}{eq:PIB_I4_offdiag} have been verified
independently by symbolic integration with SymPy for a representative set of
index pairs $(m,n)$ with $L = 10\,a_0$.  Agreement with direct numerical
quadrature is exact to machine precision in all cases tested.
\end{remark}

% ------------------------------------------------------------
\subsection{Full Hamiltonian for the double-well potential}\label{sec:PIB_H}
% ------------------------------------------------------------

For the potential $V(x) = ax^4 - bx^2 + cx$, the Hamiltonian matrix in the
PIB basis reads
\begin{equation}\label{eq:PIB_H}
  \boxed{H_{mn} = \frac{n^2\pi^2}{2L^2}\,\delta_{mn}
    + a\,I_{mn}^{(4)} - b\,I_{mn}^{(2)} + c\,I_{mn}^{(1)}\,.}
\end{equation}
The kinetic energy contributes only to the diagonal.  The potential energy
contributes to both the diagonal (through $I^{(2)}$ and $I^{(4)}$) and, when
$c \neq 0$, to off-diagonal blocks connecting states of opposite parity
(through $I^{(1)}$).

\subsubsection{Block structure for $c = 0$ (symmetric well)}

When $c = 0$, the Hamiltonian commutes with the parity operator $\hat{\Pi}$:
$[\hat{H}, \hat{\Pi}] = 0$.  The selection rule \cref{eq:PIB_sel} then
forbids all matrix elements between even-parity states ($n = 1, 3, 5, \ldots$)
and odd-parity states ($n = 2, 4, 6, \ldots$), so $H$ is block-diagonal:
\begin{equation}\label{eq:PIB_block_sym}
  H\big|_{c=0} =
  \begin{pmatrix}
    H^{\text{even}} & \mathbf{0} \\
    \mathbf{0} & H^{\text{odd}}
  \end{pmatrix}\!,
\end{equation}
where $H^{\text{even}}$ (resp.\ $H^{\text{odd}}$) couples only states with
odd (resp.\ even) quantum number.  The two sub-problems can be solved
independently, halving the computational cost.

\subsubsection{Block structure for $c \neq 0$ (asymmetric well)}

When $c \neq 0$, the linear term introduces the coupling block
$W_{mn} = c\,I_{mn}^{(1)}$ connecting the even and odd sectors:
\begin{equation}\label{eq:PIB_block_asym}
  H =
  \begin{pmatrix}
    H^{\text{even}} & W \\
    W^T & H^{\text{odd}}
  \end{pmatrix}\!,\qquad W_{mn} = c\,I_{mn}^{(1)}\neq 0\quad\text{for }m+n\text{ odd.}
\end{equation}
The full $N\times N$ matrix must be diagonalized.  However, since the elements
of $W$ decay as $|W_{mn}| \sim |c|\cdot mn/(m^2-n^2)^2 \to 0$ rapidly for
large $|m-n|$, the matrix $H$ is effectively banded in practice and convergence
is exponentially fast with $N$.

% ------------------------------------------------------------
\subsection{Analytical advantages for the asymmetric potential}\label{sec:PIB_asym}
% ------------------------------------------------------------

The asymmetric double well ($c \neq 0$) exposes a fundamental limitation of
the other basis sets that the PIB basis circumvents by construction.

\textbf{Shifted harmonic oscillator (Basis~II).}
The shifted-HO basis is centered at a chosen equilibrium point $x_e$, the
minimum of a reference parabolic potential.  For the symmetric well ($c = 0$),
placing two centers at $\pm x_e$ is natural and leads to a compact representation.
For $c \neq 0$, the two minima of the physical potential shift asymmetrically:
one becomes deeper and the other shallower.  The centers $\pm x_e$ are no
longer optimal, and the shifted-HO basis must be repositioned.  This requires
re-computing all matrix elements for every new value of $c$, and introduces
a non-trivial generalized eigenvalue problem $\mathbf{H}\mathbf{c} =
E\,\mathbf{S}\mathbf{c}$ because the shifted basis is non-orthogonal
(see \cref{sec:shifted_HO}).

\textbf{Harmonic oscillator (Basis~I).}
The centered HO basis at $x = 0$ exploits parity symmetry to block-diagonalize
$H$ for $c = 0$.  When $c \neq 0$, parity is broken and the block structure
collapses exactly as in \cref{eq:PIB_block_asym}.  The HO basis offers no
computational advantage and its Hermite-polynomial matrix elements become
increasingly costly to evaluate for high-order terms.

\textbf{PIB basis (Basis~V).}
No equilibrium point or reference frequency enters the PIB basis.  The box
$[-L/2, L/2]$ is positioned symmetrically about the origin and must only be
large enough to contain the physically relevant domain (see \cref{sec:PIB_conv}).
Consequently:
\begin{enumerate}
  \item \textbf{No re-centering.}  Changing $c$ only alters $V_{mn}$ through
        the closed-form element \cref{eq:PIB_I1}; the basis functions $\phi_n$
        and all other matrix elements are unchanged.
  \item \textbf{Uniform spatial resolution.}  Each $\phi_n$ is oscillatory
        throughout $[-L/2, L/2]$, providing equal coverage near both wells and
        in the barrier region regardless of $c$.  No single-well localization
        bias is introduced.
  \item \textbf{Standard eigenvalue problem.}  Because $\{\phi_n\}$ is
        orthonormal (\cref{eq:PIB_ortho}), the overlap matrix is the identity and
        $H_{mn}$ defines a standard (not generalized) eigenvalue problem for all $c$.
  \item \textbf{Rapid convergence.}  The off-diagonal coupling
        $|W_{mn}| \sim |c|\,mn/(m^2-n^2)^2$ decays as $m^{-3}$ for large $m$,
        guaranteeing that the additional fill-in caused by the asymmetry does not
        degrade the convergence rate.
\end{enumerate}

\begin{result}
For the asymmetric double well $V(x) = ax^4 - bx^2 + cx$ with $c \neq 0$,
the PIB-FBR provides a standard, orthonormal, fully analytic Hamiltonian that
requires no modification whatsoever as $c$ varies.  This makes it the natural
variational basis for systematic studies of the tunneling splitting as a function
of the asymmetry parameter.
\end{result}

% ------------------------------------------------------------
\subsection{Connection to the sinc-DVR (Basis~IV)}\label{sec:PIB_DVR}
% ------------------------------------------------------------

\subsubsection{The FBR$\to$DVR transformation}

Let $\{\phi_n\}_{n=1}^N$ be a truncated PIB-FBR.  Define $N$ uniformly spaced
grid points
\begin{equation}
  x_j = -\frac{L}{2} + j\,\Delta x\,,\quad j = 1, \ldots, N\,,\quad
  \Delta x = \frac{L}{N+1}\,.
\end{equation}
These are exactly the nodes of the $N+1$-point quadrature rule associated with
the PIB basis.  The transformation matrix
\begin{equation}\label{eq:PIB_U}
  U_{jn} = \sqrt{\Delta x}\;\phi_n(x_j)
\end{equation}
is orthogonal to quadrature accuracy: $U^T U = I$.  The DVR representation of
any operator $\hat{O}$ is $O^{\text{DVR}} = U\,O^{\text{FBR}}\,U^T$.

\subsubsection{Recovery of Colbert--Miller kinetic energy}

Since $T^{\text{FBR}}_{mn} = E_n^{(0)}\delta_{mn}$ (\cref{eq:PIB_T}), the
DVR kinetic energy is
\begin{equation}
  T^{\text{DVR}}_{jl}
  = \sum_{n=1}^{N} U_{jn}\,E_n^{(0)}\,U_{ln}
  = \Delta x\sum_{n=1}^{N}\phi_n(x_j)\,\frac{n^2\pi^2}{2L^2}\,\phi_n(x_l)\,.
\end{equation}
In the limit $N \to \infty$ with $\Delta x$ fixed, this sum converges to the
Colbert--Miller kinetic energy matrix~\cite{ColbertMiller1992}:
\begin{equation}
  T^{\text{DVR}}_{jj} \xrightarrow{N\to\infty}
  \frac{\pi^2}{6\,\Delta x^2}\cdot\frac{1}{2} \,=\,
  \frac{\hb^2\pi^2}{6\mu\,\Delta x^2}\bigg|_{\hb=\mu=1}\,,
\end{equation}
\begin{equation}
  T^{\text{DVR}}_{jl} \xrightarrow{N\to\infty}
  \frac{(-1)^{j-l}}{2(j-l)^2\Delta x^2}\,,\quad j\neq l\,,
\end{equation}
recovering exactly \cref{eq:T_DVR_diag,eq:T_DVR_offdiag} of \cref{sec:DVR}.
The DVR kinetic energy is thus the PIB kinetic energy expressed in the position
grid representation.

\subsubsection{Potential energy in DVR}

For the potential energy, the DVR property $\chi_j(x_l) = \delta_{jl}$ gives
$V^{\text{DVR}}_{jl} = V(x_j)\,\delta_{jl}$ (\cref{eq:V_DVR}), which is
\emph{diagonal} and requires only $N$ function evaluations.  In contrast, the
FBR-PIB potential matrix $V_{mn}^{\text{FBR}} = \sum_k c_k\,I_{mn}^{(k)}$ is
dense, but is given in closed form for any polynomial potential.

\subsubsection{Comparison of the two representations}

\begin{center}
\renewcommand{\arraystretch}{1.5}
\begin{tabular}{lll}
  \toprule
  Property & PIB-FBR & Sinc-DVR (Basis~IV) \\
  \midrule
  Basis functions & $\phi_n(x)$, analytic & $\delta$-like at grid points \\
  Overlap matrix & $\mathbf{I}$ (orthonormal) & $\mathbf{I}$ (orthonormal) \\
  Kinetic matrix $T$ & Diagonal, \cref{eq:PIB_T} & Dense, \cref{eq:T_DVR_diag,eq:T_DVR_offdiag} \\
  Potential matrix $V$ & Dense, closed form (polynomial $V$) & Diagonal, $V(x_j)\delta_{jl}$ \\
  Polynomial $V(x)$ & Exact, analytic & Exact (diagonal) \\
  Non-polynomial $V(x)$ & Quadrature needed & Diagonal, $V(x_j)\delta_{jl}$ \\
  Eigenvalue problem & Standard & Standard \\
  Convergence parameter & $N$ & $\Delta x = L/(N+1)$ \\
  Preferred when & $V$ is polynomial & $V$ is arbitrary \\
  \bottomrule
\end{tabular}
\end{center}

Both representations yield the same eigenvalues to quadrature accuracy for
fixed $N$ and $L$.  Their Hamiltonians are unitarily equivalent via \cref{eq:PIB_U}.

% ------------------------------------------------------------
\subsection{Box size and convergence}\label{sec:PIB_conv}
% ------------------------------------------------------------

\subsubsection{Choice of box half-length $L/2$}

The Dirichlet boundary conditions \cref{eq:PIB_BC} are exact only if the true
wavefunctions vanish at $x = \pm L/2$.  A necessary condition is therefore that
$x = \pm L/2$ lies deep in the classically forbidden region of the potential.
The practical criterion is
\begin{equation}\label{eq:PIB_L_crit}
  \bigl|\psi_n(\pm L/2)\bigr|^2 < \varepsilon\,,\qquad
  \varepsilon \sim 10^{-10}\text{--}10^{-12}\,.
\end{equation}
For the double-well potential $V(x) = ax^4 - bx^2 + cx$, the local minima of the
symmetric part lie at $x_{\min} = \pm\sqrt{b/(2a)}$.  A safe initial choice is
\begin{equation}\label{eq:PIB_L_dw}
  \frac{L}{2} \geq 2\sqrt{\frac{b}{a}}\,,
\end{equation}
adjusted upward for large $|c|$ which shifts the outer turning points.  The
box size should subsequently be validated by the $L$-convergence test of
\cref{sec:PIB_protocol}.

\subsubsection{Nyquist momentum criterion}

For a truncated basis of size $N$, the maximum resolvable momentum is
\begin{equation}\label{eq:PIB_kmax}
  k_{\max} = \frac{N\pi}{L}\,.
\end{equation}
To accurately represent all eigenstates with energy up to $E_{\max}$,
one requires $k_{\max} > \sqrt{2E_{\max}}$, i.e.\
\begin{equation}\label{eq:PIB_N_crit}
  N > \frac{L\sqrt{2E_{\max}}}{\pi}\,.
\end{equation}
For the ground-state doublet, $E_{\max}$ is typically of order $|V_{\min}|$
plus a few tunneling splittings, and $N \sim 20$--$50$ suffices with a
well-chosen $L$.

\subsubsection{Variational convergence}

Since $\{\phi_n\}_{n=1}^N$ is a subset of a complete orthonormal system, the
Ritz variational theorem guarantees that the $k$-th eigenvalue $E_k^{(N)}$ of
the truncated $N\times N$ Hamiltonian satisfies $E_k^{(N)} \geq E_k^{(\infty)}$
for all $N$, and $E_k^{(N)} \searrow E_k^{(\infty)}$ as $N \to \infty$.
Convergence is monotone from above and can be monitored directly.

\subsubsection{The $L$--$N$ trade-off}

For fixed computational cost (fixed $N$), the spacing $\Delta x = L/(N+1)$ grows
with $L$, reducing $k_{\max}$ and degrading the kinetic energy resolution.
Conversely, reducing $L$ below the criterion \cref{eq:PIB_L_dw} introduces
artificial boundary errors.  The optimal strategy is to choose $L$ just large
enough to satisfy \cref{eq:PIB_L_crit}, then increase $N$ until
\cref{eq:PIB_N_crit} is met.

\subsubsection{Practical convergence protocol}\label{sec:PIB_protocol}

A robust convergence protocol for the QuTu implementation is as follows:
\begin{enumerate}
  \item Choose $L/2 = 2\sqrt{b/a}$ as a starting point (\cref{eq:PIB_L_dw}).
  \item Compute eigenvalues for $N, 2N, 4N$ basis functions.
  \item Accept convergence when $|E_k^{(2N)} - E_k^{(N)}| < \delta_E$ for
        all states of interest ($\delta_E \sim 10^{-8}$\,a.u.).
  \item Validate $L$: repeat the converged calculation with $L \to 1.5L$ and
        verify eigenvalue stability.  If eigenvalues shift by more than $\delta_E$,
        increase $L$ and repeat from step~1.
\end{enumerate}
For the asymmetric double well ($c \neq 0$), the full $N \times N$ matrix is
diagonalized rather than two sub-blocks.  Since all matrix elements remain in
closed form for any $c$, the computational overhead relative to the symmetric
case is limited to the cost of additionally evaluating the elements
$I_{mn}^{(1)}$ given by \cref{eq:PIB_I1} and of solving the $N \times N$
(rather than two $N/2 \times N/2$) eigenvalue problem---a factor of four
increase in flop count that is wholly acceptable for the problem sizes of
interest.

% ============================================================
% ============================================================

\section{Basis Set~IV: DVR / Fourier Grid Hamiltonian}\label{sec:DVR}
% ============================================================
% ============================================================

The Discrete Variable Representation (DVR), specifically the Fourier Grid Hamiltonian (FGH)
method of Marston and Balint-Kurti (based on the Colbert--Miller sinc-DVR), provides a
conceptually simple and numerically robust approach to solving the Schr\"odinger equation.

% ------------------------------------------------------------
\subsection{Fourier Grid Hamiltonian: derivation}\label{sec:FGH}
% ------------------------------------------------------------

\subsubsection{Setup}

Consider the particle confined to a box $[-L/2, L/2]$ (with $L$ large enough that the
wavefunctions are negligible at the boundaries). Discretize the box with $N$ uniformly
spaced grid points:
\begin{equation}
  x_j = -\frac{L}{2} + j\,\Delta x\,,\qquad j = 1, 2, \ldots, N\,,
\end{equation}
where $\Delta x = L/(N+1)$ (particle-in-a-box boundary conditions) or $\Delta x = L/N$
(periodic boundary conditions). We use the latter convention.

The corresponding $k$-space grid is
\begin{equation}
  k_n = \frac{2\pi n}{L} = \frac{2\pi n}{N\,\Delta x}\,,\qquad
  n = -\frac{N-1}{2}, \ldots, \frac{N-1}{2}\quad\text{($N$ odd)}\,.
\end{equation}
For even $N$, $n = -N/2+1, \ldots, N/2$.

\subsubsection{Kinetic energy in momentum space}

The kinetic energy operator in momentum representation is diagonal:
\begin{equation}
  \hat{T}\ket{k_n} = \frac{\hb^2 k_n^2}{2\mu}\ket{k_n}\,.
\end{equation}

\subsubsection{Transformation to position space}

The discrete Fourier transform relates position and momentum representations:
\begin{equation}
  \braket{x_j}{k_n} = \frac{1}{\sqrt{N}}\,e^{ik_n x_j}\,.
\end{equation}
The kinetic energy matrix element in position space is:
\begin{align}\label{eq:T_DVR_sum}
  T_{jl} &= \mel{x_j}{\hat{T}}{x_l}
  = \sum_n \braket{x_j}{k_n}\frac{\hb^2 k_n^2}{2\mu}\braket{k_n}{x_l}\notag\\
  &= \frac{\hb^2}{2\mu N}\sum_n k_n^2\,e^{ik_n(x_j - x_l)}\,.
\end{align}

\subsubsection{Evaluation of the kinetic energy sum}

Define $p = j - l$ and $\phi = 2\pi/(N\Delta x)\cdot\Delta x = 2\pi/N$. Then
$k_n(x_j - x_l) = k_n\,p\,\Delta x = 2\pi np/N = np\phi$.

\textbf{Diagonal case ($p = 0$, $j = l$):}
\begin{equation}
  T_{jj} = \frac{\hb^2}{2\mu N}\sum_n k_n^2
  = \frac{\hb^2}{2\mu N}\cdot\frac{4\pi^2}{N^2\Delta x^2}\sum_n n^2\,.
\end{equation}
For $N$ odd, $n$ ranges over $-(N-1)/2$ to $(N-1)/2$:
\begin{equation}
  \sum_{n=-(N-1)/2}^{(N-1)/2}n^2 = 2\sum_{n=1}^{(N-1)/2}n^2
  = 2\cdot\frac{(N-1)/2\cdot((N-1)/2+1)\cdot(2(N-1)/2+1)}{6}
  = \frac{N(N^2-1)}{12}\,.
\end{equation}
For large $N$, $\sum n^2 \approx N^3/12$, so
\begin{equation}
  T_{jj} = \frac{\hb^2}{2\mu N}\cdot\frac{4\pi^2}{N^2\Delta x^2}\cdot\frac{N(N^2-1)}{12}
  = \frac{\hb^2\pi^2(N^2-1)}{6\mu N^2\Delta x^2}
  \xrightarrow{N\to\infty}
  \frac{\hb^2\pi^2}{6\mu\Delta x^2}\,.
\end{equation}

\textbf{Off-diagonal case ($p \neq 0$):}

We need to evaluate the sum
\begin{equation}
  \Sigma_p = \sum_{n=-(N-1)/2}^{(N-1)/2} n^2\,e^{i n p\phi}\,,\qquad \phi = 2\pi/N\,.
\end{equation}
Using the identity $n^2 e^{in\theta} = -\frac{d^2}{d\theta^2}e^{in\theta}$ and the geometric
series $\sum_n e^{in\theta} = \sin(N\theta/2)/\sin(\theta/2)$, one can show (after careful
differentiation and evaluation) that for $N$ odd and $p \neq 0$ mod $N$:
\begin{equation}
  \Sigma_p = (-1)^p\frac{N}{2}\cdot\frac{1}{\sin^2(p\pi/N)}\cdot\cos(p\pi/N)\cdot\frac{1}{1}\,.
\end{equation}
More precisely, using the exact result for the finite sum:
\begin{equation}
  T_{jl} = \frac{\hb^2}{2\mu\Delta x^2}\cdot\frac{(-1)^{j-l}}{N}\cdot\frac{\pi^2}{\sin^2\!\bigl(\pi(j-l)/N\bigr)}\cdot\frac{1}{3}\,.
\end{equation}

In the limit $N \to \infty$ (sinc-DVR), $\sin(\pi p/N) \to \pi p/N$, and the matrix elements
simplify to the well-known Colbert--Miller result.

% ------------------------------------------------------------
\subsection{Sinc-DVR (Colbert--Miller) kinetic energy}\label{sec:sinc_DVR}
% ------------------------------------------------------------

For an infinite uniform grid ($N \to \infty$) with spacing $\Delta x$, the kinetic energy
matrix elements are obtained by evaluating the integral
\begin{equation}
  T_{jl} = -\frac{\hb^2}{2\mu}\int_{-\infty}^{\infty}\text{sinc}\!\left(\frac{x-x_j}{\Delta x}\right)\frac{d^2}{dx^2}\text{sinc}\!\left(\frac{x-x_l}{\Delta x}\right)dx\,,
\end{equation}
where $\text{sinc}(u) = \sin(\pi u)/(\pi u)$ and $x_j = j\Delta x$.

The result, derived by Colbert and Miller (1992), is:

\textbf{Diagonal:}
\begin{equation}\label{eq:T_DVR_diag}
  \boxed{T_{jj} = \frac{\hb^2\pi^2}{6\mu\,\Delta x^2}\,.}
\end{equation}

\textbf{Off-diagonal ($j \neq l$):}
\begin{equation}\label{eq:T_DVR_offdiag}
  \boxed{T_{jl} = \frac{\hb^2}{\mu\,\Delta x^2}\cdot\frac{(-1)^{j-l}}{(j-l)^2}\,.}
\end{equation}

\textbf{Derivation.} Consider the Fourier representation of the sinc function:
\begin{equation}
  \text{sinc}\!\left(\frac{x - x_j}{\Delta x}\right) = \frac{\Delta x}{2\pi}\int_{-\pi/\Delta x}^{\pi/\Delta x}e^{ik(x-x_j)}\,dk\,.
\end{equation}
The second derivative of the sinc function centered at $x_l$ gives:
\begin{equation}
  \frac{d^2}{dx^2}\text{sinc}\!\left(\frac{x-x_l}{\Delta x}\right)
  = \frac{\Delta x}{2\pi}\int_{-\pi/\Delta x}^{\pi/\Delta x}(-k^2)\,e^{ik(x-x_l)}\,dk\,.
\end{equation}
The matrix element becomes:
\begin{align}
  T_{jl} &= \frac{\hb^2}{2\mu}\left(\frac{\Delta x}{2\pi}\right)^2
  \int_{-\infty}^{\infty}dx\int dk\int dk'\,k'^2\,e^{ik(x-x_j)}\,e^{ik'(x-x_l)}\notag\\
  &= \frac{\hb^2}{2\mu}\cdot\frac{\Delta x}{2\pi}\int_{-\pi/\Delta x}^{\pi/\Delta x}k^2\,e^{ik(x_j-x_l)}\,dk\,,
\end{align}
where we used the orthogonality $\int e^{i(k-k')x}dx = 2\pi\delta(k-k')$.
Setting $x_j - x_l = (j-l)\Delta x \equiv p\,\Delta x$:
\begin{equation}
  T_{jl} = \frac{\hb^2}{2\mu}\cdot\frac{\Delta x}{2\pi}\int_{-\pi/\Delta x}^{\pi/\Delta x}k^2\,e^{ikp\Delta x}\,dk\,.
\end{equation}
Substituting $u = k\Delta x$:
\begin{equation}\label{eq:T_integral}
  T_{jl} = \frac{\hb^2}{2\mu\Delta x^2}\cdot\frac{1}{2\pi}\int_{-\pi}^{\pi}u^2\,e^{ipu}\,du\,.
\end{equation}

\textbf{Diagonal ($p = 0$):}
\begin{equation}
  \frac{1}{2\pi}\int_{-\pi}^{\pi}u^2\,du = \frac{1}{2\pi}\cdot\frac{2\pi^3}{3} = \frac{\pi^2}{3}\,.
\end{equation}
Therefore $T_{jj} = \frac{\hb^2\pi^2}{6\mu\Delta x^2}$. $\checkmark$

\textbf{Off-diagonal ($p \neq 0$):}

Integrate by parts twice. Let $I_p = \frac{1}{2\pi}\int_{-\pi}^{\pi}u^2 e^{ipu}du$.

First integration by parts ($v = u^2$, $dw = e^{ipu}du$, $w = e^{ipu}/(ip)$):
\begin{align}
  I_p &= \frac{1}{2\pi}\left[\frac{u^2 e^{ipu}}{ip}\right]_{-\pi}^{\pi}
       - \frac{1}{2\pi}\int_{-\pi}^{\pi}\frac{2u\,e^{ipu}}{ip}\,du\notag\\
  &= \frac{1}{2\pi}\cdot\frac{\pi^2}{ip}\bigl(e^{ip\pi}-e^{-ip\pi}\bigr)
     - \frac{1}{i\pi p}\int_{-\pi}^{\pi}u\,e^{ipu}\,du\notag\\
  &= \frac{\pi^2}{2\pi ip}\cdot 2i\sin(p\pi)
     - \frac{1}{i\pi p}\int_{-\pi}^{\pi}u\,e^{ipu}\,du\,.
\end{align}
Since $p$ is an integer, $\sin(p\pi) = 0$, so the first term vanishes.

Second integration by parts on $\int_{-\pi}^{\pi}u\,e^{ipu}\,du$:
\begin{align}
  \int_{-\pi}^{\pi}u\,e^{ipu}\,du
  &= \left[\frac{u\,e^{ipu}}{ip}\right]_{-\pi}^{\pi} - \int_{-\pi}^{\pi}\frac{e^{ipu}}{ip}\,du\notag\\
  &= \frac{\pi e^{ip\pi}+\pi e^{-ip\pi}}{ip} - \frac{1}{ip}\underbrace{\int_{-\pi}^{\pi}e^{ipu}\,du}_{=\,2\pi\delta_{p,0}\,=\,0}\notag\\
  &= \frac{2\pi\cos(p\pi)}{ip} = \frac{2\pi(-1)^p}{ip}\,.
\end{align}
Substituting back:
\begin{equation}
  I_p = -\frac{1}{i\pi p}\cdot\frac{2\pi(-1)^p}{ip} = \frac{2(-1)^p}{p^2}\,.
\end{equation}
Therefore:
\begin{equation}
  T_{jl} = \frac{\hb^2}{2\mu\Delta x^2}\cdot\frac{2(-1)^{j-l}}{(j-l)^2}
  = \frac{\hb^2(-1)^{j-l}}{\mu\Delta x^2(j-l)^2}\,.\quad\checkmark
\end{equation}

\textbf{Connection to sinc functions.} The sinc-DVR basis functions are
$\chi_j(x) = \text{sinc}((x-x_j)/\Delta x)$. These satisfy the DVR property:
$\chi_j(x_l) = \delta_{jl}$. They form an orthonormal set on the uniform grid,
and the kinetic energy matrix in this basis yields exactly \cref{eq:T_DVR_diag,eq:T_DVR_offdiag}.

% ------------------------------------------------------------
\subsection{Potential energy matrix in DVR}\label{sec:DVR_V}
% ------------------------------------------------------------

The fundamental advantage of the DVR is that the potential energy matrix is \emph{diagonal}:
\begin{equation}\label{eq:V_DVR}
  \boxed{V_{jl} = V(x_j)\,\delta_{jl}\,.}
\end{equation}
This follows from the DVR property $\chi_j(x_l) = \delta_{jl}$:
\begin{equation}
  V_{jl} = \int\chi_j(x)\,V(x)\,\chi_l(x)\,dx \approx V(x_j)\int\chi_j(x)\chi_l(x)\,dx = V(x_j)\delta_{jl}\,,
\end{equation}
where the approximation becomes exact in the limit of a complete DVR basis.
No integrals need to be computed --- $V_{jl}$ is simply the potential evaluated at the grid points.

The full Hamiltonian matrix is therefore:
\begin{equation}
  H_{jl} = T_{jl} + V(x_j)\,\delta_{jl}\,.
\end{equation}

% ------------------------------------------------------------
\subsection{Convergence considerations}\label{sec:DVR_convergence}
% ------------------------------------------------------------

\subsubsection{Grid spacing}

The DVR can resolve momenta up to $k_{\max} = \pi/\Delta x$ (Nyquist limit).
To represent eigenstates with energy up to $E_{\max}$, the maximum kinetic energy
is approximately $E_{\max} + |V_{\min}|$ (where $V_{\min}$ is the potential minimum),
so $k_{\max} \sim \sqrt{2\mu(E_{\max}+|V_{\min}|)}$ and:
\begin{equation}
  \Delta x \lesssim \frac{\pi}{\sqrt{2\mu(E_{\max}+|V_{\min}|)}}\,.
\end{equation}

\subsubsection{Grid range}

The grid must be large enough that the wavefunctions are negligible at the boundaries:
$|\Psi(\pm L/2)|^2 \lesssim \eps_{\text{boundary}}$. For a double-well potential, a
typical choice is $L \approx 4$--$6\,\xe$.

\subsubsection{Total number of grid points}

Combining both requirements: $N = L/\Delta x$. For the NH$_3$ double well,
$N \sim 100$--$500$ typically provides convergence of the lowest $\sim 10$ levels
to machine precision.

% ------------------------------------------------------------
\subsection{DVR as a reference method}\label{sec:DVR_benchmark}
% ------------------------------------------------------------

The DVR/FGH method has several properties that make it ideal as a benchmark:
\begin{enumerate}
  \item \textbf{No basis set bias:} Unlike the HO or Gaussian bases, the DVR does not
        assume any particular functional form for the eigenstates.
  \item \textbf{Systematically improvable:} Accuracy improves monotonically with $N$, controlled
        by a single parameter ($\Delta x$ or $N$).
  \item \textbf{Exact potential representation:} The potential matrix is diagonal and exact
        (no quadrature errors for polynomial potentials).
  \item \textbf{No overlap issues:} The basis is orthonormal by construction.
  \item \textbf{Dense matrices:} The kinetic energy matrix is dense ($T_{jl} \neq 0$ for
        all $j, l$), so the computational cost scales as $\mathcal{O}(N^3)$ for diagonalization.
        This is the main disadvantage compared to sparse variational methods.
\end{enumerate}

In practice, the DVR results with sufficiently large $N$ serve as ``numerically exact''
reference values against which the variational methods (Bases~I--III) are benchmarked.

% ============================================================
% ============================================================

\section{Basis Set Comparison}\label{sec:comparison}
% ============================================================
% ============================================================

\begin{center}
\renewcommand{\arraystretch}{1.5}
\begin{tabular}{lcccc}
  \toprule
  Property & Basis I (HO) & Basis II (Shifted HO) & Basis III (Gaussian) & Basis IV (DVR) \\
  \midrule
  Orthogonal          & Yes & No  & No  & Yes \\
  Overlap $\mathbf{S}$ & $\mathbf{I}$ & Non-trivial & Non-trivial & $\mathbf{I}$ \\
  $\mathbf{H}$ structure & Banded & Dense & Dense & Dense \\
  Parity factorization & Yes (symmetric $V$) & Possible & No & No \\
  $\mathbf{V}$ matrix  & Analytic & Analytic & Analytic & Diagonal \\
  Parameters           & $\alpha$ & $\alpha, \xe$ & $\{x_i, \alpha_i\}$ & $\Delta x, N$ \\
  Typical $N$ for $10^{-8}$ & 30--50 & 15--25 & 20--40 & 200--500 \\
  Benchmark quality    & No & No & No & Yes \\
  \bottomrule
\end{tabular}
\end{center}

% ------------------------------------------------------------
\subsection{Analytical comparison for the asymmetric double well}\label{sec:basis_asym}
% ------------------------------------------------------------

The summary above compares the four bases in generic terms.  For the \emph{asymmetric}
potential $V(x) = ax^4 - bx^2 + cx$ ($c \neq 0$), several properties change qualitatively,
making the DVR (Basis~IV) analytically preferable to the shifted-HO (Basis~II) despite
Basis~II achieving faster convergence per function.  This subsection examines that
trade-off in detail.

\subsubsection{Matrix element complexity}

\textbf{DVR/PIB basis.}
The kinetic energy matrix is given by the closed-form Colbert--Miller expressions
(\cref{eq:T_DVR_diag,eq:T_DVR_offdiag}).  The potential matrix is diagonal:
$V_{jl} = V(x_j)\,\delta_{jl}$ (\cref{eq:V_DVR}).  No integral is ever computed;
the Hamiltonian is assembled solely from point evaluations of $V(x)$ and the analytic
kinetic-energy formula.  Changing the potential requires changing only the function
evaluation, not any integral routine.

\textbf{Shifted-HO basis.}
Same-center matrix elements reduce to origin-centered HO results via the binomial shift
(\cref{sec:shifted_potential}), but the cross-center blocks $\langle\varphi_m^{(L)}|x^k|
\varphi_n^{(R)}\rangle$ require: (i)~computing the exponential overlap factor
$S_{mn}^{LR}$ involving associated Laguerre polynomials (\cref{eq:S_LR}); (ii)~expanding
$x^k$ in displacements about the right center; and (iii)~contracting the result with
$S_{mn}^{LR}$.  The bookkeeping grows with $k$ and becomes substantially heavier when
$\alpha_L \neq \alpha_R$ (which is the generic asymmetric case, see below).

\subsubsection{Orthogonality and the eigenvalue problem}

The DVR basis is orthonormal, giving the standard eigenvalue problem $\mathbf{H}\mathbf{c}
= E\mathbf{c}$.

The shifted-HO basis is non-orthogonal, requiring the generalized problem
$\mathbf{H}\mathbf{c} = E\,\mathbf{S}\mathbf{c}$.  The overlap matrix $\mathbf{S}$ must
be computed, stored, and factored (e.g.\ via Cholesky or L\"owdin orthogonalization
$\mathbf{S}^{-1/2}$) before the standard solver can be applied.  This introduces an
additional source of numerical ill-conditioning: if the two well centers are close (shallow
barrier, strong asymmetry pushing $x_-$ and $x_+$ together), the cross-center overlaps
$S_{mn}^{LR}$ approach unity and $\mathbf{S}$ develops near-linear dependencies.  A
regularization or basis truncation is then required.

\begin{remark}
  The variational theorem (Hylleraas--Undheim--MacDonald) holds for the generalized
  eigenvalue problem, so monotone convergence with basis size is guaranteed in principle.
  In practice, however, near-linear dependence limits the useful basis size before
  numerical noise dominates.
\end{remark}

\subsubsection{Asymmetry-specific complications for Basis~II}

In the \emph{symmetric} case the shifted-HO basis benefits from:
\begin{itemize}
  \item $x_L = -\xe$, $x_R = +\xe$ known analytically;
  \item $\alpha_L = \alpha_R = \alpha_{\text{local}}$ by symmetry;
  \item a parity decomposition that halves the dimension of each block.
\end{itemize}

For $c \neq 0$, all three advantages are lost:
\begin{enumerate}
  \item \textbf{Shifted centers.}  The minima $x_\pm$ are roots of the depressed cubic
    $4ax^3 - 2bx + c = 0$ (\cref{eq:cubic_roots}), which must be solved (numerically
    or via the trigonometric formula) at each value of $c$.
  \item \textbf{Unequal widths.}  The curvatures $V''(x_\pm) = 12ax_\pm^2 - 2b$ differ,
    so $\omega_L \neq \omega_R$ and the natural width parameters $\alpha_L = \mu\omega_L$,
    $\alpha_R = \mu\omega_R$ must be computed separately.  With $\alpha_L \neq \alpha_R$,
    the displacement-operator formula \cref{eq:S_LR} no longer applies directly; the
    overlap integrals must be evaluated by combining two Gaussians of different widths,
    leading to a more involved computation.
  \item \textbf{No parity block structure.}  The linear term $cx$ couples even and odd
    functions (\cref{sec:parity}), so the full $N\times N$ generalized eigenvalue problem
    must be solved without any block decomposition.
  \item \textbf{Basis partitioning.}  The optimal split $N_L$ vs.\ $N_R$ between the two
    centers is not obvious a priori when the wells are unequal, and an inappropriate
    partition may under-represent one well.
  \item \textbf{Discontinuous basis under parameter variation.}  A sweep over $c$ requires
    recomputing centers, widths, and overlaps at every parameter value, making the basis
    non-uniform across the sweep.
\end{enumerate}

The DVR basis is entirely unaffected by all five points: the same grid with the same
parameters applies for any $c$.

\subsubsection{The DVR/FGH connection: elimination of all integrals}

The deepest analytical advantage of the DVR basis is the
\emph{integral-free representation} of any potential.  The Fourier Grid Hamiltonian
(\cref{sec:FGH}) exploits the fact that the PIB/sinusoidal functions define a unitary
transformation between the finite basis representation (FBR) and a uniform grid of
quadrature points $\{x_j\}$.  In the DVR, the potential matrix becomes exactly diagonal
(\cref{eq:V_DVR}):
\begin{equation}
  V_{jl} = V(x_j)\,\delta_{jl}\,.
\end{equation}
No integrals need to be computed for any potential, polynomial or otherwise.  The
Hamiltonian matrix is assembled in $\mathcal{O}(N^2)$ operations (one potential
evaluation per grid point, plus the analytic kinetic-energy formula).

The shifted-HO basis has no natural DVR.  A Gauss--Hermite quadrature exists for
each individual center, but combining two centers with different widths requires a
non-standard bi-centric quadrature that negates the simplicity of the approach.

\subsubsection{Time-dependent dynamics}

For wavepacket propagation studies (e.g.\ survival probability, expectation values as
functions of time), the DVR/PIB basis connects naturally to the \emph{split-operator}
method via the Fast Fourier Transform:
\begin{equation}
  e^{-i\hat{H}\Delta t/\hb}
  \approx e^{-i\hat{V}\Delta t/(2\hb)}\,\mathcal{F}^{-1}\,
          e^{-i\hat{T}\Delta t/\hb}\,\mathcal{F}\,
          e^{-i\hat{V}\Delta t/(2\hb)}
  + \mathcal{O}(\Delta t^3)\,,
\end{equation}
where $\mathcal{F}$ is the discrete Fourier transform.  Each step costs
$\mathcal{O}(N\log N)$.  The potential exponential is diagonal in the position
representation; the kinetic exponential is diagonal in momentum space.  This requires
no matrix diagonalization whatsoever.

For the shifted-HO basis, the propagator $e^{-i\mathbf{S}^{-1}\mathbf{H}\Delta t/\hb}$
must be computed as a dense matrix exponential, costing $\mathcal{O}(N^3)$ per
time step.

\subsubsection{Summary}

\begin{center}
\renewcommand{\arraystretch}{1.4}
\begin{tabular}{lcc}
  \toprule
  Criterion & DVR (Basis~IV) & Shifted-HO (Basis~II) \\
  \midrule
  Convergence rate (eigenstates needed) & Slower & \textbf{Faster} \\
  Matrix element simplicity & \textbf{Integral-free} & Heavy (cross-center) \\
  Eigenvalue problem type & \textbf{Standard} & Generalized \\
  Overlap conditioning & \textbf{None} & Risk near $\Delta E \to 0$ \\
  Parity exploitation ($c=0$) & No & Yes \\
  Robustness for $c \neq 0$ & \textbf{Full} & Loses parity; re-parameterize \\
  Parameter sweep ($c$ varies) & \textbf{Trivial} & Recompute basis each step \\
  Time-dependent propagation & \textbf{FFT split-op} $\mathcal{O}(N\log N)$ & Dense expm $\mathcal{O}(N^3)$ \\
  Physical transparency & Less & \textbf{More} \\
  \bottomrule
\end{tabular}
\end{center}

\textbf{Conclusion.}
The shifted-HO basis (Basis~II) is physically transparent and achieves excellent
convergence per function in the symmetric case, where the well centers and widths are
known analytically and parity halves the problem size.  In the asymmetric case these
advantages largely disappear: the basis must be redesigned at every parameter value,
the generalized eigenvalue problem introduces conditioning risk, and the integral-free
simplicity of the DVR is foregone.

The DVR/FGH basis (Basis~IV) requires modestly more grid points to reach the same
accuracy, but this cost is offset by: (i)~zero integral computation for any potential;
(ii)~a standard, well-conditioned eigenvalue problem; (iii)~complete insensitivity to
$c$; and (iv)~the ability to perform time-dependent propagation at $\mathcal{O}(N\log N)$
cost.  For the asymmetric double well---and particularly for parameter sweeps over the
asymmetry coefficient $c$---the DVR is the analytically and computationally preferred basis.

% ============================================================
% ============================================================

\section{Wavepacket Construction and Dynamics}\label{sec:wavepackets}
% ============================================================
% ============================================================

Given the eigenstates $\{\Phi_k, E_k\}_{k=0}^{N-1}$ obtained by diagonalization,
we construct localized wavepackets and study their time evolution.

% ------------------------------------------------------------
\subsection{Two-state wavepackets}\label{sec:2state}
% ------------------------------------------------------------

The simplest wavepackets use only the ground pair $(\Phi_0, \Phi_1)$:
\begin{equation}\label{eq:Psi_pm}
  \Psi_\pm(x) = \frac{1}{\sqrt{2}}\bigl(\Phi_0(x) \pm \Phi_1(x)\bigr)\,.
\end{equation}

\subsubsection{Localization in the symmetric case}

For the symmetric double well, $\Phi_0$ has even parity and $\Phi_1$ has odd parity.
Both are symmetric about each well (i.e., $|\Phi_0(x)|$ and $|\Phi_1(x)|$ have peaks
near $\pm\xe$), but $\Phi_1$ changes sign at $x = 0$.

$\Psi_+$: Near $x = +\xe$, both $\Phi_0$ and $\Phi_1$ are positive, so they add
constructively. Near $x = -\xe$, $\Phi_0 > 0$ but $\Phi_1 < 0$ (odd parity), so
they cancel. Therefore, $\Psi_+$ is localized in the \emph{right well}.

$\Psi_-$: By the same argument, $\Psi_-$ is localized in the \emph{left well}.

\textbf{Verification:}
\begin{align}
  \langle\Psi_\pm|\hat{x}|\Psi_\pm\rangle
  &= \frac{1}{2}\bigl(\mel{\Phi_0}{\hat{x}}{\Phi_0} + \mel{\Phi_1}{\hat{x}}{\Phi_1}
     \pm 2\mel{\Phi_0}{\hat{x}}{\Phi_1}\bigr)\notag\\
  &= \pm\mel{\Phi_0}{\hat{x}}{\Phi_1}\,,
\end{align}
where we used $\mel{\Phi_k}{\hat{x}}{\Phi_k} = 0$ for even potentials (parity).
Since $\mel{\Phi_0}{\hat{x}}{\Phi_1} > 0$ (by convention), $\langle x\rangle_+ > 0$
(right well) and $\langle x\rangle_- < 0$ (left well). $\checkmark$

\subsubsection{Time evolution}

\begin{align}
  \Psi_+(x, t) &= \frac{1}{\sqrt{2}}\bigl(e^{-iE_0 t/\hb}\Phi_0(x) + e^{-iE_1 t/\hb}\Phi_1(x)\bigr)\notag\\
  &= \frac{e^{-i\bar{E}t/\hb}}{\sqrt{2}}\bigl(e^{i\Delta t/(2\hb)}\Phi_0(x) + e^{-i\Delta t/(2\hb)}\Phi_1(x)\bigr)\,,
\end{align}
where $\bar{E} = (E_0+E_1)/2$ and $\Delta = E_1 - E_0$ is the tunnel splitting.

The probability density evolves as:
\begin{align}
  |\Psi_+(x,t)|^2 &= \frac{1}{2}\bigl[\Phi_0^2(x) + \Phi_1^2(x) + 2\Phi_0(x)\Phi_1(x)\cos(\Delta t/\hb)\bigr]\,.
\end{align}
At $t = 0$: $|\Psi_+|^2 = \frac{1}{2}(\Phi_0+\Phi_1)^2$ (localized right).
At $t = \pi\hb/\Delta$: $\cos(\Delta t/\hb) = \cos(\pi) = -1$, so $|\Psi_+|^2 = \frac{1}{2}(\Phi_0-\Phi_1)^2$ (localized left).

The \emph{tunneling period} (time for complete transfer from one well to the other) is:
\begin{equation}\label{eq:T_tunnel}
  \boxed{\ttun = \frac{\pi\hb}{\Delta} = \frac{\pi\hb}{E_1 - E_0} = \frac{h}{2(E_1-E_0)}\,.}
\end{equation}

The \emph{recurrence time} (time for the wavepacket to return to its initial state) is:
\begin{equation}\label{eq:t_rec}
  \boxed{\trec = \frac{2\pi\hb}{\Delta} = \frac{h}{E_1 - E_0} = 2\,\ttun\,.}
\end{equation}

% ------------------------------------------------------------
\subsection{Four-state wavepackets with mixing angle $\vartheta$}\label{sec:4state}
% ------------------------------------------------------------

A richer family of wavepackets incorporates the first excited pair $(\Phi_2, \Phi_3)$:
\begin{equation}\label{eq:Psi_theta}
  \Psi(\vartheta, x) = \frac{\cos\vartheta}{\sqrt{2}}\bigl(\Phi_0(x) + \Phi_1(x)\bigr)
  + \frac{\sin\vartheta}{\sqrt{2}}\bigl(\Phi_2(x) + \Phi_3(x)\bigr)\,,
\end{equation}
with $\vartheta \in [0, \pi/2]$.

\subsubsection{Expansion coefficients}

The coefficients in the eigenstate expansion $\Psi(\vartheta) = \sum_k c_k\Phi_k$ are:
\begin{equation}\label{eq:4state_ck}
  c_0 = c_1 = \frac{\cos\vartheta}{\sqrt{2}}\,,\qquad
  c_2 = c_3 = \frac{\sin\vartheta}{\sqrt{2}}\,.
\end{equation}
Normalization: $\sum_k |c_k|^2 = \cos^2\vartheta + \sin^2\vartheta = 1$. $\checkmark$

\subsubsection{Mean energy}

\begin{align}
  \langle E\rangle(\vartheta)
  &= \sum_k |c_k|^2 E_k
  = \frac{\cos^2\vartheta}{2}(E_0+E_1) + \frac{\sin^2\vartheta}{2}(E_2+E_3)\notag\\
  &= \cos^2\vartheta\,\bar{E}_{01} + \sin^2\vartheta\,\bar{E}_{23}\,,
\end{align}
where $\bar{E}_{01} = (E_0+E_1)/2$ and $\bar{E}_{23} = (E_2+E_3)/2$ are the mean
energies of the ground and excited doublets, respectively.

\subsubsection{Time evolution}

\begin{equation}
  \Psi(\vartheta, x, t) = \sum_{k=0}^{3}c_k\,e^{-iE_k t/\hb}\,\Phi_k(x)\,.
\end{equation}

\subsubsection{Survival probability}

Using \cref{eq:4state_ck} and the general formula $P(t) = |\sum_k |c_k|^2 e^{-iE_k t/\hb}|^2$:
\begin{align}
  A(t) &= \sum_k |c_k|^2 e^{-iE_k t/\hb}
  = \frac{\cos^2\vartheta}{2}\bigl(e^{-iE_0 t/\hb}+e^{-iE_1 t/\hb}\bigr)
  + \frac{\sin^2\vartheta}{2}\bigl(e^{-iE_2 t/\hb}+e^{-iE_3 t/\hb}\bigr)\notag\\
  &= \cos^2\vartheta\,e^{-i\bar{E}_{01}t/\hb}\cos\frac{\Delta_{01}t}{2\hb}
  + \sin^2\vartheta\,e^{-i\bar{E}_{23}t/\hb}\cos\frac{\Delta_{23}t}{2\hb}\,,
\end{align}
where $\Delta_{01} = E_1 - E_0$ and $\Delta_{23} = E_3 - E_2$.

The survival probability $P(t) = |A(t)|^2$ expands to:
\begin{align}\label{eq:P_4state}
  P(t) &= \cos^4\vartheta\cos^2\frac{\Delta_{01}t}{2\hb}
  + \sin^4\vartheta\cos^2\frac{\Delta_{23}t}{2\hb}\notag\\
  &\quad + 2\cos^2\vartheta\sin^2\vartheta\cos\frac{\Delta_{01}t}{2\hb}\cos\frac{\Delta_{23}t}{2\hb}
  \cos\frac{(\bar{E}_{23}-\bar{E}_{01})t}{\hb}\,.
\end{align}

\textbf{Derivation of \cref{eq:P_4state}:}

Write $A(t) = A_1 + A_2$ with $A_1 = \cos^2\vartheta\,e^{-i\bar{E}_{01}t/\hb}\cos(\Delta_{01}t/(2\hb))$
and $A_2 = \sin^2\vartheta\,e^{-i\bar{E}_{23}t/\hb}\cos(\Delta_{23}t/(2\hb))$.
\begin{align}
  P(t) &= |A_1|^2 + |A_2|^2 + 2\text{Re}(A_1^*A_2)\notag\\
  &= \cos^4\vartheta\cos^2\frac{\Delta_{01}t}{2\hb}
  + \sin^4\vartheta\cos^2\frac{\Delta_{23}t}{2\hb}\notag\\
  &\quad + 2\cos^2\vartheta\sin^2\vartheta\cos\frac{\Delta_{01}t}{2\hb}\cos\frac{\Delta_{23}t}{2\hb}
  \text{Re}\bigl(e^{i(\bar{E}_{01}-\bar{E}_{23})t/\hb}\bigr)\,.
\end{align}
Since $\text{Re}(e^{i\phi}) = \cos\phi$, this gives \cref{eq:P_4state}. $\square$

Alternatively, using the component form with all six frequency pairs:
\begin{align}
  P(t) &= \sum_k |c_k|^4 + 2\sum_{k<l}|c_k|^2|c_l|^2\cos\frac{(E_l-E_k)t}{\hb}\notag\\
  &= \frac{\cos^4\vartheta + \sin^4\vartheta}{2}
  + \frac{\cos^4\vartheta}{2}\cos\frac{\Delta_{01}t}{\hb}
  + \frac{\sin^4\vartheta}{2}\cos\frac{\Delta_{23}t}{\hb}\notag\\
  &\quad + \frac{\cos^2\vartheta\sin^2\vartheta}{2}\left[
    \cos\frac{(E_2-E_0)t}{\hb} + \cos\frac{(E_2-E_1)t}{\hb}\right.\notag\\
  &\qquad\qquad\qquad\qquad\qquad\left.
    + \cos\frac{(E_3-E_0)t}{\hb} + \cos\frac{(E_3-E_1)t}{\hb}
  \right]\,.
\end{align}
These two forms are equivalent via product-to-sum identities.

\textbf{Limiting cases:}
\begin{itemize}
  \item $\vartheta = 0$: $P(t) = \cos^2(\Delta_{01}t/(2\hb))$, the two-state result.
  \item $\vartheta = \pi/2$: $P(t) = \cos^2(\Delta_{23}t/(2\hb))$, tunneling in the excited doublet.
  \item $\vartheta = \pi/4$: Equal weight on both doublets, $P(t)$ shows beating between
        $\Delta_{01}$ and $\Delta_{23}$ modulated by $\bar{E}_{23}-\bar{E}_{01}$.
\end{itemize}

% ------------------------------------------------------------
\subsection{Localized states in the asymmetric case}\label{sec:localized_asym}
% ------------------------------------------------------------

When $\varepsilon \neq 0$ (see \cref{sec:V_asym}), the two lowest eigenstates
$\Phi_0, \Phi_1$ of $\hat{H}^{\text{asym}}$ are no longer related by parity.
Localized wavepackets are still constructed as superpositions of these two states,
\begin{equation}\label{eq:psi_LR_asym}
  \Psi_L = \cos\theta\,\Phi_0 + \sin\theta\,\Phi_1\,,
  \qquad
  \Psi_R = -\sin\theta\,\Phi_0 + \cos\theta\,\Phi_1\,,
\end{equation}
but the mixing angle $\theta$ is now determined by the effective two-level Hamiltonian.

\subsubsection{Two-level effective Hamiltonian}

In the localized-state basis $\{\Psi_L^{(0)}, \Psi_R^{(0)}\}$ built from the
\emph{unperturbed} ($\varepsilon=0$) symmetric eigenstates via \cref{eq:Psi_pm},
the projection of $\hat{H}^{\text{asym}}$ onto the two lowest levels takes the form:
\begin{equation}
  H_{\text{eff}} = \bar{E}\,\mathbf{I}
  + \begin{pmatrix}
    +\varepsilon_{\text{eff}} & -\Delt/2 \\
    -\Delt/2 & -\varepsilon_{\text{eff}}
  \end{pmatrix},
\end{equation}
where $\bar{E} = (E_0^{(0)}+E_1^{(0)})/2$ is the mean unperturbed energy,
$\Delt = E_1^{(0)}-E_0^{(0)}$ is the symmetric tunnel splitting, and the
effective bias is
\begin{equation}
  \varepsilon_{\text{eff}} = \varepsilon\,\mel{\Phi_0}{\hat{x}}{\Phi_1}\,.
\end{equation}
The diagonal elements of $\hat{x}$ vanish by parity for the unperturbed eigenstates:
$\mel{\Phi_i}{\hat{x}}{\Phi_i} = 0$.

\subsubsection{Diagonalization and mixing angle}

The eigenvalues of $H_{\text{eff}}$ (relative to $\bar{E}$) are
\begin{equation}
  E_\pm = \pm\Omg\,,\qquad
  \Omg = \sqrt{\varepsilon_{\text{eff}}^2 + (\Delt/2)^2}\,.
\end{equation}
The total asymmetric energy gap is
\begin{equation}\label{eq:gap_asym}
  \Delta E = 2\Omg = \sqrt{\Delt^2 + 4\varepsilon_{\text{eff}}^2}\,.
\end{equation}
The mixing angle $\theta$ satisfies
\begin{equation}\label{eq:mixing_angle}
  \boxed{\tan(2\theta) = \frac{\Delt}{2\varepsilon_{\text{eff}}}\,,}
  \qquad
  \cos(2\theta) = \frac{\varepsilon_{\text{eff}}}{\Omg}\,,
  \qquad
  \sin(2\theta) = \frac{\Delt}{2\Omg}\,.
\end{equation}

\subsubsection{Left-well survival probability}

Starting from $\Psi_L$ (left-well localized state), the probability of finding
the particle in the left well at time $t$ is
\begin{equation}\label{eq:P_L_asym}
  \boxed{P_L(t) = 1 - \frac{\Delt^2}{\Delta E^2}\sin^2\!\left(\frac{\Omg\,t}{\hb}\right)\,.}
\end{equation}
The maximum transfer probability is $(\Delt/\Delta E)^2 \leq 1$, which is
strictly less than 1 when $\varepsilon_{\text{eff}} \neq 0$. The tunneling
efficiency is thus $\eta = (\Delt/\Delta E)^2$.

\subsubsection{Limiting cases}

\begin{description}
  \item[Symmetric limit $\varepsilon \to 0$:] $\varepsilon_{\text{eff}} \to 0$,
    $\tan(2\theta)\to\infty$, $\theta\to\pi/4$.  The localized states reduce to
    $(\Phi_0\pm\Phi_1)/\sqrt{2}$ and complete transfer ($P_L$ oscillates 0 to 1) is restored.
  \item[Strong-asymmetry limit $|\varepsilon|\to\infty$:] $\varepsilon_{\text{eff}}\gg\Delt/2$,
    $\theta\to 0$. The localized states tend to the energy eigenstates
    ($\Psi_L\to\Phi_0$, $\Psi_R\to\Phi_1$) and tunneling is suppressed:
    $P_L(t)\approx 1-(\Delt/(2\varepsilon_{\text{eff}}))^2\sin^2(\Omg t/\hb)\approx 1$.
  \item[Resonance $\varepsilon_{\text{eff}} = \Delt/2$:] $\theta=\pi/8$,
    $\Delta E=\Delt\sqrt{2}$. Maximum transfer probability is $1/2$.
\end{description}

\subsubsection{Well selection: sign convention and basis independence}

The sign structure of \cref{eq:psi_LR_asym} is not arbitrary---it directly selects
the well. To see why, note that the mixing angle $\theta$ and the labelling
$\Psi_L / \Psi_R$ arise from \emph{inverting} the diagonalization of $H_{\text{eff}}$.
The rotation that maps localized (diabatic) states $\{|L\rangle, |R\rangle\}$ to
energy eigenstates $\{\Phi_0, \Phi_1\}$ is

\begin{equation}
  \begin{pmatrix}\Phi_0 \\ \Phi_1\end{pmatrix}
  =
  \begin{pmatrix}\cos\theta & -\sin\theta \\ \sin\theta & \cos\theta\end{pmatrix}
  \begin{pmatrix}\Psi_L \\ \Psi_R\end{pmatrix}.
\end{equation}

Inverting this rotation recovers \cref{eq:psi_LR_asym}: the $(+\cos\theta, +\sin\theta)$
combination concentrates the probability density near the left minimum $x_-$, while the
$(-\sin\theta, +\cos\theta)$ combination concentrates it near $x_+$. Swapping the two
combinations selects the right well. The computational basis (HO, shifted-HO, DVR, etc.)
used to obtain $\Phi_0$ and $\Phi_1$ is irrelevant: once the eigenstates are known, the
rotation is uniquely determined by $\theta$.

\begin{remark}
  In practice one identifies which combination corresponds to which well by evaluating
  $\langle\hat{x}\rangle_{\Psi_L} = \cos^2\theta\,x_{00} + \sin^2\theta\,x_{11}
   + \sin(2\theta)\,x_{01}$ (using \cref{eq:xt_general}) or simply by inspecting
  whether the probability density is peaked at $x_-$ or $x_+$.
\end{remark}

\subsubsection{Mechanism of localization: symmetric vs.\ asymmetric case}

The two cases differ fundamentally in how localization is achieved.

\textbf{Symmetric case ($\varepsilon = 0$, $\theta = \pi/4$).}
The eigenstates $\Phi_0$ and $\Phi_1$ carry definite parity ($+$ and $-$ respectively).
Both have equal weight in both wells, but they are out of phase across the barrier.
Localization is produced by \emph{parity-based destructive interference}: near $x = +\xe$,
$\Phi_0$ and $\Phi_1$ add constructively in $\Psi_+$; near $x = -\xe$, the opposite
sign of $\Phi_1$ causes cancellation.  The result is an equal superposition
$\Psi_\pm = (\Phi_0 \pm \Phi_1)/\sqrt{2}$ that is perfectly localized (within the
two-state approximation).

\textbf{Asymmetric case ($\varepsilon \neq 0$, $\theta \neq \pi/4$).}
Neither eigenstate carries definite parity. Instead, $\Phi_0$ is already biased toward
the deeper well and $\Phi_1$ toward the shallower one.  Localization is now achieved by a
\emph{rotation in the two-dimensional eigenstate subspace}: the angle $\theta$ adjusts
the superposition coefficients so that one combination $\Psi_L$ has its density concentrated
near $x_-$ and the other $\Psi_R$ near $x_+$. The parity-cancellation mechanism plays no
role; all that matters is the correct rotation angle.

This distinction has a practical consequence: parity cannot be exploited to simplify the
eigenvalue problem when $\varepsilon \neq 0$, and the identification of the ``left'' vs.\
``right'' localized state requires computing $\theta$ rather than relying on symmetry arguments.

\subsubsection{Tunneling efficiency and the localization--tunneling complementarity}

The tunneling efficiency
\begin{equation}\label{eq:eta_tun}
  \boxed{\eta \equiv \left(\frac{\Delt}{\Delta E}\right)^2 = \sin^2(2\theta) \leq 1}
\end{equation}
measures the maximum fraction of the population that can transfer between wells.
There is a fundamental complementarity between the quality of initial localization and
the extent of tunneling: a strongly localized initial state (large $|\varepsilon_{\text{eff}}|$)
is nearly an eigenstate and therefore barely tunnels.

\begin{center}
\renewcommand{\arraystretch}{1.4}
\begin{tabular}{lcccl}
  \toprule
  Regime & $\varepsilon_{\text{eff}}/\Delt$ & $\theta$ & $\eta$ & Physical picture \\
  \midrule
  Symmetric & $0$ & $\pi/4$ & $1$ & Complete transfer \\
  Weak asymmetry & $\ll 1$ & $\approx\pi/4$ & $\approx 1$ & Near-complete transfer \\
  Resonance & $1/2$ & $\pi/8$ & $1/2$ & Half transfer \\
  Strong asymmetry & $\gg 1$ & $\approx 0$ & $\approx 0$ & Tunneling suppressed \\
  \bottomrule
\end{tabular}
\end{center}

In the strong-asymmetry limit, $P_L(t) \approx 1 - \eta\sin^2(\Omg t/\hb)$ with
$\eta \approx (\Delt/2\varepsilon_{\text{eff}})^2 \ll 1$: the oscillation frequency
$\Omg \approx |\varepsilon_{\text{eff}}|$ is large but the amplitude of probability
transfer is negligibly small. Physically, the bias energy mismatch between the two
wells detunes the tunneling resonance, just as in the spin-$\tfrac{1}{2}$ analogy
(Rabi problem with detuning $2\varepsilon_{\text{eff}}$ and coupling $\Delt$).

% ------------------------------------------------------------
\subsection{Expectation value $\langle\hat{x}\rangle(t)$}\label{sec:x_expect}
% ------------------------------------------------------------

\subsubsection{General formula}

For a wavepacket $\Psi(t) = \sum_k c_k e^{-iE_k t/\hb}\Phi_k$:
\begin{align}\label{eq:xt_general}
  \langle\hat{x}\rangle(t)
  &= \sum_{j,k}c_j^*c_k\,e^{i(E_j-E_k)t/\hb}\mel{\Phi_j}{\hat{x}}{\Phi_k}\notag\\
  &= \sum_k |c_k|^2\mel{\Phi_k}{\hat{x}}{\Phi_k}
  + 2\sum_{j<k}\text{Re}\bigl[c_j^*c_k\,e^{-i(E_k-E_j)t/\hb}\bigr]\mel{\Phi_j}{\hat{x}}{\Phi_k}\,.
\end{align}
For real coefficients $c_j$:
\begin{equation}\label{eq:xt_real}
  \langle\hat{x}\rangle(t) = \sum_k c_k^2\,x_{kk}
  + 2\sum_{j<k}c_j c_k\,x_{jk}\cos\frac{(E_k-E_j)t}{\hb}\,,
\end{equation}
where $x_{jk} \equiv \mel{\Phi_j}{\hat{x}}{\Phi_k}$.

\subsubsection{Eigenstate matrix elements of $\hat{x}$}

Each eigenstate is expanded in the basis: $\Phi_k(x) = \sum_n C_{n}^{(k)}\varphi_n(x)$.
The matrix element is:
\begin{align}
  x_{jk} &= \mel{\Phi_j}{\hat{x}}{\Phi_k}
  = \sum_{m,n}C_m^{(j)}C_n^{(k)}\mel{\varphi_m}{\hat{x}}{\varphi_n}\,.
\end{align}
For the HO basis:
\begin{equation}
  \mel{\varphi_m}{\hat{x}}{\varphi_n} = \frac{1}{\sqrt{2\alpha}}\bigl(\sqrt{n}\,\delta_{m,n-1}+\sqrt{n+1}\,\delta_{m,n+1}\bigr)\,,
\end{equation}
so:
\begin{equation}\label{eq:xjk_HO}
  \boxed{x_{jk} = \frac{1}{\sqrt{2\alpha}}\sum_n \sqrt{n+1}\Bigl[C_n^{(j)}C_{n+1}^{(k)} + C_{n+1}^{(j)}C_n^{(k)}\Bigr]\,.}
\end{equation}

\subsubsection{Selection rules for the symmetric case}

For the symmetric potential, $\Phi_k$ has definite parity $(-1)^k$.
Since $\hat{x}$ is an odd operator:
\begin{itemize}
  \item $x_{kk} = \mel{\Phi_k}{\hat{x}}{\Phi_k} = 0$ for all $k$ (diagonal elements vanish).
  \item $x_{jk} \neq 0$ only when $\Phi_j$ and $\Phi_k$ have opposite parity, i.e., $j+k$ is odd.
\end{itemize}

For the two-state wavepacket $\Psi_+ = (\Phi_0 + \Phi_1)/\sqrt{2}$:
\begin{equation}
  \langle\hat{x}\rangle(t) = x_{01}\cos\frac{\Delta_{01}t}{\hb}\,,
\end{equation}
which oscillates between $+|x_{01}|$ (right well) and $-|x_{01}|$ (left well) with
period $\trec = 2\pi\hb/\Delta_{01}$.

For the four-state wavepacket:
\begin{align}
  \langle\hat{x}\rangle(t)
  &= \cos^2\vartheta\,x_{01}\cos\frac{\Delta_{01}t}{\hb}
  + \sin^2\vartheta\,x_{23}\cos\frac{\Delta_{23}t}{\hb}\notag\\
  &\quad + \frac{\cos\vartheta\sin\vartheta}{\sqrt{2}}\,x_{03}\cos\frac{(E_3-E_0)t}{\hb}\notag\\
  &\quad + \frac{\cos\vartheta\sin\vartheta}{\sqrt{2}}\,x_{12}\cos\frac{(E_2-E_1)t}{\hb}\notag\\
  &\quad + \frac{\cos\vartheta\sin\vartheta}{\sqrt{2}}\,x_{01}^{(23)}\cos\frac{(E_2-E_0)t}{\hb}\notag\\
  &\quad + \frac{\cos\vartheta\sin\vartheta}{\sqrt{2}}\,x_{01}^{(31)}\cos\frac{(E_3-E_1)t}{\hb}\,,
\end{align}
where $x_{01}^{(23)}$ and $x_{01}^{(31)}$ represent cross-doublet matrix elements
$\mel{\Phi_0}{\hat{x}}{\Phi_2}$ and $\mel{\Phi_1}{\hat{x}}{\Phi_3}$.
Note: $\mel{\Phi_0}{\hat{x}}{\Phi_2} = 0$ and $\mel{\Phi_1}{\hat{x}}{\Phi_3} = 0$
by parity (both pairs have the same parity). Therefore, only opposite-parity pairs
$(0,1)$, $(0,3)$, $(1,2)$, $(2,3)$ contribute.

% ------------------------------------------------------------
\subsection{Turning points}\label{sec:turning_points}
% ------------------------------------------------------------

For a wavepacket with total energy $E_{\text{total}} = \langle E\rangle(\vartheta)$,
the classical turning points are the solutions of $V(x_{\text{tp}}) = E_{\text{total}}$:
\begin{equation}
  a\,x_{\text{tp}}^4 - b\,x_{\text{tp}}^2 = E_{\text{total}}\,.
\end{equation}
This is a quadratic equation in $u = x_{\text{tp}}^2$:
\begin{equation}
  a\,u^2 - b\,u - E_{\text{total}} = 0\,.
\end{equation}
By the quadratic formula:
\begin{equation}
  u = \frac{b \pm \sqrt{b^2 + 4a\,E_{\text{total}}}}{2a}\,.
\end{equation}

\textbf{Case 1: $E_{\text{total}} > 0$ (above the barrier).}
The discriminant $b^2 + 4aE > 0$. The ``$+$'' root gives $u_+ > 0$ (since $b > 0$
and $\sqrt{b^2+4aE} > 0$). The ``$-$'' root gives $u_- = (b - \sqrt{b^2+4aE})/(2a) < 0$
(since $\sqrt{b^2+4aE} > b$ when $E > 0$). A negative $u$ has no real $x$ solutions.
Therefore, there are only \emph{two} turning points:
\begin{equation}
  x_{\text{tp}} = \pm\sqrt{u_+} = \pm\sqrt{\frac{b+\sqrt{b^2+4aE_{\text{total}}}}{2a}}\,.
\end{equation}

\textbf{Case 2: $-\Vb < E_{\text{total}} < 0$ (below the barrier, within the wells).}
Now $|4aE| < b^2$ (since $\Vb = b^2/(4a)$), so $b^2 + 4aE > 0$ and both roots are positive:
$u_+ > u_- > 0$. There are \emph{four} turning points:
\begin{equation}\label{eq:4tp}
  \boxed{x_{\text{tp}} = \pm\sqrt{u_+}\,,\quad \pm\sqrt{u_-}\,,}
\end{equation}
where $u_\pm = \frac{b \pm \sqrt{b^2+4aE_{\text{total}}}}{2a}$.
The particle oscillates between $-\sqrt{u_+}$ and $-\sqrt{u_-}$ in the left well,
or between $+\sqrt{u_-}$ and $+\sqrt{u_+}$ in the right well. The classically forbidden
region is $|x| < \sqrt{u_-}$ (under the barrier).

\textbf{Case 3: $E_{\text{total}} = -\Vb$ (at the bottom of the wells).}
$u_+ = u_- = b/(2a) = \xe^2$, so $x_{\text{tp}} = \pm\xe$ (the particle sits at the minima).

% ------------------------------------------------------------
\subsection{Gaussian wavepacket as initial state}\label{sec:Gaussian_WP}
% ------------------------------------------------------------

A minimum-uncertainty Gaussian wavepacket centered at $x_0$ with initial momentum $p_0$ is
\begin{equation}\label{eq:Psi_Gauss}
  \Psi_G(x, 0) = \left(\frac{2\alpha_0}{\pi}\right)^{1/4}
  \exp\bigl[-\alpha_0(x-x_0)^2\bigr]\exp(ip_0 x/\hb)\,,
\end{equation}
with $\alpha_0 > 0$ determining the spatial width ($\Delta x = 1/\sqrt{4\alpha_0}$).

\subsubsection{Expansion in eigenstates}

The coefficients in the eigenstate expansion are
\begin{equation}
  c_k = \braket{\Phi_k}{\Psi_G} = \int_{-\infty}^{\infty}\Phi_k^*(x)\,\Psi_G(x,0)\,dx\,.
\end{equation}

\textbf{For the HO basis:} $\Phi_k(x) = \sum_n C_n^{(k)}\varphi_n(x;\alpha)$, so
\begin{equation}
  c_k = \sum_n C_n^{(k)}\braket{\varphi_n}{\Psi_G}\,.
\end{equation}
The overlap of a HO eigenstate with a Gaussian is:
\begin{align}
  \braket{\varphi_n}{\Psi_G}
  &= \mathcal{N}_n\left(\frac{2\alpha_0}{\pi}\right)^{1/4}
  \int_{-\infty}^{\infty}H_n(\sqrt{\alpha}\,x)\,e^{-\alpha x^2/2}\,e^{-\alpha_0(x-x_0)^2}\,e^{ip_0 x/\hb}\,dx\,.
\end{align}
For zero initial momentum ($p_0 = 0$), the Gaussians can be combined using the
product formula (\cref{sec:Gauss_product}):
\begin{equation}
  e^{-\alpha x^2/2}\,e^{-\alpha_0(x-x_0)^2}
  = e^{-Q}\,e^{-A(x-X)^2}\,,
\end{equation}
with $A = \alpha/2 + \alpha_0$, $X = \alpha_0 x_0/A$, $Q = \alpha_0(\alpha/2)(x_0^2)/A$.
The integral becomes:
\begin{equation}
  \braket{\varphi_n}{\Psi_G} = \mathcal{N}_n\left(\frac{2\alpha_0}{\pi}\right)^{1/4}e^{-Q}
  \int_{-\infty}^{\infty}H_n(\sqrt{\alpha}\,x)\,e^{-A(x-X)^2}\,dx\,.
\end{equation}
Substituting $x = u + X$:
\begin{equation}
  = \mathcal{N}_n\left(\frac{2\alpha_0}{\pi}\right)^{1/4}e^{-Q}
  \int_{-\infty}^{\infty}H_n(\sqrt{\alpha}(u+X))\,e^{-Au^2}\,du\,.
\end{equation}
Using the Hermite polynomial addition formula
$H_n(\sqrt{\alpha}(u+X)) = \sum_{k=0}^{n}\binom{n}{k}(2\sqrt{\alpha}\,X)^{n-k}H_k(\sqrt{\alpha}\,u)$
and the orthogonality $\int H_k(\sqrt{\alpha}\,u)e^{-Au^2}du$, this integral can be
evaluated analytically. The result involves nested sums, and for practical purposes
is best computed numerically.

\textbf{For the Gaussian basis:} If $\alpha_0$ and $x_0$ match one of the basis
functions $G_i$, then $c_k$ reduces to the $i$-th component of the $k$-th eigenvector
(up to normalization). Otherwise, $c_k = \sum_i C_i^{(k)}\braket{G_i}{\Psi_G}/\sqrt{S_{ii}}$
involves the overlap of two Gaussians, which is analytically trivial.

% ============================================================
% ============================================================

\section{Physical Observables}\label{sec:observables}
% ============================================================
% ============================================================

We collect here the complete working formulas for all physical observables computed by QuTu.

% ------------------------------------------------------------
\subsection{Energy levels and tunneling splitting}
% ------------------------------------------------------------

The energy eigenvalues $E_0 < E_1 < E_2 < \cdots$ are obtained by diagonalization.
The \emph{tunneling splitting} of the $k$-th doublet is
\begin{equation}
  \boxed{\Delta_k = E_{2k+1} - E_{2k}\,,\qquad k = 0, 1, 2, \ldots}
\end{equation}
The ground-state splitting is $\Delta_0 = E_1 - E_0$.

The inversion splitting in spectroscopic units (cm$^{-1}$) is
\begin{equation}
  \boxed{\Delta\nu = \frac{E_1 - E_0}{hc} = \frac{\Delta_0}{2\pi\hb c}\,,}
\end{equation}
where $c$ is the speed of light. In atomic units ($\hb = 1$, energies in $E_h$):
$\Delta\nu\,[\text{cm}^{-1}] = \Delta_0\,[\text{a.u.}] \times 219474.63\,\text{cm}^{-1}/E_h$.

% ------------------------------------------------------------
\subsection{Survival probability (general $N$-state)}\label{sec:Ps_general}
% ------------------------------------------------------------

For a normalized wavepacket $\Psi(0) = \sum_k c_k\Phi_k$ with $\sum_k|c_k|^2 = 1$:
\begin{equation}\label{eq:Ps_general}
  \boxed{P_s(t) = \left|\sum_k |c_k|^2\,e^{-iE_k t/\hb}\right|^2
  = \sum_k |c_k|^4 + 2\sum_{k<l}|c_k|^2|c_l|^2\cos\frac{(E_l-E_k)t}{\hb}\,.}
\end{equation}

Properties:
\begin{itemize}
  \item $P_s(0) = 1$.
  \item $0 \leq P_s(t) \leq 1$.
  \item Time-averaged value: $\overline{P_s} = \sum_k |c_k|^4 \leq 1$.
  \item For two states: $\overline{P_s} = 1/2$.
  \item For $N$ equally weighted states ($|c_k| = 1/\sqrt{N}$): $\overline{P_s} = 1/N$.
\end{itemize}

% ------------------------------------------------------------
\subsection{Position expectation value $\langle\hat{x}\rangle(t)$}
% ------------------------------------------------------------

\begin{equation}\label{eq:xt_boxed}
  \boxed{\langle\hat{x}\rangle(t) = \sum_k |c_k|^2 x_{kk}
  + 2\sum_{j<k}\text{Re}\bigl[c_j^*c_k\,e^{-i(E_k-E_j)t/\hb}\bigr]x_{jk}\,,}
\end{equation}
where $x_{jk} = \mel{\Phi_j}{\hat{x}}{\Phi_k}$.
For the symmetric potential, $x_{kk} = 0$ and only opposite-parity pairs contribute.

% ------------------------------------------------------------
\subsection{Momentum expectation value $\langle\hat{p}\rangle(t)$}
% ------------------------------------------------------------

By Ehrenfest's theorem, $\langle\hat{p}\rangle = \mu\,d\langle\hat{x}\rangle/dt$.
Differentiating \cref{eq:xt_boxed} (for real $c_j$):
\begin{equation}
  \boxed{\langle\hat{p}\rangle(t) = 2\mu\sum_{j<k}c_j c_k\,\frac{E_k-E_j}{\hb}\,x_{jk}\sin\frac{(E_k-E_j)t}{\hb}\,.}
\end{equation}

\textbf{Derivation.} Starting from
\begin{equation}
  \langle\hat{x}\rangle(t) = 2\sum_{j<k}c_jc_k\,x_{jk}\cos\frac{(E_k-E_j)t}{\hb}\,,
\end{equation}
differentiate:
\begin{align}
  \frac{d\langle\hat{x}\rangle}{dt}
  &= -2\sum_{j<k}c_jc_k\,x_{jk}\cdot\frac{E_k-E_j}{\hb}\sin\frac{(E_k-E_j)t}{\hb}\,.
\end{align}
Then $\langle\hat{p}\rangle = \mu\,d\langle\hat{x}\rangle/dt$:
\begin{equation}
  \langle\hat{p}\rangle(t) = -\frac{2\mu}{\hb}\sum_{j<k}c_jc_k\,(E_k-E_j)\,x_{jk}\sin\frac{(E_k-E_j)t}{\hb}\,.
\end{equation}
Note: the sign depends on convention. Here $\langle\hat{p}\rangle(0) = 0$ for a
real wavepacket at rest, which is consistent since $\sin(0) = 0$. $\checkmark$

% ------------------------------------------------------------
\subsection{Position variance $\Delta x^2(t)$}
% ------------------------------------------------------------

The variance is
\begin{equation}
  \Delta x^2(t) = \langle\hat{x}^2\rangle(t) - \langle\hat{x}\rangle^2(t)\,,
\end{equation}
where
\begin{equation}
  \langle\hat{x}^2\rangle(t) = \sum_{j,k}c_j^*c_k\,e^{i(E_j-E_k)t/\hb}\mel{\Phi_j}{\hat{x}^2}{\Phi_k}\,.
\end{equation}
The matrix elements $\mel{\Phi_j}{\hat{x}^2}{\Phi_k}$ are computed analogously to $x_{jk}$
using the HO matrix elements of $x^2$ (\cref{eq:x2_HO}).

% ------------------------------------------------------------
\subsection{Tunneling period and recurrence time}
% ------------------------------------------------------------

For the two-state wavepacket:
\begin{align}
  &\text{Tunneling period (one-way):}\quad \boxed{\ttun = \frac{\pi\hb}{E_1-E_0} = \frac{h}{2\Delta_0}\,,}\\[6pt]
  &\text{Recurrence time (round-trip):}\quad \boxed{\trec = \frac{2\pi\hb}{E_1-E_0} = \frac{h}{\Delta_0}\,.}
\end{align}

For the four-state wavepacket, exact recurrence requires $\Delta_{01}/\Delta_{23}$ to be
rational. The approximate recurrence time is $\trec \approx h/\Delta_{01}$ (dominated by the
smaller splitting).

% ------------------------------------------------------------
\subsection{Inversion splitting in cm$^{-1}$}
% ------------------------------------------------------------

\begin{equation}
  \boxed{\Delta\nu = \frac{E_1 - E_0}{hc}
  = (E_1 - E_0)\,[\text{a.u.}]\times 219474.63\;\text{cm}^{-1}/E_h\,.}
\end{equation}

For NH$_3$, the experimental value is $\Delta\nu \approx 0.79\,\text{cm}^{-1}$ for the
ground-state inversion splitting and $\Delta\nu \approx 35.8\,\text{cm}^{-1}$ for the
$\nu_2 = 1$ excited state.

% ============================================================
% ============================================================

\section{Numerical Implementation}\label{sec:numerical_impl}
% ============================================================
% ============================================================

\subsection{Variational method: algorithm}

The following steps summarize the complete numerical procedure implemented in QuTu:

\begin{enumerate}
  \item \textbf{Read physical parameters}: $\xe$, $\Vb$, reduced mass $\mu$
    (e.g., $\mu = 3m_H m_N/(3m_H+m_N)$ for NH$_3$), asymmetry parameter $\varepsilon$,
    and number of basis functions $N$.
  \item \textbf{Compute $\alpha_{\text{opt}}$} from \cref{sec:alpha_opt}:
    use the deep-well estimate $\alpha_{\text{opt}} = 2\sqrt{\mu b}$ as starting point,
    then refine numerically.
  \item \textbf{Symmetric case} ($\varepsilon = 0$): build $\mathbf{H}_{\text{even}}$
    (dimension $\lceil N/2\rceil$) and $\mathbf{H}_{\text{odd}}$ (dimension
    $\lfloor N/2\rfloor$) using \cref{eq:T_HO,eq:x2_HO,eq:x4_HO} restricted to
    even/odd quantum numbers; diagonalize both with a real symmetric eigensolver
    (e.g., LAPACK \texttt{DSYEV}); interleave eigenvalues as $E_0 < E_1 < E_2 < \cdots$
  \item \textbf{Asymmetric case} ($\varepsilon \neq 0$): build the full $N\times N$
    matrix $\mathbf{H}^{\text{asym}}$ from \cref{eq:H_asym} (adding
    $\varepsilon\mel{m}{\hat{x}}{n}$ from \cref{eq:x1_HO}); diagonalize with
    \texttt{DSYEV}.
  \item \textbf{Reconstruct wavefunctions}: from the eigenvectors
    $\{C_n^{(k)}\}$, evaluate $\Phi_k(x) = \sum_n C_n^{(k)}\varphi_n(x;\alpha)$
    on a real-space grid.
  \item \textbf{Construct localized wavepackets}: symmetric case via
    \cref{eq:Psi_pm}; asymmetric case via \cref{eq:psi_LR_asym,eq:mixing_angle}.
  \item \textbf{Compute observables}: survival probability from \cref{sec:Ps_general},
    $\langle\hat{x}\rangle(t)$ from \cref{sec:x_expect}, recurrence time
    $\trec = h/\Delta_0$, and additional observables as listed in
    \cref{sec:observables}.
\end{enumerate}

\subsection{Convergence criterion}

The energy $E_k^{(N)}$ computed with $N$ basis functions converges to the exact
value as $N\to\infty$. A practical convergence criterion is
\begin{equation}
  |E_k^{(N)} - E_k^{(N_{\max})}| < \delta_{\text{conv}}\,,
\end{equation}
where $N_{\max}$ is a large reference basis and $\delta_{\text{conv}}$ is a chosen
tolerance (e.g., $10^{-6}~E_h$). For the quartic double-well, the lowest levels
typically converge at $N \sim 20$--$50$.

\subsection{Unit conversion table}

The code works internally in atomic units. Key conversion factors:
\begin{center}
\renewcommand{\arraystretch}{1.3}
\begin{tabular}{lll}
  \toprule
  Quantity & Atomic unit & Practical unit \\
  \midrule
  Length   & $a_0 = 0.52918\;\text{\AA}$ & angstrom (\AA) \\
  Energy   & $E_h = 1\;\text{hartree}$   & $1\,E_h = 219474.63\;\text{cm}^{-1}$ \\
  Mass     & $m_e$                        & $1\;\text{u} = 1822.89\,m_e$ \\
  Time     & $\hb/E_h$                    & $1\;\hb/E_h = 24.19\;\text{as}$ \\
  \bottomrule
\end{tabular}
\end{center}

% ============================================================
% ============================================================

% ============================================================
% ============================================================

\section{Additional Observables}\label{sec:additional_obs}
% ============================================================
% ============================================================

This section develops a comprehensive suite of observables beyond the survival probability $P(t)$, the expectation value $\expval{\hat{x}}(t)$, and the recurrence time $\trec$ derived in the preceding sections.
Each observable is presented with its physical motivation, a complete analytical derivation, and the final formula ready for numerical implementation.

Throughout, we use the eigenstates $\ket{\Phi_k}$ with energies $E_k$ and basis expansion coefficients $C_n^{(k)}$, the transition frequencies $\omega_{kl} = (E_l - E_k)/\hb$, and the wavepacket $\ket{\Psi(t)} = \sum_k c_k\,e^{-iE_k t/\hb}\ket{\Phi_k}$ with real coefficients $c_k$.

% ============================================================
\subsection{Position and Momentum Phase Space}\label{sec:phase_space}
% ============================================================

% ------------------------------------------------------------
\subsubsection{Expectation value of momentum $\expval{\hat{p}}(t)$}\label{sec:pexp}
% ------------------------------------------------------------

\begin{remark}[Why implement?]
The momentum expectation value $\expval{\hat{p}}(t)$ completes the phase-space portrait of the tunneling dynamics.
Together with $\expval{\hat{x}}(t)$, it defines a classical-like trajectory $(x(t),p(t))$ in phase space.
Moreover, through Ehrenfest's theorem, $\expval{\hat{p}}(t)$ provides an independent consistency check: $d\expval{\hat{x}}/dt = \expval{\hat{p}}/\mu$.
A nonzero $\expval{\hat{p}}(t)$ at the barrier region signals active probability flux through the classically forbidden zone.
\end{remark}

\paragraph{Matrix elements of $\hat{p}$ in the HO basis.}

From \cref{eq:x_and_ddx}, the momentum operator is
\begin{equation}
  \hat{p} = i\hb\sqrt{\frac{\alp}{2}}\left(\adag - \hat{a}\right).
\end{equation}
Acting on a basis state:
\begin{equation}
  \hat{p}\ket{m} = i\hb\sqrt{\frac{\alp}{2}}\left(\sqrt{m+1}\ket{m+1} - \sqrt{m}\ket{m-1}\right).
\end{equation}
Therefore the matrix elements are
\begin{equation}\label{eq:p_mel}
  \mel{n}{\hat{p}}{m}
  = i\hb\sqrt{\frac{\alp}{2}}\left[\sqrt{m+1}\,\delta_{n,m+1} - \sqrt{m}\,\delta_{n,m-1}\right].
\end{equation}
In compact form, using $m_* = \minn$:
\begin{equation}\label{eq:p_mel_compact}
  \mel{n}{\hat{p}}{m} =
  \begin{cases}
    +\,i\hb\sqrt{\dfrac{\alp}{2}}\sqrt{m_*+1} & \text{if } n = m+1,\\[8pt]
    -\,i\hb\sqrt{\dfrac{\alp}{2}}\sqrt{m_*+1} & \text{if } n = m-1,\\[8pt]
    0 & \text{otherwise.}
  \end{cases}
\end{equation}
The matrix is \emph{tridiagonal} and \emph{purely imaginary antisymmetric}: $\mel{n}{\hat{p}}{m} = -\mel{m}{\hat{p}}{n}$.

\paragraph{Eigenstate matrix elements.}

The matrix elements between eigenstates are
\begin{align}
  \mel{\Phi_i}{\hat{p}}{\Phi_j}
  &= \sum_{n,l} C_n^{(i)} C_l^{(j)} \mel{n}{\hat{p}}{l}\notag\\
  &= i\hb\sqrt{\frac{\alp}{2}} \sum_n \sqrt{n+1}\left[C_{n+1}^{(i)} C_n^{(j)} - C_n^{(i)} C_{n+1}^{(j)}\right].
  \label{eq:p_eigenstate}
\end{align}
Since the eigenstates can be chosen real, the quantity $P_{ij} \coloneqq \mel{\Phi_i}{\hat{p}}{\Phi_j}$ is \emph{purely imaginary} and antisymmetric: $P_{ij} = -P_{ji}$, with $P_{ii} = 0$ for all $i$.
We write $P_{ij} = i\,\tilde{P}_{ij}$ where $\tilde{P}_{ij} \in \mathbb{R}$ and $\tilde{P}_{ij} = -\tilde{P}_{ji}$:
\begin{equation}\label{eq:Ptilde}
  \tilde{P}_{ij} = \hb\sqrt{\frac{\alp}{2}} \sum_n \sqrt{n+1}\left[C_{n+1}^{(i)} C_n^{(j)} - C_n^{(i)} C_{n+1}^{(j)}\right].
\end{equation}

\paragraph{Derivation of $\expval{\hat{p}}(t)$.}

Starting from the wavepacket expansion:
\begin{align}
  \expval{\hat{p}}(t)
  &= \mel{\Psi(t)}{\hat{p}}{\Psi(t)}
   = \sum_{i,j} c_i c_j\,e^{i(E_i - E_j)t/\hb}\,\mel{\Phi_i}{\hat{p}}{\Phi_j}\notag\\
  &= \sum_{i,j} c_i c_j\,e^{i(E_i - E_j)t/\hb}\,i\,\tilde{P}_{ij}.
\end{align}
Since $P_{ii} = 0$, the diagonal terms vanish identically.
For the off-diagonal terms, we pair $(i,j)$ with $(j,i)$:
\begin{align}
  \expval{\hat{p}}(t)
  &= i\sum_{i<j} c_i c_j\,\tilde{P}_{ij}\left[e^{i\omega_{ji}t} - e^{-i\omega_{ji}t}\right]
     \qquad\text{(using $\tilde{P}_{ji} = -\tilde{P}_{ij}$)}\notag\\
  &= i\sum_{i<j} c_i c_j\,\tilde{P}_{ij}\cdot 2i\sin(\omega_{ji}t)\notag\\
  &= -2\sum_{i<j} c_i c_j\,\tilde{P}_{ij}\,\sin\!\left(\frac{(E_j - E_i)t}{\hb}\right).
\end{align}

\begin{result}[Momentum expectation value]\label{res:pexp}
\begin{equation}\label{eq:pexp}
  \boxed{
  \expval{\hat{p}}(t)
  = -2\sum_{i<j} c_i c_j\,\tilde{P}_{ij}\,\sin\!\left(\frac{(E_j - E_i)t}{\hb}\right),
  }
\end{equation}
where $\tilde{P}_{ij}$ is defined by \cref{eq:Ptilde}.
Note the \emph{sine} dependence (contrast with the \emph{cosine} in $\expval{\hat{x}}(t)$), reflecting the canonical conjugacy of $\hat{x}$ and $\hat{p}$.
\end{result}

\paragraph{Ehrenfest consistency check.}

From \cref{eq:xt_real}, taking the time derivative:
\begin{equation}
  \frac{d\expval{\hat{x}}}{dt}
  = 2\sum_{i<j} c_i c_j\,\mel{\Phi_i}{\hat{x}}{\Phi_j}\,\frac{(E_j - E_i)}{\hb}\,\sin\!\left(\frac{(E_j - E_i)t}{\hb}\right).
\end{equation}
Ehrenfest's theorem requires $\mu\,d\expval{\hat{x}}/dt = \expval{\hat{p}}$, which yields the identity
\begin{equation}\label{eq:ehrenfest_check}
  \tilde{P}_{ij} = -\mu\,\omega_{ji}\,\mel{\Phi_i}{\hat{x}}{\Phi_j},
\end{equation}
providing a nontrivial numerical validation of the eigenvector coefficients.

\paragraph{Selection rules (symmetric case).}
Since $\hat{p}$ is an \emph{odd} operator under parity ($\hat{P}\hat{p}\hat{P}^{-1} = -\hat{p}$), one has:
\begin{itemize}
  \item $\mel{\Phi_i}{\hat{p}}{\Phi_i} = 0$ for all $i$ (diagonal elements vanish);
  \item $\tilde{P}_{ij} \neq 0$ only when $\Phi_i$ and $\Phi_j$ have \emph{opposite} parity.
\end{itemize}
Both properties hold identically to $\hat{x}$.
For the asymmetric case, neither selection rule applies.

\paragraph{Two-state symmetric wavepacket.}
For $\Psi_L = (\Phi_0 + \Phi_1)/\sqrt{2}$:
\begin{equation}
  \expval{\hat{p}}(t) = -\tilde{P}_{01}\,\sin\!\left(\frac{\Delt\,t}{\hb}\right),
\end{equation}
oscillating with amplitude $|\tilde{P}_{01}|$ and period $\trec$, with a $\pi/2$ phase lag relative to $\expval{\hat{x}}(t)$.

\paragraph{Computational cost.}
Precomputing $\tilde{P}_{ij}$ for $K$ eigenstates: $\mathcal{O}(K^2 N)$.
Time evolution at each time step: $\mathcal{O}(K^2)$.
Applies to both symmetric and asymmetric cases.

% ------------------------------------------------------------
\subsubsection{Position and momentum variances}\label{sec:variances}
% ------------------------------------------------------------

\begin{remark}[Why implement?]
The variances $\sigma_x^2(t)$ and $\sigma_p^2(t)$ quantify the spatial and momentum \emph{spread} of the wavepacket.
During tunneling, $\sigma_x^2(t)$ undergoes dramatic changes: it is small when the particle is localized in one well, reaches a maximum during barrier traversal, and the Heisenberg uncertainty product $\sigma_x\sigma_p$ reveals the ``cost'' of tunneling in terms of quantum uncertainty.
A squeezed state with $\sigma_x\sigma_p > \hb/2$ indicates non-Gaussian deformation of the wavepacket during the tunneling process.
\end{remark}

\paragraph{Position variance $\sigma_x^2(t)$.}

The variance is
\begin{equation}
  \sigma_x^2(t) = \expval{\hat{x}^2}(t) - \expval{\hat{x}}^2(t).
\end{equation}
The expectation value $\expval{\hat{x}}(t)$ is already given by \cref{eq:xt_real}.
For $\expval{\hat{x}^2}(t)$, the derivation follows the same pattern with $\hat{x}^2$ replacing $\hat{x}$:
\begin{equation}\label{eq:x2exp}
  \expval{\hat{x}^2}(t)
  = \sum_i |c_i|^2 \mel{\Phi_i}{\hat{x}^2}{\Phi_i}
  + 2\sum_{i<j} c_i c_j\,\mel{\Phi_i}{\hat{x}^2}{\Phi_j}\cos\!\left(\frac{(E_j - E_i)t}{\hb}\right),
\end{equation}
where the eigenstate matrix elements are computed using the basis expansion and the known matrix elements of $\hat{x}^2$ from \cref{eq:x2_HO}:
\begin{align}
  \mel{\Phi_i}{\hat{x}^2}{\Phi_j}
  &= \sum_{n,l} C_n^{(i)} C_l^{(j)} \mel{n}{\hat{x}^2}{l}\notag\\
  &= \frac{1}{\alp}\sum_n C_n^{(i)}C_n^{(j)}\left(n + \half\right)\notag\\
  &\quad + \frac{1}{2\alp}\sum_n \sqrt{(n+1)(n+2)}\left[C_n^{(i)}C_{n+2}^{(j)} + C_{n+2}^{(i)}C_n^{(j)}\right].
  \label{eq:x2_eigenstate}
\end{align}
Note that $\hat{x}^2$ is an \emph{even} operator, so in the symmetric case $\mel{\Phi_i}{\hat{x}^2}{\Phi_j} \neq 0$ only when $\Phi_i$ and $\Phi_j$ have the \emph{same} parity. In particular, the diagonal elements $\mel{\Phi_i}{\hat{x}^2}{\Phi_i}$ are always nonzero.

\paragraph{Matrix elements of $\hat{p}^2$ in the HO basis.}

We derive $\mel{n}{\hat{p}^2}{m}$ from the ladder operator representation.
Starting from $\hat{p} = i\hb\sqrt{\alp/2}(\adag - \hat{a})$:
\begin{align}
  \hat{p}^2
  &= \left[i\hb\sqrt{\frac{\alp}{2}}\right]^2 (\adag - \hat{a})^2
   = -\frac{\hb^2\alp}{2}\,(\adag - \hat{a})^2.
\end{align}
Expanding $(\adag - \hat{a})^2$:
\begin{align}
  (\adag - \hat{a})^2
  &= (\adag)^2 - \adag\hat{a} - \hat{a}\adag + \hat{a}^2\notag\\
  &= (\adag)^2 + \hat{a}^2 - \hat{N} - (\hat{N}+1)\notag\\
  &= (\adag)^2 + \hat{a}^2 - 2\hat{N} - 1,
\end{align}
where $\hat{N} = \adag\hat{a}$ is the number operator.
Therefore:
\begin{equation}\label{eq:p2_op}
  \hat{p}^2 = -\frac{\hb^2\alp}{2}\left[(\adag)^2 + \hat{a}^2 - 2\hat{N} - 1\right]
  = \frac{\hb^2\alp}{2}\left[2\hat{N} + 1 - \hat{a}^2 - (\adag)^2\right].
\end{equation}
Acting on $\ket{m}$:
\begin{align}
  \hat{p}^2\ket{m}
  &= \frac{\hb^2\alp}{2}\Bigl[(2m+1)\ket{m}
    - \sqrt{m(m-1)}\ket{m-2}\notag\\
  &\qquad\qquad\quad - \sqrt{(m+1)(m+2)}\ket{m+2}\Bigr].
\end{align}
The matrix elements are:
\begin{equation}\label{eq:p2_mel}
  \mel{n}{\hat{p}^2}{m} =
  \begin{cases}
    \dfrac{\hb^2\alp}{2}(2n+1) & n = m,\\[8pt]
    -\dfrac{\hb^2\alp}{2}\sqrt{(m_*+1)(m_*+2)} & |n-m| = 2,\\[8pt]
    0 & \text{otherwise.}
  \end{cases}
\end{equation}

\begin{remark}
The kinetic energy matrix \cref{eq:T_HO} is immediately recovered: $T_{nm} = \mel{n}{\hat{p}^2}{m}/(2\mu)$.
The matrix $\mel{n}{\hat{p}^2}{m}$ is real, symmetric, positive definite, and pentadiagonal---exactly the same structure as $\mel{n}{\hat{x}^2}{m}$ but with opposite sign on the off-diagonal elements.
\end{remark}

\paragraph{Momentum variance $\sigma_p^2(t)$.}

The computation of $\expval{\hat{p}^2}(t)$ follows the same structure as $\expval{\hat{x}^2}(t)$.
The eigenstate matrix elements are:
\begin{align}
  \mel{\Phi_i}{\hat{p}^2}{\Phi_j}
  &= \sum_{n,l} C_n^{(i)} C_l^{(j)} \mel{n}{\hat{p}^2}{l}\notag\\
  &= \frac{\hb^2\alp}{2}\sum_n C_n^{(i)}C_n^{(j)}(2n+1)\notag\\
  &\quad - \frac{\hb^2\alp}{2}\sum_n \sqrt{(n+1)(n+2)}\left[C_n^{(i)}C_{n+2}^{(j)} + C_{n+2}^{(i)}C_n^{(j)}\right].
  \label{eq:p2_eigenstate}
\end{align}
The expectation value is:
\begin{equation}\label{eq:p2exp}
  \expval{\hat{p}^2}(t)
  = \sum_i |c_i|^2 \mel{\Phi_i}{\hat{p}^2}{\Phi_i}
  + 2\sum_{i<j} c_i c_j\,\mel{\Phi_i}{\hat{p}^2}{\Phi_j}\cos\!\left(\frac{(E_j - E_i)t}{\hb}\right),
\end{equation}
and the momentum variance is
\begin{equation}\label{eq:sigmap}
  \sigma_p^2(t) = \expval{\hat{p}^2}(t) - \expval{\hat{p}}^2(t).
\end{equation}

\begin{result}[Heisenberg uncertainty product]\label{res:uncertainty}
At every time $t$, the uncertainty relation
\begin{equation}\label{eq:heisenberg}
  \boxed{\sigma_x(t)\,\sigma_p(t) \geq \frac{\hb}{2}}
\end{equation}
must be satisfied.
For the harmonic oscillator ground state, equality holds ($\sigma_x\sigma_p = \hb/2$).
During tunneling, the product exceeds $\hb/2$ due to the non-Gaussian distortion of the wavepacket, providing a quantitative measure of quantum delocalization.
\end{result}

\paragraph{Computational cost.}
Precomputing $\mel{\Phi_i}{\hat{x}^2}{\Phi_j}$ and $\mel{\Phi_i}{\hat{p}^2}{\Phi_j}$: $\mathcal{O}(K^2 N)$ each.
Time evolution: $\mathcal{O}(K^2)$ per time step (same cost as $\expval{\hat{x}}(t)$).
The squaring in $\expval{\hat{x}}^2(t)$ and $\expval{\hat{p}}^2(t)$ is performed after computing the expectation values, adding negligible cost.
Applies to both symmetric and asymmetric cases.

% ============================================================
\subsection{Well Occupation and Barrier Penetration}\label{sec:well_occupation}
% ============================================================

% ------------------------------------------------------------
\subsubsection{Well occupation probabilities $P_L(t)$ and $P_R(t)$}\label{sec:PLR}
% ------------------------------------------------------------

\begin{remark}[Why implement?]
The well occupation probabilities $P_L(t)$ and $P_R(t)$ provide the most direct characterization of tunneling: they measure the fraction of probability in each well as a function of time.
Unlike $\expval{\hat{x}}(t)$, which gives only the average position, $P_L$ and $P_R$ capture the full probability transfer.
The barrier occupation $P_{\text{barrier}}(t)$ directly quantifies the classically forbidden probability, which is the hallmark of quantum tunneling.
These are the quantities most frequently quoted in the experimental literature.
\end{remark}

\paragraph{Definitions.}

We partition the configuration space into three regions separated by the barrier.
For the symmetric potential, the barrier maximum is at $x = 0$, and we define:
\begin{align}
  P_L(t) &= \int_{-\infty}^{0} |\Psi(x,t)|^2\,dx, \label{eq:PL_def}\\
  P_R(t) &= \int_{0}^{\infty} |\Psi(x,t)|^2\,dx,\label{eq:PR_def}\\
  P_{\text{barrier}}(t) &= \int_{x_2(E)}^{x_3(E)} |\Psi(x,t)|^2\,dx,\label{eq:Pbar_def}
\end{align}
where $x_2(E)$ and $x_3(E)$ are the inner classical turning points at energy $E$ (the classically forbidden region).
Normalization gives $P_L(t) + P_R(t) = 1$ (or $P_L + P_{\text{barrier}} + P_R = 1$ if the barrier region is defined separately as a subset of one of the half-spaces).

\paragraph{Partial overlap matrices.}

We define the partial overlap matrices:
\begin{align}
  L_{kl} &= \int_{-\infty}^{0} \Phi_k(x)\,\Phi_l(x)\,dx,\label{eq:Lkl}\\
  R_{kl} &= \int_{0}^{\infty} \Phi_k(x)\,\Phi_l(x)\,dx,\label{eq:Rkl}\\
  B_{kl} &= \int_{x_2}^{x_3} \Phi_k(x)\,\Phi_l(x)\,dx.\label{eq:Bkl}
\end{align}
These are real symmetric matrices.
By orthonormality of the eigenstates over the full real line:
\begin{equation}\label{eq:LR_completeness}
  L_{kl} + R_{kl} = \int_{-\infty}^{\infty} \Phi_k(x)\,\Phi_l(x)\,dx = \delta_{kl}.
\end{equation}

\paragraph{Derivation of $P_L(t)$.}

Inserting the wavepacket expansion into \cref{eq:PL_def}:
\begin{align}
  P_L(t)
  &= \int_{-\infty}^{0} \left|\sum_k c_k\,e^{-iE_k t/\hb}\,\Phi_k(x)\right|^2 dx\notag\\
  &= \sum_{k,l} c_k c_l\,e^{i(E_k - E_l)t/\hb}\int_{-\infty}^{0} \Phi_k(x)\,\Phi_l(x)\,dx\notag\\
  &= \sum_{k,l} c_k c_l\,e^{i(E_k - E_l)t/\hb}\,L_{kl}.
\end{align}
Separating diagonal and off-diagonal terms:
\begin{align}
  P_L(t)
  &= \sum_k |c_k|^2\,L_{kk}
  + \sum_{k \neq l} c_k c_l\,e^{i\omega_{kl}t}\,L_{kl}\notag\\
  &= \sum_k |c_k|^2\,L_{kk}
  + 2\sum_{k<l} c_k c_l\,L_{kl}\,\cos\!\left(\frac{(E_l - E_k)t}{\hb}\right).
\end{align}

\begin{result}[Well occupation probabilities]\label{res:well_occ}
\begin{align}
  P_L(t) &= \sum_k |c_k|^2\,L_{kk}
  + 2\sum_{k<l} c_k c_l\,L_{kl}\,\cos\!\left(\frac{(E_l - E_k)t}{\hb}\right),\label{eq:PL}\\
  P_R(t) &= \sum_k |c_k|^2\,R_{kk}
  + 2\sum_{k<l} c_k c_l\,R_{kl}\,\cos\!\left(\frac{(E_l - E_k)t}{\hb}\right),\label{eq:PR_occ}\\
  P_{\text{barrier}}(t) &= \sum_k |c_k|^2\,B_{kk}
  + 2\sum_{k<l} c_k c_l\,B_{kl}\,\cos\!\left(\frac{(E_l - E_k)t}{\hb}\right).\label{eq:Pbarrier}
\end{align}
One verifies $P_L(t) + P_R(t) = 1$ using \cref{eq:LR_completeness} and normalization $\sum_k |c_k|^2 = 1$.
\end{result}

\paragraph{Symmetric case: simplifications.}

When the potential is symmetric, eigenstates have definite parity: $\Phi_k(-x) = (-1)^k\Phi_k(x)$.
This implies:
\begin{enumerate}[label=(\roman*)]
  \item \textbf{Diagonal elements}: $L_{kk} = R_{kk} = \frac{1}{2}$ for all $k$ (by the even symmetry of $|\Phi_k(x)|^2$).
  \item \textbf{Off-diagonal elements}: $L_{kl} = -R_{kl}$ when $k$ and $l$ have opposite parity, and $L_{kl} = R_{kl}$ when they have the same parity.
\end{enumerate}
The time-independent part becomes $\sum_k|c_k|^2/2 = 1/2$, as expected for a wavepacket initially localized in one well.

\paragraph{Two-state symmetric wavepacket.}

For $c_0 = c_1 = 1/\sqrt{2}$ (left-localized initial state):
\begin{equation}
  P_L(t) = \frac{1}{2} + L_{01}\cos\!\left(\frac{\Delt\,t}{\hb}\right)
  = \frac{1}{2}\left[1 + \cos\!\left(\frac{\Delt\,t}{\hb}\right)\right]
  = \cos^2\!\left(\frac{\Delt\,t}{2\hb}\right),
\end{equation}
where we used $L_{01} = 1/2$ (provable from parity: $\int_{-\infty}^{0}\Phi_0\Phi_1\,dx = 1/2$, since $\Phi_0\Phi_1$ is odd and $\int_{-\infty}^{\infty}\Phi_0\Phi_1\,dx = 0$, so $\int_{-\infty}^{0} = -\int_{0}^{\infty}$ and each equals $\pm 1/2$ of the norm).

\paragraph{Numerical computation of $L_{kl}$, $R_{kl}$, $B_{kl}$.}

These integrals are computed numerically using the spatial grid already available for wavefunction visualization.
If the eigenstates $\Phi_k(x_i)$ are tabulated on an $N_x$-point grid with spacing $\Delta x$:
\begin{equation}
  L_{kl} \approx \Delta x\sum_{x_i < 0} \Phi_k(x_i)\,\Phi_l(x_i),
  \qquad
  R_{kl} \approx \Delta x\sum_{x_i > 0} \Phi_k(x_i)\,\Phi_l(x_i).
\end{equation}
Alternatively, Simpson's rule or the trapezoidal rule may be used for higher accuracy.

\paragraph{Computational cost.}
Precomputing $L_{kl}$, $R_{kl}$, $B_{kl}$ for $K$ eigenstates: $\mathcal{O}(K^2 N_x)$ (one-time cost).
Time evolution: $\mathcal{O}(K^2)$ per time step.
Applies to both symmetric and asymmetric cases.

% ------------------------------------------------------------
\subsubsection{Localization measure $\mathcal{L}(t)$ and tunneling efficiency $\eta$}\label{sec:localization}
% ------------------------------------------------------------

\begin{remark}[Why implement?]
The localization measure $\mathcal{L}(t)$ provides a single-number characterization of the particle's position: $\mathcal{L} = -1$ means completely left, $\mathcal{L} = +1$ means completely right.
The tunneling efficiency $\eta$ quantifies what fraction of the probability can actually transfer between wells, which is the central question in asymmetric tunneling.
In the symmetric case $\eta = 1$ (complete transfer), while in the asymmetric case $\eta < 1$ measures the ``damage'' done by asymmetry.
\end{remark}

\paragraph{Definition.}

The localization measure is defined as
\begin{equation}\label{eq:localization}
  \mathcal{L}(t) = P_R(t) - P_L(t) = 1 - 2P_L(t).
\end{equation}
Using \cref{eq:PL}:
\begin{align}
  \mathcal{L}(t)
  &= \sum_k |c_k|^2(R_{kk} - L_{kk})
  + 2\sum_{k<l} c_k c_l\,(R_{kl} - L_{kl})\cos\!\left(\frac{(E_l - E_k)t}{\hb}\right).
  \label{eq:L_full}
\end{align}

\paragraph{Time average.}

The time average of the cosine terms vanishes (assuming nondegenerate spectrum):
\begin{equation}\label{eq:L_time_avg}
  \overline{\mathcal{L}} \coloneqq \lim_{T\to\infty}\frac{1}{T}\int_0^T \mathcal{L}(t)\,dt
  = \sum_k |c_k|^2(R_{kk} - L_{kk}).
\end{equation}
This is \emph{zero} in the symmetric case (since $R_{kk} = L_{kk} = 1/2$) but generically \emph{nonzero} in the asymmetric case, reflecting the static bias favoring one well.

\paragraph{Tunneling efficiency.}

The tunneling efficiency is defined as
\begin{equation}\label{eq:eta_def}
  \eta = \frac{\mathcal{L}_{\max} - \mathcal{L}_{\min}}{2} \in [0,1],
\end{equation}
where $\mathcal{L}_{\max}$ and $\mathcal{L}_{\min}$ are the maximum and minimum values of $\mathcal{L}(t)$ over all $t \geq 0$.

\paragraph{Two-level analysis: derivation of $\eta$.}

For a two-state wavepacket in the asymmetric case, from \cref{eq:P_L_asym}:
\begin{equation}
  P_L(t) = 1 - \frac{\Delt^2}{\Delt E^2}\sin^2\!\left(\frac{\Omg\,t}{\hb}\right),
\end{equation}
where $\Delt E = 2\Omg = \sqrt{\Delt^2 + 4\eps_{\text{eff}}^2}$.
Therefore:
\begin{align}
  \mathcal{L}(t) &= 1 - 2P_L(t)
  = -1 + \frac{2\Delt^2}{\Delt E^2}\sin^2\!\left(\frac{\Omg\,t}{\hb}\right)\notag\\
  &= -1 + \frac{\Delt^2}{\Delt E^2}\left[1 - \cos\!\left(\frac{2\Omg\,t}{\hb}\right)\right]\notag\\
  &= \left(\frac{\Delt^2}{\Delt E^2} - 1\right) - \frac{\Delt^2}{\Delt E^2}\cos\!\left(\frac{\Delt E\,t}{\hb}\right).
\end{align}
Using $\Delt E^2 = \Delt^2 + 4\eps_{\text{eff}}^2$, the constant term is $-4\eps_{\text{eff}}^2/\Delt E^2$.
Thus:
\begin{equation}\label{eq:L_twolevel}
  \mathcal{L}(t) = -\frac{4\eps_{\text{eff}}^2}{\Delt E^2}
  - \frac{\Delt^2}{\Delt E^2}\cos\!\left(\frac{\Delt E\,t}{\hb}\right).
\end{equation}
The extrema are:
\begin{align}
  \mathcal{L}_{\max} &= -\frac{4\eps_{\text{eff}}^2}{\Delt E^2} + \frac{\Delt^2}{\Delt E^2}
  = \frac{\Delt^2 - 4\eps_{\text{eff}}^2}{\Delt E^2},\\
  \mathcal{L}_{\min} &= -\frac{4\eps_{\text{eff}}^2}{\Delt E^2} - \frac{\Delt^2}{\Delt E^2}
  = -1.
\end{align}
Therefore:
\begin{equation}\label{eq:eta_twolevel}
  \boxed{\eta = \frac{\mathcal{L}_{\max} - \mathcal{L}_{\min}}{2}
  = \frac{1}{2}\left(\frac{\Delt^2 - 4\eps_{\text{eff}}^2}{\Delt E^2} + 1\right)
  = \frac{\Delt^2}{\Delt E^2}
  = \left(\frac{\Delt}{\Delt E}\right)^2.}
\end{equation}
\begin{description}
  \item[Symmetric limit ($\eps_{\text{eff}} = 0$):] $\Delt E = \Delt$, $\eta = 1$. Complete transfer.
  \item[Strong asymmetry ($\eps_{\text{eff}} \gg \Delt/2$):] $\Delt E \approx 2\eps_{\text{eff}}$, $\eta \approx (\Delt/2\eps_{\text{eff}})^2 \ll 1$. Tunneling is suppressed.
  \item[Resonance ($\eps_{\text{eff}} = \Delt/2$):] $\Delt E = \Delt\sqrt{2}$, $\eta = 1/2$.
\end{description}

\paragraph{Computational cost.}
No additional precomputation beyond $P_L(t)$ and $P_R(t)$; $\eta$ is extracted from the time series.
Applies primarily to the asymmetric case (in the symmetric case $\eta = 1$ trivially).

% ============================================================
\subsection{Spectral and Information Quantities}\label{sec:spectral}
% ============================================================

% ------------------------------------------------------------
\subsubsection{Tunnel splittings $\Delt_k$ and spectral structure}\label{sec:splittings}
% ------------------------------------------------------------

\begin{remark}[Why implement?]
The tunnel splittings $\Delt_k$ are the fundamental energetic fingerprint of quantum tunneling.
Each doublet $(E_{2k}, E_{2k+1})$ corresponds to a different transverse excitation of the double well, and $\Delt_k$ determines the tunneling rate at that energy.
The systematic behavior of $\Delt_k$ with $k$ encodes the energy dependence of the barrier transmission and can be compared with WKB predictions to validate the quality of the variational basis.
\end{remark}

\paragraph{Definitions.}

For a deep symmetric double well, the low-lying eigenvalues come in nearly degenerate pairs.
The \emph{tunnel splitting} of the $k$-th doublet is
\begin{equation}\label{eq:Delta_k}
  \Delt_k = E_{2k+1} - E_{2k}, \qquad k = 0, 1, 2, \ldots
\end{equation}
and the \emph{mean pair energy} is
\begin{equation}\label{eq:Ebar_k}
  \bar{E}_k = \frac{E_{2k} + E_{2k+1}}{2}.
\end{equation}
The \emph{tunneling period} associated with the $k$-th doublet is
\begin{equation}\label{eq:t_r_k}
  \trec^{(k)} = \frac{h}{\Delt_k} = \frac{2\pi\hb}{\Delt_k}.
\end{equation}

\paragraph{WKB prediction.}

The semiclassical (WKB) estimate of the tunnel splitting is
\begin{equation}\label{eq:Delta_WKB}
  \Delt_k^{\text{WKB}} = \frac{\hb\omega_k}{\pi}\exp\!\left(-\frac{S_k}{\hb}\right),
\end{equation}
where $\omega_k = \sqrt{V''(\xe)/\mu}$ is the local harmonic frequency at the minimum (to leading order, the same for all doublets), and $S_k$ is the WKB action integral evaluated at energy $\bar{E}_k$:
\begin{equation}\label{eq:Sk}
  S_k = \int_{x_2(\bar{E}_k)}^{x_3(\bar{E}_k)} \sqrt{2\mu\bigl(V(x) - \bar{E}_k\bigr)}\,dx.
\end{equation}
Here $x_2$ and $x_3$ are the inner classical turning points (see \cref{sec:WKB_action}).

\paragraph{Ratio test.}

The ratio of consecutive splittings
\begin{equation}\label{eq:ratio_test}
  r_k = \frac{\Delt_{k+1}}{\Delt_k}
\end{equation}
measures how rapidly tunneling rates increase with energy.
The WKB prediction is $r_k^{\text{WKB}} = \exp[(S_k - S_{k+1})/\hb]$, which is greater than 1 since $S_{k+1} < S_k$ (the action decreases as energy approaches the barrier top).

\paragraph{Computational cost.}
Extracting $\Delt_k$ is $\mathcal{O}(1)$ once the eigenvalues are known.
Computing $S_k$ requires numerical quadrature: $\mathcal{O}(N_x)$ per doublet.
Applies to both symmetric and asymmetric cases (in the latter, $\Delt_k$ is replaced by the avoided-crossing gap; see \cref{sec:avoided_crossing}).

% ------------------------------------------------------------
\subsubsection{Participation ratio and inverse participation ratio}\label{sec:IPR}
% ------------------------------------------------------------

\begin{remark}[Why implement?]
The participation ratio (PR) answers a fundamental question: \emph{how many eigenstates effectively participate in the wavepacket dynamics?}
For a two-state wavepacket, PR~$=2$; for a highly delocalized wavepacket in the energy basis, PR~$\gg 1$.
The inverse participation ratio (IPR) additionally equals the long-time average of the survival probability $P(t)$, connecting a static property of the initial state to the asymptotic dynamics.
\end{remark}

\paragraph{Definitions.}

The \emph{inverse participation ratio} is
\begin{equation}\label{eq:IPR}
  \text{IPR} = \sum_k |c_k|^4,
\end{equation}
and the \emph{participation ratio} is its reciprocal:
\begin{equation}\label{eq:PR}
  \text{PR} = \frac{1}{\text{IPR}} = \frac{1}{\sum_k |c_k|^4}.
\end{equation}

\paragraph{Connection to the time-averaged survival probability.}

We show that $\text{IPR} = \overline{P(t)}$.
Starting from \cref{eq:Ps_general}:
\begin{equation}
  P(t) = \sum_k |c_k|^4 + 2\sum_{k<l}|c_k|^2|c_l|^2\cos\!\left(\frac{(E_l - E_k)t}{\hb}\right).
\end{equation}
Taking the time average over an infinite interval:
\begin{align}
  \overline{P(t)}
  &= \lim_{T\to\infty}\frac{1}{T}\int_0^T P(t)\,dt\notag\\
  &= \sum_k |c_k|^4
  + 2\sum_{k<l}|c_k|^2|c_l|^2\,\underbrace{\lim_{T\to\infty}\frac{1}{T}\int_0^T\cos(\omega_{kl}t)\,dt}_{=\,0\;\text{for}\;\omega_{kl}\neq 0}\notag\\
  &= \sum_k |c_k|^4 = \text{IPR}.
  \label{eq:IPR_time_avg}
\end{align}
The second equality follows because $\omega_{kl} = (E_l - E_k)/\hb \neq 0$ for $k \neq l$ (assuming a nondegenerate spectrum), and the time average of a cosine with nonzero frequency vanishes:
\begin{equation}
  \lim_{T\to\infty}\frac{1}{T}\int_0^T \cos(\omega t)\,dt
  = \lim_{T\to\infty}\frac{\sin(\omega T)}{\omega T} = 0.
\end{equation}

\begin{result}[IPR--survival probability relation]\label{res:IPR}
\begin{equation}\label{eq:IPR_Pbar}
  \boxed{\text{IPR} = \sum_k |c_k|^4 = \overline{P(t)} = \lim_{T\to\infty}\frac{1}{T}\int_0^T P(t)\,dt.}
\end{equation}
The participation ratio $\text{PR} = 1/\text{IPR}$ gives the effective number of eigenstates contributing to the dynamics.
\end{result}

\paragraph{Special cases.}

\begin{description}
  \item[Two-state wavepacket ($c_0 = c_1 = 1/\sqrt{2}$):] $\text{IPR} = 2\cdot(1/2)^2 = 1/2$, $\text{PR} = 2$.
  \item[Four-state wavepacket ($\vartheta = 45^\circ$):] $c_k = 1/2$ for $k = 0,\ldots,3$, $\text{IPR} = 4\cdot(1/4)^2 = 1/4$, $\text{PR} = 4$.
  \item[Single eigenstate ($c_k = \delta_{k,0}$):] $\text{IPR} = 1$, $\text{PR} = 1$. No dynamics.
  \item[Uniform distribution ($c_k = 1/\sqrt{K}$ for $K$ states):] $\text{IPR} = 1/K$, $\text{PR} = K$.
\end{description}

\paragraph{Energy entropy.}

A complementary measure of the wavepacket's complexity in the energy basis is the \emph{energy entropy}:
\begin{equation}\label{eq:S_E}
  S_E = -\sum_k |c_k|^2 \ln|c_k|^2.
\end{equation}
This is a static quantity (independent of time) that measures the Shannon information content of the initial state preparation.
For a uniform distribution over $K$ states: $S_E = \ln K$.
For a two-state wavepacket: $S_E = \ln 2$.
Unlike PR, $S_E$ is sensitive to the detailed shape of the distribution $|c_k|^2$, not just its second moment.

\paragraph{Computational cost.}
IPR, PR, and $S_E$ are $\mathcal{O}(K)$ operations on the expansion coefficients; no precomputation needed.
Applies to both symmetric and asymmetric cases.

% ------------------------------------------------------------
\subsubsection{Zeno time $t_Z$}\label{sec:zeno}
% ------------------------------------------------------------

\begin{remark}[Why implement?]
The Zeno time $t_Z$ characterizes the \emph{short-time} quadratic decay of the survival probability: for $t \ll t_Z$, the system has barely evolved and $P(t) \approx 1$.
It sets the timescale for the \emph{quantum Zeno effect}: if the system is measured at intervals shorter than $t_Z$, the evolution is effectively frozen.
More practically, $t_Z$ determines the numerical time step required to resolve the initial decay and validates the physical consistency of the computed dynamics.
\end{remark}

\paragraph{Derivation.}

The survival amplitude is
\begin{equation}
  A(t) = \sum_k |c_k|^2\,e^{-iE_k t/\hb}.
\end{equation}
Taylor expanding the exponential around $t = 0$:
\begin{align}
  A(t) &= \sum_k |c_k|^2\left[1 - \frac{iE_k t}{\hb} - \frac{E_k^2 t^2}{2\hb^2} + \mathcal{O}(t^3)\right]\notag\\
  &= 1 - \frac{i\expval{E}t}{\hb} - \frac{\expval{E^2}t^2}{2\hb^2} + \mathcal{O}(t^3),
\end{align}
where we define the energy moments
\begin{equation}
  \expval{E} = \sum_k |c_k|^2 E_k, \qquad \expval{E^2} = \sum_k |c_k|^2 E_k^2.
\end{equation}
The survival probability is $P(t) = |A(t)|^2 = A(t)A^*(t)$.
Computing:
\begin{align}
  A^*(t) &= 1 + \frac{i\expval{E}t}{\hb} - \frac{\expval{E^2}t^2}{2\hb^2} + \mathcal{O}(t^3),\\
  P(t) &= |A(t)|^2 = 1 - \frac{(\expval{E^2} - \expval{E}^2)t^2}{\hb^2} + \mathcal{O}(t^4)\notag\\
  &= 1 - \frac{(\Delt E)^2 t^2}{\hb^2} + \mathcal{O}(t^4),
\end{align}
where $(\Delt E)^2 = \expval{E^2} - \expval{E}^2$ is the \emph{energy variance} of the wavepacket.

\begin{result}[Zeno time]\label{res:zeno}
The short-time behavior of the survival probability is
\begin{equation}\label{eq:P_short}
  P(t) \approx 1 - \left(\frac{t}{t_Z}\right)^2, \qquad t \to 0,
\end{equation}
with the \emph{Zeno time}
\begin{equation}\label{eq:tZ}
  \boxed{t_Z = \frac{\hb}{\Delt E}, \qquad (\Delt E)^2 = \sum_k |c_k|^2 E_k^2 - \left(\sum_k |c_k|^2 E_k\right)^2.}
\end{equation}
\end{result}

\paragraph{Physical interpretation.}
The quadratic decay $P(t) \approx 1 - (t/t_Z)^2$ is a universal feature of quantum mechanics, arising from the uncertainty relation $\Delt E\cdot\Delt t \geq \hb/2$.
If the system is repeatedly measured at intervals $\tau \ll t_Z$, the survival probability after each measurement is $P(\tau) \approx 1 - (\tau/t_Z)^2$, and after $n = t/\tau$ measurements:
\begin{equation}
  P_{\text{Zeno}}(t) \approx \left[1 - \left(\frac{\tau}{t_Z}\right)^2\right]^{t/\tau}
  \xrightarrow{\tau\to 0} 1.
\end{equation}
This is the \emph{quantum Zeno effect}: frequent measurements freeze the evolution.

\paragraph{Two-state wavepacket.}
For $c_0 = c_1 = 1/\sqrt{2}$:
\begin{equation}
  (\Delt E)^2 = \frac{E_0^2 + E_1^2}{2} - \left(\frac{E_0 + E_1}{2}\right)^2 = \frac{(E_1 - E_0)^2}{4} = \frac{\Delt^2}{4},
\end{equation}
so $t_Z = 2\hb/\Delt = \trec/\pi$.

\paragraph{Computational cost.}
Computing $t_Z$: $\mathcal{O}(K)$ using the eigenvalues and expansion coefficients.
This is a single scalar, computed once.
Applies to both symmetric and asymmetric cases.

% ------------------------------------------------------------
\subsubsection{Revival time $t_{\text{rev}}$ and collapse time $t_{\text{coll}}$}\label{sec:revival}
% ------------------------------------------------------------

\begin{remark}[Why implement?]
The revival time $t_{\text{rev}}$ is the timescale at which a complex multi-state wavepacket reconstructs itself after dephasing.
Unlike the recurrence time $\trec$ (which applies to a two-state system), $t_{\text{rev}}$ governs the long-time behavior of multi-state wavepackets and determines the temporal structure of ``quantum carpets.''
The collapse time $t_{\text{coll}}$ quantifies how quickly the initial oscillations are washed out by destructive interference among different frequency components.
Together, $t_{\text{coll}}$ and $t_{\text{rev}}$ define the three-act dynamics: coherent oscillation, collapse, and revival.
\end{remark}

\paragraph{Energy expansion around the mean quantum number.}

Consider a wavepacket populated around a mean quantum number $\bar{k}$.
Taylor-expanding the energies to second order:
\begin{equation}\label{eq:E_taylor}
  E_k \approx \bar{E} + E'(k - \bar{k}) + \frac{1}{2}E''(k - \bar{k})^2,
\end{equation}
where $E' = dE_k/dk|_{\bar{k}}$ and $E'' = d^2E_k/dk^2|_{\bar{k}}$ are evaluated at $k = \bar{k}$.
The phase accumulated by the $k$-th component after time $t$ is
\begin{equation}\label{eq:phase_k}
  \phi_k(t) = \frac{E_k t}{\hb} \approx \frac{\bar{E}t}{\hb} + \frac{E'(k-\bar{k})t}{\hb} + \frac{E''(k-\bar{k})^2 t}{2\hb}.
\end{equation}

\paragraph{Classical period.}

The first-order term $E'(k-\bar{k})t/\hb$ generates a linear phase progression.
All components rephase when this term is a multiple of $2\pi$ for adjacent $k$ values, giving the \emph{classical period}:
\begin{equation}\label{eq:t_cl}
  t_{\text{cl}} = \frac{2\pi\hb}{|E'|}.
\end{equation}
For a double well, this corresponds to the classical oscillation period within a single well.

\paragraph{Revival time.}

The second-order term introduces a quadratic phase.
A \emph{full revival} occurs when the quadratic phase $E''(k-\bar{k})^2 t/(2\hb)$ is a multiple of $2\pi$ for all participating $k$ values.
The most stringent condition is for $|k - \bar{k}| = 1$:
\begin{equation}
  \frac{E''\,t_{\text{rev}}}{2\hb} = 2\pi
  \quad\Longrightarrow\quad
  \boxed{t_{\text{rev}} = \frac{4\pi\hb}{|E''|}.}
  \label{eq:t_rev}
\end{equation}
At $t = t_{\text{rev}}$, the wavepacket approximately reforms its initial shape.

\paragraph{Collapse time.}

The initial coherent oscillations are destroyed when the accumulated quadratic phase spread across the participating states reaches $\sim\pi$.
For a wavepacket with PR participating states, the spread in $k$ is $\delta k \sim \sqrt{\text{PR}}$, and collapse occurs when:
\begin{equation}
  \frac{|E''|(\delta k)^2\,t_{\text{coll}}}{2\hb} \sim \pi
  \quad\Longrightarrow\quad
  \boxed{t_{\text{coll}} \approx \frac{t_{\text{rev}}}{2\pi\sqrt{\text{PR}}} = \frac{2\hb}{|E''|\sqrt{\text{PR}}}.}
  \label{eq:t_coll}
\end{equation}
The relation $t_{\text{coll}} \ll t_{\text{rev}}$ (provided $\text{PR} \gg 1$) guarantees the three-stage dynamics.

\paragraph{Application to the double well.}

For the double-well system, the relevant ``spectrum'' for multi-doublet dynamics is the set of tunnel splittings $\{\Delt_k\}$.
The second derivative of the energy with respect to the doublet index $k$ is approximately
\begin{equation}
  E'' \approx \Delt_{k+1} - \Delt_k,
\end{equation}
since $E_{2k+1} - E_{2k} = \Delt_k$ and the inter-doublet spacing is approximately constant.
Therefore:
\begin{align}
  t_{\text{rev}} &\approx \frac{4\pi\hb}{|\Delt_1 - \Delt_0|},\label{eq:t_rev_DW}\\
  t_{\text{coll}} &\approx \frac{2\hb}{|\Delt_1 - \Delt_0|\sqrt{\text{PR}}}.\label{eq:t_coll_DW}
\end{align}

\paragraph{Computational cost.}
Computing $t_{\text{rev}}$ and $t_{\text{coll}}$: $\mathcal{O}(1)$ from the eigenvalues and PR.
Applies to both symmetric and asymmetric cases (multi-state wavepackets only; for two-state wavepackets, $t_{\text{rev}}$ is undefined since there is no quadratic dispersion).

% ------------------------------------------------------------
\subsubsection{Spectral autocorrelation $C(\omega)$}\label{sec:spectral_autocorr}
% ------------------------------------------------------------

\begin{remark}[Why implement?]
The spectral autocorrelation $C(\omega)$ is the Fourier transform of the survival probability and provides the complete frequency fingerprint of the tunneling dynamics.
Each peak in $C(\omega)$ corresponds to a transition frequency $\omega_{kl}$, with height proportional to the product of occupation weights $|c_k|^2|c_l|^2$.
This is the theoretical analog of a spectroscopic measurement and allows direct identification of which eigenstates participate in the dynamics.
It is also the most robust diagnostic for identifying near-degeneracies and resonances.
\end{remark}

\paragraph{Definition and derivation.}

The spectral autocorrelation is defined as the Fourier transform of $P(t)$:
\begin{equation}\label{eq:C_omega_def}
  C(\omega) = \int_0^{\infty} P(t)\,e^{i\omega t}\,dt.
\end{equation}
Substituting \cref{eq:Ps_general} and using the distributional identity $\int_0^\infty \cos(\omega_0 t)\,e^{i\omega t}\,dt = \pi[\delta(\omega + \omega_0) + \delta(\omega - \omega_0)]$ (for $\omega_0 > 0$), we obtain:
\begin{align}
  C(\omega) &= \pi\sum_k |c_k|^4\,\delta(\omega)\notag\\
  &\quad + \pi\sum_{k<l}|c_k|^2|c_l|^2\left[\delta(\omega - \omega_{kl}) + \delta(\omega + \omega_{kl})\right],
  \label{eq:C_omega}
\end{align}
where $\omega_{kl} = (E_l - E_k)/\hb > 0$ for $l > k$.

\begin{result}[Spectral autocorrelation]\label{res:spectral}
The spectral autocorrelation is a sum of delta functions:
\begin{equation}\label{eq:C_omega_final}
  \boxed{C(\omega) = \pi\,\text{IPR}\cdot\delta(\omega) + \pi\sum_{k<l}|c_k|^2|c_l|^2\left[\delta(\omega - \omega_{kl}) + \delta(\omega + \omega_{kl})\right].}
\end{equation}
The peak at $\omega = 0$ has weight $\text{IPR} = \sum_k|c_k|^4 = \overline{P(t)}$.
Each pair of peaks at $\omega = \pm\omega_{kl}$ has weight $|c_k|^2|c_l|^2$.
\end{result}

\paragraph{Practical computation.}

In practice, $P(t)$ is sampled at $N_t$ discrete time points $t_j = j\,\delta t$, and $C(\omega)$ is obtained via a discrete Fourier transform (FFT).
The delta functions are broadened into peaks of width $\sim 2\pi/T_{\max}$, where $T_{\max} = N_t\,\delta t$ is the total simulation time.
Resolving the tunnel splitting $\Delt_0$ requires $T_{\max} \gg \trec = 2\pi\hb/\Delt_0$, i.e., the simulation must cover several recurrence periods.

The peak positions in $C(\omega)$ give the transition frequencies $\omega_{kl}$ (from which eigenvalue differences $E_l - E_k = \hb\omega_{kl}$ are extracted), and the peak heights give the products $|c_k|^2|c_l|^2$ (from which individual $|c_k|^2$ can sometimes be reconstructed).

\paragraph{Computational cost.}
FFT of the $P(t)$ time series: $\mathcal{O}(N_t \log N_t)$.
No precomputation of matrix elements is needed beyond what is already required for $P(t)$.
Applies to both symmetric and asymmetric cases.

% ============================================================
\subsection{Asymmetric Case: Avoided Crossings and WKB}\label{sec:asym_advanced}
% ============================================================

% ------------------------------------------------------------
\subsubsection{WKB tunneling action integral $S(E)$}\label{sec:WKB_action}
% ------------------------------------------------------------

\begin{remark}[Why implement?]
The WKB action integral $S(E)$ is the single most important semiclassical quantity in tunneling theory: it controls the exponential suppression of barrier penetration through the Gamow factor $e^{-S/\hb}$.
Computing $S(E)$ from the exact potential and comparing with the numerically determined tunnel splittings $\Delt_k$ provides a stringent test of both the variational calculation (exact) and the WKB approximation (semiclassical).
It also enables extrapolation to parameter regimes where the full diagonalization becomes impractical.
\end{remark}

\paragraph{Definition.}

The WKB tunneling action integral at energy $E < \Vb$ is
\begin{equation}\label{eq:S_WKB}
  S(E) = \int_{x_2(E)}^{x_3(E)} \sqrt{2\mu\bigl(V(x) - E\bigr)}\,dx,
\end{equation}
where $x_2(E)$ and $x_3(E)$ are the inner classical turning points, defined by $V(x_2) = V(x_3) = E$ with $x_2 < 0 < x_3$.

\paragraph{Turning points for the symmetric quartic potential.}

For $V(x) = (\Vb/\xe^4)(x^2 - \xe^2)^2$ (with $V(\pm\xe) = 0$ and $V(0) = \Vb$), the equation $V(x) = E$ becomes:
\begin{equation}
  \frac{\Vb}{\xe^4}(x^2 - \xe^2)^2 = E
  \quad\Longrightarrow\quad
  (x^2 - \xe^2)^2 = \frac{\xe^4\,E}{\Vb}.
\end{equation}
Taking the square root (both signs):
\begin{equation}
  x^2 - \xe^2 = \pm\xe^2\sqrt{\frac{E}{\Vb}}
  \quad\Longrightarrow\quad
  x^2 = \xe^2\left(1 \pm \sqrt{\frac{E}{\Vb}}\right).
\end{equation}
This gives four turning points (for $E < \Vb$):
\begin{align}
  \text{Outer turning points:}\quad & x = \pm\xe\sqrt{1 + \sqrt{E/\Vb}},\label{eq:tp_outer}\\
  \text{Inner turning points:}\quad & x = \pm\xe\sqrt{1 - \sqrt{E/\Vb}}.\label{eq:tp_inner}
\end{align}
The classically forbidden (barrier) region is between the two inner turning points:
\begin{equation}\label{eq:barrier_region}
  \boxed{x_2(E) = -\xe\sqrt{1 - \sqrt{E/\Vb}}, \qquad x_3(E) = +\xe\sqrt{1 - \sqrt{E/\Vb}}.}
\end{equation}
Note the symmetry: $x_3 = -x_2$.
As $E \to 0$: $x_{2,3} \to \mp\xe$ (full barrier width).
As $E \to \Vb$: $x_{2,3} \to 0$ (barrier width shrinks to zero).

\paragraph{WKB tunnel splitting formula.}

The semiclassical (instanton) formula for the tunnel splitting is
\begin{equation}\label{eq:WKB_splitting}
  \boxed{\Delt^{\text{WKB}}(E) = \frac{\hb\omega_{\text{local}}}{\pi}\exp\!\left(-\frac{S(E)}{\hb}\right),}
\end{equation}
where $\omega_{\text{local}} = \sqrt{V''(\xe)/\mu} = \sqrt{4b/\mu}$ is the harmonic frequency at the potential minimum (see \cref{eq:alpha_local}).

This formula is derived from the connection formulas at both turning points.
In the forbidden region, the WKB wavefunction decays as $\psi \sim \exp(-\int\kappa\,dx)$ where $\kappa(x) = \sqrt{2\mu(V-E)}/\hb$.
The transmission amplitude through the barrier is $\sim e^{-S/\hb}$, and the splitting arises from the symmetric/antisymmetric boundary conditions at both turning points, contributing the prefactor $\hb\omega/\pi$ from the classical frequency of oscillation in each well.

\paragraph{Asymmetric case.}

When $V(x) = V_{\text{sym}}(x) + \eps x$, the turning points are no longer symmetric about the origin.
One must solve $V(x) = E$ numerically for each energy $E$.
The action integral \cref{eq:S_WKB} is then computed by numerical quadrature over the barrier region $[x_2(E), x_3(E)]$.

In the asymmetric case, the tunneling probability depends on direction:
\begin{align}
  S_{L\to R}(E) &= \int_{x_2}^{x_3} \sqrt{2\mu(V(x) - E)}\,dx,\label{eq:S_WKB_LR}\\
  S_{R\to L}(E) &= S_{L\to R}(E) \qquad\text{(same action, different tunneling rates)}.
\end{align}
Although the action integral itself is the same (it depends on the same classically forbidden region), the effective tunneling rate differs between directions because the asymmetry changes which energy levels are in resonance.

\paragraph{Numerical computation.}

$S(E)$ is computed by standard numerical quadrature (e.g., Simpson's rule) on the existing spatial grid:
\begin{equation}
  S(E) \approx \sum_{x_i \in [x_2,x_3]} \sqrt{2\mu(V(x_i) - E)}\,\Delta x,
\end{equation}
where the sum runs over grid points in the classically forbidden region.
Care must be taken near the turning points where the integrand vanishes (square-root singularity in the derivative); a change of variable $x = x_2 + (x_3 - x_2)\sin^2\theta$ regularizes the integral.

\paragraph{Computational cost.}
Computing $S(E)$ for one energy: $\mathcal{O}(N_x)$.
Computing for all doublets: $\mathcal{O}(K\cdot N_x)$.
Applies to both symmetric and asymmetric cases.

% ------------------------------------------------------------
\subsubsection{Avoided crossing map $\Delt(\eps)$}\label{sec:avoided_crossing}
% ------------------------------------------------------------

\begin{remark}[Why implement?]
The avoided crossing map---the energy levels $E_k(\eps)$ as a function of the asymmetry parameter $\eps$---provides the complete characterization of asymmetric tunneling.
It reveals the \emph{resonance conditions}: specific values of $\eps$ where a left-well level nearly coincides with a right-well level, leading to enhanced tunneling.
Away from these resonances, tunneling is exponentially suppressed.
The minimum gap at each avoided crossing equals twice the tunneling matrix element $2|V_{LR}|$, which is the coupling strength between the two wells.
This map is the ``phase diagram'' of the tunneling problem.
\end{remark}

\paragraph{Construction.}

For a set of asymmetry values $\{\eps_j\}_{j=1}^{N_\eps}$, one constructs and diagonalizes the full Hamiltonian $H^{\text{asym}}(\eps_j)$ from \cref{eq:H_asym} to obtain the eigenvalues $E_k(\eps_j)$ for $k = 0, 1, \ldots, K-1$.

The resulting curves $E_k(\eps)$ exhibit the following structure:
\begin{itemize}
  \item At $\eps = 0$: the eigenvalues come in nearly degenerate pairs $(E_{2k}, E_{2k+1})$ split by $\Delt_k$.
  \item As $|\eps|$ increases: the degeneracy of each pair is progressively lifted. Each eigenstate becomes preferentially localized in one well.
  \item At specific values $\eps = \eps^*$: two levels from \emph{different} pairs approach each other, forming an \emph{avoided crossing} (by the Wigner--von~Neumann non-crossing theorem for a one-parameter family of Hermitian matrices).
\end{itemize}

\paragraph{Avoided crossing gap.}

Near an avoided crossing between levels $k$ and $l$ at $\eps = \eps^*$, the two levels behave as
\begin{equation}
  E_\pm(\eps) = \frac{E_k^{(0)}(\eps) + E_l^{(0)}(\eps)}{2}
  \pm \sqrt{\left(\frac{E_k^{(0)}(\eps) - E_l^{(0)}(\eps)}{2}\right)^2 + |V_{kl}|^2},
\end{equation}
where $E_k^{(0)}(\eps)$ and $E_l^{(0)}(\eps)$ are the diabatic (unperturbed) energies and $V_{kl}$ is the tunneling matrix element.
The minimum gap occurs at $\eps = \eps^*$ where $E_k^{(0)}(\eps^*) = E_l^{(0)}(\eps^*)$:
\begin{equation}\label{eq:avoided_gap}
  \boxed{\Delt_{\text{res}} = E_+(\eps^*) - E_-(\eps^*) = 2|V_{kl}|.}
\end{equation}

\paragraph{Resonance condition.}

For the lowest doublet, the resonance condition within the two-level model (\cref{eq:gap_asym}) is
\begin{equation}
  \eps_{\text{eff}}(\eps^*) = 0
  \quad\Longleftrightarrow\quad
  \eps^*\mel{\Phi_0}{\hat{x}}{\Phi_1} = 0,
\end{equation}
which is satisfied at $\eps^* = 0$ (the symmetric case).
For higher doublets, the resonance conditions involve level crossings between different pairs:
\begin{equation}
  E_{2k}^{(0)}(\eps^*) = E_{2l+1}^{(0)}(\eps^*),
\end{equation}
where $E_k^{(0)}$ denotes the energy in the diabatic (single-well) approximation.

\paragraph{Quantities to extract from the map.}

\begin{enumerate}[label=(\roman*)]
  \item \textbf{Resonance positions} $\eps_j^*$: values of $\eps$ at each avoided crossing.
  \item \textbf{Avoided-crossing gaps} $\Delt_{\text{res}}(\eps_j^*)$: minimum energy separation at each resonance.
  \item \textbf{Tunneling suppression}: for $|\eps - \eps^*| \gg \Delt_{\text{res}}/\xe$, tunneling is exponentially suppressed and the eigenstates are localized.
\end{enumerate}

\paragraph{Computational cost.}
$N_\eps$ full diagonalizations: $\mathcal{O}(N_\eps \cdot N^3)$.
Typically $N_\eps \sim 50$--$200$ is sufficient for a smooth map.
Applies only to the asymmetric case.

% ============================================================
\subsection{Phase Space Representations}\label{sec:phase_space_rep}
% ============================================================

% ------------------------------------------------------------
\subsubsection{Wigner quasi-probability distribution $W(x,p,t)$}\label{sec:wigner}
% ------------------------------------------------------------

\begin{remark}[Why implement?]
The Wigner function $W(x,p,t)$ is the closest quantum analog of a classical phase-space distribution.
It simultaneously encodes position and momentum information and is uniquely qualified to reveal the non-classical features of tunneling: \emph{negative regions} in $W$ signal quantum interference effects that have no classical counterpart.
During tunneling, the Wigner function develops characteristic interference fringes between the two wells, providing a vivid visualization of the ``quantum bridge'' that connects the classically disconnected wells.
It is also the natural starting point for semiclassical approximations.
\end{remark}

\paragraph{Definition.}

The Wigner quasi-probability distribution is defined as
\begin{equation}\label{eq:wigner_def}
  W(x,p,t) = \frac{1}{\pi\hb}\int_{-\infty}^{\infty} \Psi^*(x+y,t)\,\Psi(x-y,t)\,e^{2ipy/\hb}\,dy.
\end{equation}
Its marginals reproduce the correct probability distributions:
\begin{equation}
  \int_{-\infty}^{\infty} W(x,p,t)\,dp = |\Psi(x,t)|^2, \qquad
  \int_{-\infty}^{\infty} W(x,p,t)\,dx = |\tilde{\Psi}(p,t)|^2,
\end{equation}
and it is normalized: $\int\!\!\int W\,dx\,dp = 1$.
However, $W$ can take \emph{negative} values---the hallmark of non-classicality.

\paragraph{Decomposition in the eigenstate basis.}

Inserting the wavepacket expansion:
\begin{equation}\label{eq:wigner_expansion}
  W(x,p,t) = \sum_{k,l} c_k\,c_l\,e^{-i(E_k - E_l)t/\hb}\,W_{kl}(x,p),
\end{equation}
where the \emph{cross-Wigner functions} are
\begin{equation}\label{eq:cross_wigner}
  W_{kl}(x,p) = \frac{1}{\pi\hb}\int_{-\infty}^{\infty} \Phi_k(x+y)\,\Phi_l(x-y)\,e^{2ipy/\hb}\,dy.
\end{equation}
Note that $W_{kl}^* = W_{lk}$, and $W_{kk}$ is the Wigner function of the $k$-th eigenstate.

The time evolution in \cref{eq:wigner_expansion} has the same structure as all other observables: the spatial dependence is precomputed in $W_{kl}(x,p)$, and the time evolution is carried entirely by the phase factors $e^{-i\omega_{kl}t}$.
Separating into real and imaginary parts for real $c_k$:
\begin{equation}\label{eq:wigner_time}
  W(x,p,t) = \sum_k |c_k|^2\,W_{kk}(x,p)
  + 2\sum_{k<l} c_k c_l\,\text{Re}\!\left[W_{kl}(x,p)\,e^{-i\omega_{kl}t}\right].
\end{equation}

\paragraph{Wigner function of HO eigenstates.}

For the harmonic oscillator basis states $\ket{n}$ with parameter $\alp$, the Wigner function has the well-known closed form:
\begin{equation}\label{eq:wigner_HO}
  W_n(x,p) = \frac{(-1)^n}{\pi\hb}\,e^{-2\xi}\,L_n(4\xi),
\end{equation}
where $L_n$ is the $n$-th Laguerre polynomial, and
\begin{equation}\label{eq:xi_def}
  \xi = \frac{1}{\hb}\left(\frac{p^2}{2\mu\omega_\alp} + \frac{\mu\omega_\alp\,x^2}{2}\right)
  = \frac{\alp x^2}{2} + \frac{p^2}{2\hb^2\alp},
\end{equation}
with $\omega_\alp = \hb\alp/\mu$ being the angular frequency associated with the basis parameter $\alp$.

\paragraph{Derivation.}

We verify \cref{eq:wigner_HO} for the ground state $n = 0$.
With $\varphi_0(x) = (\alp/\pi)^{1/4}e^{-\alp x^2/2}$:
\begin{align}
  W_0(x,p) &= \frac{1}{\pi\hb}\int_{-\infty}^{\infty} \varphi_0(x+y)\,\varphi_0(x-y)\,e^{2ipy/\hb}\,dy\notag\\
  &= \frac{1}{\pi\hb}\sqrt{\frac{\alp}{\pi}}\int_{-\infty}^{\infty} e^{-\alp(x+y)^2/2}\,e^{-\alp(x-y)^2/2}\,e^{2ipy/\hb}\,dy\notag\\
  &= \frac{1}{\pi\hb}\sqrt{\frac{\alp}{\pi}}\,e^{-\alp x^2}\int_{-\infty}^{\infty} e^{-\alp y^2}\,e^{2ipy/\hb}\,dy.
\end{align}
The remaining Gaussian integral evaluates to $\sqrt{\pi/\alp}\,\exp(-p^2/(\hb^2\alp))$:
\begin{equation}
  W_0(x,p) = \frac{1}{\pi\hb}\,\exp\!\left(-\alp x^2 - \frac{p^2}{\hb^2\alp}\right)
  = \frac{1}{\pi\hb}\,e^{-2\xi}\,L_0(4\xi),
\end{equation}
since $L_0(\cdot) = 1$, confirming \cref{eq:wigner_HO} for $n = 0$.
Higher orders follow from the generating function of Laguerre polynomials.

\paragraph{Cross-Wigner functions for HO states.}

For $n \neq m$ (say $n > m$), the cross-Wigner function of two HO eigenstates is
\begin{equation}\label{eq:cross_wigner_HO}
  W_{nm}(x,p) = \frac{(-1)^m}{\pi\hb}\,e^{-2\xi}\,\left(\frac{2}{\hb}\right)^{(n-m)/2}
  \sqrt{\frac{m!}{n!}}\,(q_x + iq_p)^{n-m}\,L_m^{(n-m)}(4\xi),
\end{equation}
where $L_m^{(\beta)}$ is the associated Laguerre polynomial, and
\begin{equation}
  q_x = \sqrt{\alp}\,x, \qquad q_p = \frac{p}{\hb\sqrt{\alp}}.
\end{equation}

\paragraph{Wigner function of an eigenstate.}

Since the eigenstates are expanded in the HO basis, $\Phi_k(x) = \sum_n C_n^{(k)}\varphi_n(x)$, the eigenstate Wigner function is
\begin{equation}\label{eq:wigner_eigenstate}
  W_{kk}(x,p) = \sum_{n,m} C_n^{(k)} C_m^{(k)}\,W_{nm}^{\text{HO}}(x,p),
\end{equation}
where $W_{nm}^{\text{HO}}$ are the HO cross-Wigner functions from \cref{eq:cross_wigner_HO}.
This double sum can be computed efficiently by truncating at the number of basis functions $N$.

\paragraph{Implementation strategy.}

\begin{enumerate}[label=(\roman*)]
  \item Precompute the HO cross-Wigner functions $W_{nm}^{\text{HO}}(x_i,p_j)$ on a 2D grid $(x_i,p_j)$ for $n,m = 0,\ldots,N-1$: cost $\mathcal{O}(N^2 N_x N_p)$.
  \item Contract with eigenvector coefficients to form $W_{kl}(x_i,p_j)$: cost $\mathcal{O}(K^2 N^2 N_x N_p)$.
  \item Time evolution via \cref{eq:wigner_time}: $\mathcal{O}(K^2 N_x N_p)$ per time step.
\end{enumerate}
The dominant cost is step (ii), which can be reduced by exploiting the sparsity of the eigenvectors.

\paragraph{Computational cost.}
Precomputation: $\mathcal{O}(K^2 N^2 N_x N_p)$ (potentially expensive for large grids).
Time evolution: $\mathcal{O}(K^2 N_x N_p)$ per time step.
Applies to both symmetric and asymmetric cases.

% ------------------------------------------------------------
\subsubsection{Husimi $Q$-distribution $Q(x,p,t)$}\label{sec:husimi}
% ------------------------------------------------------------

\begin{remark}[Why implement?]
The Husimi $Q$-function is a ``smoothed'' version of the Wigner function that is guaranteed to be non-negative, making it easier to interpret as a probability distribution in phase space.
It is obtained by Gaussian convolution of $W(x,p)$ (or equivalently by projecting onto coherent states), and its non-negativity makes it ideal for contour-plot visualizations.
While it sacrifices the sub-Planck resolution of the Wigner function, it clearly reveals the classical-like trajectories and the tunneling pathways without the complication of interference oscillations.
\end{remark}

\paragraph{Definition.}

The Husimi $Q$-distribution is defined as
\begin{equation}\label{eq:husimi_def}
  Q(x,p,t) = \frac{1}{2\pi\hb}\left|\braket{z}{\Psi(t)}\right|^2 \geq 0,
\end{equation}
where $\ket{z}$ is a coherent state centered at the phase-space point $(x,p)$, with complex amplitude
\begin{equation}\label{eq:z_def}
  z = \sqrt{\frac{\mu\omega_\alp}{2\hb}}\,x + \frac{ip}{\sqrt{2\mu\hb\omega_\alp}}
  = \sqrt{\frac{\alp}{2}}\,x + \frac{ip}{\hb\sqrt{2\alp}}.
\end{equation}
Here we use the same $\alp$ as the basis parameter to simplify the overlap integrals.

\paragraph{Coherent state overlap with HO eigenstates.}

The overlap of a coherent state $\ket{z}$ with the $n$-th HO eigenstate is
\begin{equation}\label{eq:coh_overlap}
  \braket{z}{n} = e^{-|z|^2/2}\,\frac{z^{*n}}{\sqrt{n!}}.
\end{equation}
This formula is exact for any $n$ and follows from the definition $\ket{z} = e^{-|z|^2/2}\sum_{n=0}^{\infty}(z^n/\sqrt{n!})\ket{n}$.

\paragraph{Computation of $Q(x,p,t)$.}

Inserting the wavepacket expansion:
\begin{align}
  \braket{z}{\Psi(t)}
  &= \sum_k c_k\,e^{-iE_k t/\hb}\braket{z}{\Phi_k}\notag\\
  &= \sum_k c_k\,e^{-iE_k t/\hb}\sum_n C_n^{(k)}\braket{z}{n}\notag\\
  &= e^{-|z|^2/2}\sum_k c_k\,e^{-iE_k t/\hb}\sum_n C_n^{(k)}\frac{z^{*n}}{\sqrt{n!}}.
  \label{eq:husimi_expansion}
\end{align}
Define the \emph{Bargmann representation} of each eigenstate:
\begin{equation}\label{eq:bargmann}
  \mathcal{B}_k(z^*) = \sum_n C_n^{(k)}\frac{z^{*n}}{\sqrt{n!}}.
\end{equation}
Then:
\begin{equation}\label{eq:Q_final}
  \boxed{Q(x,p,t) = \frac{e^{-|z|^2}}{2\pi\hb}\left|\sum_k c_k\,e^{-iE_k t/\hb}\,\mathcal{B}_k(z^*)\right|^2.}
\end{equation}

\paragraph{Relation to the Wigner function.}

The Husimi function is the Gaussian convolution (coarse-graining) of the Wigner function:
\begin{equation}\label{eq:Q_W_relation}
  Q(x,p,t) = \int\!\!\int W(x',p',t)\,G(x-x',p-p')\,dx'\,dp',
\end{equation}
where the Gaussian kernel is
\begin{equation}
  G(x,p) = \frac{1}{\pi\hb}\exp\!\left(-\alp x^2 - \frac{p^2}{\hb^2\alp}\right)
  = \frac{2}{\pi\hb}\,e^{-2\xi}.
\end{equation}
This smoothing eliminates the negative regions of $W$ and limits the phase-space resolution to $\sim\hb$.

\paragraph{Implementation strategy.}

\begin{enumerate}[label=(\roman*)]
  \item Precompute $\mathcal{B}_k(z^*)$ on a 2D grid of $(x_i,p_j)$ points using \cref{eq:bargmann}: cost $\mathcal{O}(K\cdot N\cdot N_x N_p)$.
  \item At each time step, evaluate the sum $\sum_k c_k e^{-iE_k t/\hb}\mathcal{B}_k(z^*)$ and take the modulus squared: cost $\mathcal{O}(K\cdot N_x N_p)$.
\end{enumerate}
The Husimi function is computationally cheaper than the Wigner function since it does not require the double sum over eigenstates.

\paragraph{Computational cost.}
Precomputation: $\mathcal{O}(K\cdot N\cdot N_x N_p)$.
Time evolution: $\mathcal{O}(K\cdot N_x N_p)$ per time step.
Applies to both symmetric and asymmetric cases.

% ============================================================
\subsection{Implementation Priority}\label{sec:priority}
% ============================================================

The following table ranks all new observables by the ratio of physical insight to implementation effort, providing a suggested implementation order.
The ``Effort'' column estimates the coding complexity (L = low, M = medium, H = high), and the ``Insight'' column rates the physical value of the observable (1--5 scale, with 5 being most insightful).

\begin{table}[H]
\centering
\renewcommand{\arraystretch}{1.4}
\begin{tabular}{clcccl}
  \toprule
  \textbf{Rank} & \textbf{Observable} & \textbf{Effort} & \textbf{Insight} & \textbf{Ratio} & \textbf{Case}\\
  \midrule
  1 & $\Delt_k$, $\bar{E}_k$ (tunnel splittings) & L & 5 & 5.0 & Both\\
  2 & IPR, PR, $S_E$ (participation ratio) & L & 4 & 4.0 & Both\\
  3 & Zeno time $t_Z$ & L & 4 & 4.0 & Both\\
  4 & $\expval{\hat{p}}(t)$ (momentum) & L & 4 & 4.0 & Both\\
  5 & $P_L(t)$, $P_R(t)$ (well occupation) & M & 5 & 3.3 & Both\\
  6 & $\mathcal{L}(t)$, $\eta$ (localization/efficiency) & L & 4 & 4.0 & Asym.\\
  7 & $C(\omega)$ (spectral autocorrelation) & M & 4 & 2.7 & Both\\
  8 & $\sigma_x^2(t)$, $\sigma_p^2(t)$ (variances) & M & 3 & 2.0 & Both\\
  9 & $t_{\text{rev}}$, $t_{\text{coll}}$ (revival/collapse) & L & 3 & 3.0 & Both\\
  10 & $S(E)$ (WKB action) & M & 4 & 2.7 & Both\\
  11 & $\Delt(\eps)$ (avoided crossing map) & H & 5 & 2.5 & Asym.\\
  12 & $Q(x,p,t)$ (Husimi function) & H & 4 & 2.0 & Both\\
  13 & $W(x,p,t)$ (Wigner function) & H & 5 & 2.5 & Both\\
  \bottomrule
\end{tabular}
\caption{Implementation priority for additional observables, ranked by insight-to-effort ratio.
\textbf{Effort}: L = few lines of code, no new infrastructure; M = moderate, requires new matrix elements or numerical integration; H = significant, requires 2D grids and/or repeated diagonalization.
\textbf{Case}: Both = symmetric and asymmetric; Asym.\ = primarily useful for the asymmetric case.}
\label{tab:priority}
\end{table}

\begin{remark}
The first nine observables (ranks 1--9) can be implemented with minimal effort using quantities already available from the diagonalization step (eigenvalues $E_k$, expansion coefficients $c_k$, and eigenvector coefficients $C_n^{(k)}$).
Together, they provide a comprehensive characterization of the tunneling dynamics.
The phase-space representations (ranks 12--13) offer the richest visualization but require substantial computational investment in 2D grid calculations.
A practical strategy is to implement all low-effort observables first, then add the phase-space representations as visualization tools.
\end{remark}


\section{Summary of Key Results}\label{sec:summary}
% ============================================================
% ============================================================

\begin{center}
\renewcommand{\arraystretch}{1.8}
\begin{longtable}{p{4.5cm}p{10cm}}
  \toprule
  \textbf{Quantity} & \textbf{Formula} \\
  \midrule
  \endfirsthead
  \toprule
  \textbf{Quantity} & \textbf{Formula} \\
  \midrule
  \endhead
  \bottomrule
  \endfoot
  \multicolumn{2}{l}{\textit{Potential}} \\
  Symmetric DW
    & $V = ax^4 - bx^2$;\quad $x_{\min} = \pm\sqrt{b/(2a)}$;\quad $\Vb = b^2/(4a)$ \\
  Asymmetric DW
    & $V = ax^4 - bx^2 + cx$ \\
  $a, b$ from $\xe, \Vb$
    & $a = \Vb/\xe^4$;\quad $b = 2\Vb/\xe^2$ \\
  Local frequency
    & $\omega_{\text{loc}} = \sqrt{4b/\mu}$ \\
  Reduced mass (NH$_3$)
    & $\mu = 3m_Hm_N/(3m_H+m_N)$ \\[4pt]
  \midrule
  \multicolumn{2}{l}{\textit{Basis I: HO}} \\
  $\mel{m}{x}{n}$
    & $\sqrt{m_*\!+\!1}/\sqrt{2\alpha}\;\delta_{|m-n|,1}$ \\
  $\mel{m}{x^2}{n}$
    & $(2n\!+\!1)/(2\alpha)\;\delta_{mn}$ $+$ $\sqrt{(m_*\!+\!1)(m_*\!+\!2)}/(2\alpha)\;\delta_{|m-n|,2}$ \\
  $\mel{m}{x^4}{n}$ (diag.)
    & $3(2n^2\!+\!2n\!+\!1)/(4\alpha^2)$ \\
  $T_{nn}$
    & $\alpha(2n\!+\!1)/(4\mu)$ \\
  $\alpha_{\text{local}}$
    & $2\sqrt{\mu b}$ \\
  $\alpha_{\text{quartic}}$
    & $(6\mu\Vb/\xe^4)^{1/3}$ \\[4pt]
  \midrule
  \multicolumn{2}{l}{\textit{Basis III: Gaussians}} \\
  Overlap $S_{ij}$
    & $\sqrt{\pi/\Aij}\,e^{-\Qij}$ \\
  $\Aij$, $\Xij$, $\Qij$
    & $\alpha_i\!+\!\alpha_j$;\; $(\alpha_ix_i\!+\!\alpha_jx_j)/\Aij$;\; $\alpha_i\alpha_j(x_i\!-\!x_j)^2/\Aij$ \\
  $V_{ij}$ (sym.)
    & $e^{-\Qij}\sqrt{\pi/\Aij}\{a[\Xij^4+3\Xij^2/\Aij+3/(4\Aij^2)]-b[\Xij^2+1/(2\Aij)]\}$ \\
  $T_{ij}$
    & $(\alpha_j/\mu)S_{ij}[1-\alpha_j/\Aij-2\alpha_j\alpha_i^2(x_i-x_j)^2/\Aij^2]$ \\[4pt]
  \midrule
  \multicolumn{2}{l}{\textit{Basis IV: DVR}} \\
  $T_{jj}$
    & $\hb^2\pi^2/(6\mu\Delta x^2)$ \\
  $T_{jl}\;(j\neq l)$
    & $\hb^2(-1)^{j-l}/[\mu\Delta x^2(j-l)^2]$ \\
  $V_{jl}$
    & $V(x_j)\,\delta_{jl}$ \\[4pt]
  \midrule
  \multicolumn{2}{l}{\textit{Dynamics}} \\
  Tunneling period
    & $\ttun = \pi\hb/(E_1-E_0) = h/[2(E_1-E_0)]$ \\
  Recurrence time
    & $\trec = 2\pi\hb/(E_1-E_0) = h/(E_1-E_0)$ \\
  Survival prob.
    & $\sum_k|c_k|^4+2\sum_{k<l}|c_k|^2|c_l|^2\cos[(E_l-E_k)t/\hb]$ \\
  $\langle\hat{x}\rangle(t)$
    & $\sum_k|c_k|^2x_{kk}+2\sum_{j<k}\text{Re}[c_j^*c_ke^{-i\omega_{jk}t}]x_{jk}$ \\
  $\langle\hat{p}\rangle(t)$
    & $\mu\,d\langle\hat{x}\rangle/dt$ \\
  Splitting (cm$^{-1}$)
    & $\Delta\nu = (E_1-E_0)\times 219474.63$ \\[4pt]
  \bottomrule
\end{longtable}
\end{center}

% ============================================================
% ============================================================

\section{Extension to Two and Three Dimensions}\label{sec:multidim}
% ============================================================
% ============================================================

This section develops the complete analytical framework for extending the
one-dimensional double-well tunneling problem to two and three spatial
dimensions. We treat three classes of potentials in increasing order of
complexity: (i)~fully separable potentials, (ii)~radially symmetric
(non-separable) potentials, and (iii)~the general polynomial case.  In every
case, the basis consists of tensor products of the one-dimensional harmonic
oscillator (HO) eigenstates already employed in Section~\ref{sec:HO}, and
\emph{all} Hamiltonian matrix elements remain analytically expressible in
closed form.

% ============================================================
\subsection{Tensor-Product Basis Construction}\label{sec:basis_multidim}
% ============================================================

\subsubsection{One-dimensional review}

Recall from Section~\ref{sec:HO} that the 1D HO basis function is
\begin{equation}\label{eq:1d_basis_review}
    \phi_n(x) = N_n\, H_n\!\bigl(\sqrt{\alpha_x}\, x\bigr)\,
    e^{-\alpha_x x^2/2},
    \qquad
    N_n = \left(\frac{\alpha_x}{\pi}\right)^{\!1/4}
    \frac{1}{\sqrt{2^n\, n!}},
\end{equation}
where $\alpha_x = m\omega_x/\hb$ is the width parameter along the $x$-axis.
In Dirac notation we write $\ket{n_x}_x \equiv \phi_{n_x}(x)$, satisfying
\begin{equation}
    {}_x\!\braket{n'_x}{n_x}_x = \delta_{n'_x, n_x},
    \qquad
    \sum_{n_x=0}^{\infty} \ket{n_x}_x {}_x\!\bra{n_x} = \hat{\mathbb{1}}_x.
\end{equation}
The associated ladder operators $\hat{a}_x,\hat{a}_x^\dagger$ satisfy
\begin{equation}\label{eq:ladder_x_md}
    \hat{x} = \frac{1}{\sqrt{2\alpha_x}}\bigl(\hat{a}_x + \hat{a}_x^\dagger\bigr),
    \qquad
    [\hat{a}_x, \hat{a}_x^\dagger] = 1.
\end{equation}

\subsubsection{Two-dimensional basis}

For $d=2$ we introduce an independent HO basis along $y$ with its own width
parameter $\alpha_y = m\omega_y/\hb$:
\begin{equation}\label{eq:2d_basis}
    \ket{n_x, n_y} \equiv \ket{n_x}_x \otimes \ket{n_y}_y,
    \qquad n_x, n_y = 0, 1, 2, \dots
\end{equation}
The position-space representation is the product
\begin{equation}
    \Phi_{n_x n_y}(x,y) = \phi_{n_x}(x;\alpha_x)\;\phi_{n_y}(y;\alpha_y).
\end{equation}
This set is orthonormal,
\begin{equation}\label{eq:2d_ortho}
    \braket{n'_x, n'_y}{n_x, n_y}
    = \delta_{n'_x,n_x}\,\delta_{n'_y,n_y},
\end{equation}
and complete in $L^2(\mathbb{R}^2)$.  The $x$- and $y$-ladder operators
commute:
\begin{equation}\label{eq:commuting_ladders}
    [\hat{a}_x, \hat{a}_y] = [\hat{a}_x, \hat{a}_y^\dagger]
    = [\hat{a}_x^\dagger, \hat{a}_y] = [\hat{a}_x^\dagger, \hat{a}_y^\dagger] = 0.
\end{equation}

\paragraph{Truncation and labeling.}
In practice we truncate to $n_x = 0, 1, \dots, N_x - 1$ and
$n_y = 0, 1, \dots, N_y - 1$, yielding a basis of size
$N_\mathrm{2D} = N_x \times N_y$.  A convenient row-major index is
\begin{equation}\label{eq:2d_index}
    I(n_x, n_y) = n_x \, N_y + n_y,
    \qquad I = 0, 1, \dots, N_\mathrm{2D} - 1.
\end{equation}

\subsubsection{Three-dimensional basis}

The extension to $d=3$ is immediate:
\begin{equation}\label{eq:3d_basis}
    \ket{n_x, n_y, n_z} = \ket{n_x}_x \otimes \ket{n_y}_y \otimes \ket{n_z}_z,
\end{equation}
with width parameter $\alpha_z = m\omega_z/\hb$ along $z$.  The truncated
basis has size
\begin{equation}\label{eq:3d_size}
    N_\mathrm{3D} = N_x \times N_y \times N_z,
\end{equation}
and the single-index map is
\begin{equation}\label{eq:3d_index}
    I(n_x, n_y, n_z) = n_x\, N_y N_z + n_y\, N_z + n_z.
\end{equation}

\subsubsection{General $d$-dimensional basis}

For arbitrary dimension $d$,
\begin{equation}\label{eq:dd_basis}
    \ket{\mathbf{n}} = \ket{n_1, n_2, \dots, n_d}
    = \bigotimes_{k=1}^{d} \ket{n_k}_k,
\end{equation}
with total basis size $N_\mathrm{tot} = \prod_{k=1}^{d} N_k$ and ladder
operators satisfying $[\hat{a}_j, \hat{a}_k^\dagger] = \delta_{jk}$.

% ============================================================
\subsection{Separable 2D Case: Double Well Plus Transverse Harmonic
Confinement}\label{sec:separable_2d}
% ============================================================

\subsubsection{Hamiltonian}

Consider the physically motivated potential where the particle tunnels along
$x$ while harmonically confined along $y$:
\begin{equation}\label{eq:V_sep_2d}
    V(x,y) = \underbrace{-a\,x^2 + b\,x^4}_{V_\mathrm{DW}(x)}
    \;+\; \underbrace{\tfrac{1}{2}m\omega_y^2\, y^2}_{V_\mathrm{HO}(y)}.
\end{equation}
The full Hamiltonian is
\begin{equation}\label{eq:H_sep_2d}
    \hat{H} = \underbrace{\hat{T}_x + V_\mathrm{DW}(\hat{x})}_{\hat{H}_x}
    \;+\; \underbrace{\hat{T}_y + \tfrac{1}{2}m\omega_y^2\,\hat{y}^2}_{\hat{H}_y},
\end{equation}
where $\hat{T}_x = -\frac{\hb^2}{2m}\pdv[2]{}{x}$ and
$\hat{T}_y = -\frac{\hb^2}{2m}\pdv[2]{}{y}$.

\subsubsection{Factorization of matrix elements}

Because $\hat{H} = \hat{H}_x \otimes \hat{\mathbb{1}}_y
+ \hat{\mathbb{1}}_x \otimes \hat{H}_y$ and the basis is a tensor product,
the matrix element factorizes completely:
\begin{equation}\label{eq:sep_factorize}
    \boxed{
    \mel{n'_x, n'_y}{\hat{H}}{n_x, n_y}
    = \mel{n'_x}{\hat{H}_x}{n_x}\,\delta_{n'_y,n_y}
    \;+\; \delta_{n'_x,n_x}\,\mel{n'_y}{\hat{H}_y}{n_y}.
    }
\end{equation}

\paragraph{$x$-sector (double-well).}
This is identical to the 1D problem of Section~\ref{sec:HO}:
\begin{equation}\label{eq:Hx_element}
    \mel{n'_x}{\hat{H}_x}{n_x} = \mel{n'_x}{\hat{T}_x}{n_x}
    - a\,\mel{n'_x}{\hat{x}^2}{n_x}
    + b\,\mel{n'_x}{\hat{x}^4}{n_x},
\end{equation}
with every term given analytically by the formulas in
Section~\ref{sec:HO_matrix_elements}.

\paragraph{$y$-sector (harmonic oscillator).}
When $\alpha_y = m\omega_y/\hb$, the $y$-Hamiltonian is diagonal:
\begin{equation}\label{eq:Hy_exact}
    \mel{n'_y}{\hat{H}_y}{n_y} = \hb\omega_y\!\left(n_y + \tfrac{1}{2}\right)
    \delta_{n'_y, n_y}.
\end{equation}
For a variational $\alpha_y \neq m\omega_y/\hb$:
\begin{equation}\label{eq:Hy_variational}
    \mel{n'_y}{\hat{H}_y}{n_y}
    = \mel{n'_y}{\hat{T}_y}{n_y}
    + \tfrac{1}{2}m\omega_y^2\,\mel{n'_y}{\hat{y}^2}{n_y},
\end{equation}
where both terms are given by the 1D formulas (Section~\ref{sec:HO_matrix_elements}).

\subsubsection{Kronecker-product matrix form}

The full 2D Hamiltonian matrix takes the elegant Kronecker-product form
\begin{equation}\label{eq:H_kron}
    \mathbf{H}_\mathrm{2D} = \mathbf{H}_x \otimes \mathbf{I}_{N_y}
    + \mathbf{I}_{N_x} \otimes \mathbf{H}_y,
\end{equation}
where $\mathbf{H}_x$ is the $N_x \times N_x$ 1D double-well matrix,
$\mathbf{H}_y$ is the diagonal $N_y \times N_y$ matrix with entries
$\hb\omega_y(n_y + \tfrac{1}{2})$, and $\mathbf{I}_N$ is the $N\times N$
identity.

\paragraph{Key consequence.}
The eigenvalues are
\begin{equation}\label{eq:sep_eigenvalues}
    E_{k, n_y} = E_k^{(x)} + \hb\omega_y\!\left(n_y + \tfrac{1}{2}\right),
\end{equation}
so the tunnel splitting $\Delta E = E_1^{(x)} - E_0^{(x)} = \Delta E_\mathrm{1D}$
is \emph{independent} of the transverse quantum number $n_y$.  This provides
an important benchmark.

% ============================================================
\subsection{Non-Separable 2D Case: Radially Symmetric Double
Well}\label{sec:radial_2d}
% ============================================================

\subsubsection{Potential and its expansion}

The radially symmetric 2D double-well potential is
\begin{equation}\label{eq:V_radial_2d}
    V(x,y) = -a\bigl(x^2 + y^2\bigr) + b\bigl(x^2 + y^2\bigr)^2.
\end{equation}
Expanding the quartic term:
\begin{equation}\label{eq:V_radial_2d_expand}
    V(x,y) = -a\,x^2 - a\,y^2 + b\,x^4 + 2b\,x^2 y^2 + b\,y^4.
\end{equation}
Compared to the separable case, the new ingredient is the \emph{cross term}
$2b\,x^2 y^2$, which couples the $x$ and $y$ degrees of freedom.

\subsubsection{Derivation of the cross-term matrix element}

\paragraph{Step 1: Operator factorization.}
Since $\hat{x}$ acts only on $x$-space and $\hat{y}$ only on $y$-space,
\begin{equation}\label{eq:x2y2_factor_op}
    \hat{x}^2\hat{y}^2
    = \bigl(\hat{x}^2 \otimes \hat{\mathbb{1}}_y\bigr)
    \bigl(\hat{\mathbb{1}}_x \otimes \hat{y}^2\bigr)
    = \hat{x}^2 \otimes \hat{y}^2.
\end{equation}

\paragraph{Step 2: Matrix element factorization.}
\begin{align}\label{eq:x2y2_mel}
    \mel{n'_x, n'_y}{\hat{x}^2\hat{y}^2}{n_x, n_y}
    &= \bigl({}_x\!\bra{n'_x} \otimes {}_y\!\bra{n'_y}\bigr)
       \bigl(\hat{x}^2 \otimes \hat{y}^2\bigr)
       \bigl(\ket{n_x}_x \otimes \ket{n_y}_y\bigr) \notag\\
    &= {}_x\!\mel{n'_x}{\hat{x}^2}{n_x}_x
       \;\times\;
       {}_y\!\mel{n'_y}{\hat{y}^2}{n_y}_y.
\end{align}
\begin{equation}\label{eq:x2y2_boxed}
    \boxed{
    \mel{n'_x, n'_y}{\hat{x}^2\hat{y}^2}{n_x, n_y}
    = \mel{n'_x}{\hat{x}^2}{n_x} \cdot \mel{n'_y}{\hat{y}^2}{n_y}.
    }
\end{equation}
This is the central result: the cross term factorizes into a \emph{product}
(not a sum) of 1D matrix elements, each known in closed form from
Section~\ref{sec:HO_matrix_elements}.  The factorization has been verified
numerically with SymPy for multiple index combinations.

\paragraph{Step 3: Explicit evaluation.}
From the ladder operator expansion
$\hat{q}^2 = \tfrac{1}{2\alpha_q}\bigl(\hat{a}_q^2 + 2\hat{n}_q + 1
+ (\hat{a}_q^\dagger)^2\bigr)$, the non-zero matrix elements
(selection rule $\Delta n_q \in \{0,\pm 2\}$) are
\begin{align}
    \mel{n_q}{\hat{q}^2}{n_q} &= \frac{2n_q + 1}{2\alpha_q},
    \label{eq:q2_diag_md} \\[4pt]
    \mel{n_q + 2}{\hat{q}^2}{n_q} &= \frac{\sqrt{(n_q+1)(n_q+2)}}{2\alpha_q},
    \label{eq:q2_up_md} \\[4pt]
    \mel{n_q - 2}{\hat{q}^2}{n_q} &= \frac{\sqrt{n_q(n_q-1)}}{2\alpha_q}.
    \label{eq:q2_down_md}
\end{align}

\paragraph{Step 4: All 9 non-zero cross-term channels.}
The product~\eqref{eq:x2y2_boxed} is non-zero only when both factors are
non-zero, i.e., $\Delta n_x \in \{0, \pm 2\}$ \textbf{and}
$\Delta n_y \in \{0, \pm 2\}$ simultaneously ($3 \times 3 = 9$ channels).
Setting $f_0(n) \coloneqq 2n+1$, $f_+(n) \coloneqq \sqrt{(n+1)(n+2)}$,
$f_-(n) \coloneqq \sqrt{n(n-1)}$:

\begin{center}
\renewcommand{\arraystretch}{1.4}
\begin{tabular}{ccc}
\toprule
$\Delta n_x$ & $\Delta n_y$ &
  Matrix element $\bigl[\times\,(4\alpha_x\alpha_y)^{-1}\bigr]$ \\
\midrule
$0$    & $0$    & $f_0(n_x)\,f_0(n_y)$ \\
$0$    & $+2$   & $f_0(n_x)\,f_+(n_y)$ \\
$0$    & $-2$   & $f_0(n_x)\,f_-(n_y)$ \\
$+2$   & $0$    & $f_+(n_x)\,f_0(n_y)$ \\
$+2$   & $+2$   & $f_+(n_x)\,f_+(n_y)$ \\
$+2$   & $-2$   & $f_+(n_x)\,f_-(n_y)$ \\
$-2$   & $0$    & $f_-(n_x)\,f_0(n_y)$ \\
$-2$   & $+2$   & $f_-(n_x)\,f_+(n_y)$ \\
$-2$   & $-2$   & $f_-(n_x)\,f_-(n_y)$ \\
\bottomrule
\end{tabular}
\end{center}

\subsubsection{Complete 2D radial Hamiltonian matrix element}

Assembling all terms from
$\hat{H} = \hat{T}_x + \hat{T}_y - a\hat{x}^2 - a\hat{y}^2
+ b\hat{x}^4 + 2b\hat{x}^2\hat{y}^2 + b\hat{y}^4$:
\begin{align}\label{eq:H_2d_radial_full}
    \mel{n'_x, n'_y}{\hat{H}}{n_x, n_y}
    &= \Bigl[\mel{n'_x}{\hat{T}_x}{n_x} - a\,\mel{n'_x}{\hat{x}^2}{n_x}
    + b\,\mel{n'_x}{\hat{x}^4}{n_x}\Bigr]\,\delta_{n'_y,n_y}
    \notag\\
    &\quad+ \Bigl[\mel{n'_y}{\hat{T}_y}{n_y} - a\,\mel{n'_y}{\hat{y}^2}{n_y}
    + b\,\mel{n'_y}{\hat{y}^4}{n_y}\Bigr]\,\delta_{n'_x,n_x}
    \notag\\
    &\quad+ 2b\,\mel{n'_x}{\hat{x}^2}{n_x}\;\mel{n'_y}{\hat{y}^2}{n_y}.
\end{align}
In Kronecker-product matrix notation:
\begin{equation}\label{eq:H_2d_radial_kron}
    \mathbf{H}_\mathrm{2D}^\mathrm{rad}
    = \mathbf{H}_x^\mathrm{DW} \otimes \mathbf{I}_{N_y}
    + \mathbf{I}_{N_x} \otimes \mathbf{H}_y^\mathrm{DW}
    + 2b\;\mathbf{X}^2 \otimes \mathbf{Y}^2,
\end{equation}
where $\mathbf{H}_q^\mathrm{DW}$ is the 1D double-well Hamiltonian matrix and
$\mathbf{X}^2$, $\mathbf{Y}^2$ are the 1D matrices of $\hat{x}^2$,
$\hat{y}^2$.

% ============================================================
\subsection{Three-Dimensional Extension}\label{sec:3d_ext}
% ============================================================

\subsubsection{Separable 3D: double well along $z$, harmonic in $x$ and $y$}

The natural 3D generalization of the ammonia-like problem is tunneling along
one axis ($z$) with harmonic confinement in the transverse plane:
\begin{equation}\label{eq:V_sep_3d}
    V(x,y,z) = \tfrac{1}{2}m\omega_x^2 x^2 + \tfrac{1}{2}m\omega_y^2 y^2
    - a\,z^2 + b\,z^4.
\end{equation}
The Hamiltonian factorizes as
$\hat{H} = \hat{H}_x^\mathrm{HO} + \hat{H}_y^\mathrm{HO} + \hat{H}_z^\mathrm{DW}$,
giving eigenvalues
\begin{equation}
    E_{n_x,n_y,k} = \hb\omega_x\!\left(n_x + \tfrac{1}{2}\right)
    + \hb\omega_y\!\left(n_y + \tfrac{1}{2}\right) + E_k^{(z)}.
\end{equation}
All matrix elements are products of 1D elements times Kronecker deltas.

\subsubsection{Spherically symmetric 3D double well}

The isotropic 3D case is
\begin{equation}\label{eq:V_radial_3d}
    V(\mathbf{r}) = -a\,r^2 + b\,r^4,
    \qquad r^2 = x^2 + y^2 + z^2.
\end{equation}
Expanding $r^4 = (x^2+y^2+z^2)^2$:
\begin{equation}\label{eq:r4_expand}
    r^4 = x^4 + y^4 + z^4 + 2x^2 y^2 + 2y^2 z^2 + 2x^2 z^2.
\end{equation}
The Hamiltonian is
\begin{align}\label{eq:H_3d_radial}
    \hat{H} &= \hat{T}_x + \hat{T}_y + \hat{T}_z
    - a\bigl(\hat{x}^2 + \hat{y}^2 + \hat{z}^2\bigr) \notag\\
    &\quad+ b\bigl(\hat{x}^4 + \hat{y}^4 + \hat{z}^4
    + 2\hat{x}^2\hat{y}^2 + 2\hat{y}^2\hat{z}^2 + 2\hat{x}^2\hat{z}^2\bigr).
\end{align}
Each cross term $\hat{q}_i^2\hat{q}_j^2$ factorizes exactly as in 2D
(eq.~\eqref{eq:x2y2_boxed}), with the third coordinate as a spectator:
\begin{equation}\label{eq:cross_3d}
    \mel{n'_x,n'_y,n'_z}{\hat{x}^2\hat{y}^2}{n_x,n_y,n_z}
    = \mel{n'_x}{\hat{x}^2}{n_x}\;\mel{n'_y}{\hat{y}^2}{n_y}\;
    \delta_{n'_z,n_z}.
\end{equation}
The complete 3D matrix element is
\begin{align}\label{eq:H_3d_mel}
    &\mel{n'_x,n'_y,n'_z}{\hat{H}}{n_x,n_y,n_z} \notag\\
    &\quad= \sum_{q\in\{x,y,z\}}\Bigl[
        \mel{n'_q}{\hat{T}_q}{n_q} - a\,\mel{n'_q}{\hat{q}^2}{n_q}
        + b\,\mel{n'_q}{\hat{q}^4}{n_q}\Bigr]
        \prod_{p\neq q}\delta_{n'_p,n_p}
    \notag\\
    &\quad+ 2b\Bigl[
        \mel{n'_x}{\hat{x}^2}{n_x}\,\mel{n'_y}{\hat{y}^2}{n_y}\,\delta_{n'_z,n_z}
        + \mel{n'_y}{\hat{y}^2}{n_y}\,\mel{n'_z}{\hat{z}^2}{n_z}\,\delta_{n'_x,n_x}
        + \mel{n'_x}{\hat{x}^2}{n_x}\,\mel{n'_z}{\hat{z}^2}{n_z}\,\delta_{n'_y,n_y}
    \Bigr],
\end{align}
or in Kronecker-product notation:
\begin{align}\label{eq:H_3d_kron}
    \mathbf{H}_\mathrm{3D}^\mathrm{rad}
    &= \mathbf{H}_x^\mathrm{DW}\otimes\mathbf{I}_y\otimes\mathbf{I}_z
    + \mathbf{I}_x\otimes\mathbf{H}_y^\mathrm{DW}\otimes\mathbf{I}_z
    + \mathbf{I}_x\otimes\mathbf{I}_y\otimes\mathbf{H}_z^\mathrm{DW}
    \notag\\
    &\quad+ 2b\Bigl(
        \mathbf{X}^2\otimes\mathbf{Y}^2\otimes\mathbf{I}_z
        + \mathbf{I}_x\otimes\mathbf{Y}^2\otimes\mathbf{Z}^2
        + \mathbf{X}^2\otimes\mathbf{I}_y\otimes\mathbf{Z}^2
    \Bigr).
\end{align}

% ============================================================
\subsection{Selection Rules and Sparsity}\label{sec:selection_rules}
% ============================================================

\subsubsection{1D selection rules (review)}

From the ladder operator algebra (Section~\ref{sec:HO_matrix_elements}):
\begin{itemize}
    \item $\hat{x}^2$: $\Delta n \in \{0, \pm 2\}$ \quad
          (at most 3 non-zero elements per column).
    \item $\hat{x}^4$: $\Delta n \in \{0, \pm 2, \pm 4\}$ \quad
          (at most 5).
    \item $\hat{T}$: $\Delta n \in \{0, \pm 2\}$ \quad (at most 3).
\end{itemize}
The 1D Hamiltonian has bandwidth~4 ($|n'-n|\leq 4$), giving at most
\textbf{9 non-zero elements per interior row}.

\subsubsection{2D selection rules --- separable case}

The separable Hamiltonian~\eqref{eq:H_sep_2d} couples $\ket{n_x, n_y}$ to
$\ket{n'_x, n'_y}$ only when
\begin{equation}\label{eq:sel_sep_2d}
    \bigl(|\Delta n_x| \leq 4 \;\text{and}\; \Delta n_y = 0\bigr)
    \quad\text{or}\quad
    \bigl(\Delta n_x = 0 \;\text{and}\; |\Delta n_y| \leq 2\bigr).
\end{equation}
Maximum non-zero elements per row: $9 + 5 - 1 = 13$ (variational $\alpha_y$).

\subsubsection{2D selection rules --- radial case}

The cross term $\hat{x}^2\hat{y}^2$ introduces transitions with
$\Delta n_x \in \{0,\pm 2\}$ \textbf{and} $\Delta n_y \in \{0,\pm 2\}$
simultaneously.  A matrix element is non-zero when at least one of the
following holds:
\begin{enumerate}
    \item $|\Delta n_x| \leq 4$ and $\Delta n_y = 0$
          (from $\hat{T}_x$, $\hat{x}^2$, $\hat{x}^4$),
    \item $\Delta n_x = 0$ and $|\Delta n_y| \leq 4$
          (from $\hat{T}_y$, $\hat{y}^2$, $\hat{y}^4$),
    \item $|\Delta n_x| \leq 2$ and $|\Delta n_y| \leq 2$
          (from $\hat{x}^2\hat{y}^2$).
\end{enumerate}
Condition~3 introduces genuinely new elements at
$(\Delta n_x, \Delta n_y) = (\pm 2, \pm 2)$ not reached by conditions~1
or~2.  Maximum non-zero elements per row: $\sim 21$.

\subsubsection{3D selection rules}

Non-zero channels for the 3D radial Hamiltonian are the union of:
\begin{enumerate}
    \item Single-axis: $|\Delta n_q| \leq 4$ for one axis,
          $\Delta n_p = 0$ for the other two.
    \item Cross terms: $|\Delta n_{q_i}| \leq 2$ and
          $|\Delta n_{q_j}| \leq 2$, $\Delta n_{q_k} = 0$.
\end{enumerate}
Maximum non-zero elements per row: $\sim 37$.

% ============================================================
\subsection{Computational Complexity}\label{sec:complexity}
% ============================================================

\subsubsection{Basis size scaling}

Let $N_b$ be the number of basis functions per dimension.

\begin{center}
\renewcommand{\arraystretch}{1.3}
\begin{tabular}{lccc}
\toprule
 & 1D & 2D & 3D \\
\midrule
Basis size & $N_b$ & $N_b^2$ & $N_b^3$ \\
Matrix dimension & $N_b \times N_b$ & $N_b^2 \times N_b^2$ & $N_b^3 \times N_b^3$ \\
\midrule
$N_b = 10$ & $10$ & $100$ & $1\,000$ \\
$N_b = 20$ & $20$ & $400$ & $8\,000$ \\
$N_b = 30$ & $30$ & $900$ & $27\,000$ \\
$N_b = 50$ & $50$ & $2\,500$ & $125\,000$ \\
\bottomrule
\end{tabular}
\end{center}

\subsubsection{Matrix storage}

\paragraph{Dense storage.}
Storing all $N_\mathrm{tot}^2$ doubles (8 bytes each):

\begin{center}
\renewcommand{\arraystretch}{1.3}
\begin{tabular}{lccc}
\toprule
$N_b$ & 1D & 2D & 3D \\
\midrule
$10$ & $0.8\;\mathrm{kB}$ & $80\;\mathrm{kB}$ & $8\;\mathrm{MB}$ \\
$20$ & $3.2\;\mathrm{kB}$ & $1.3\;\mathrm{MB}$ & $512\;\mathrm{MB}$ \\
$30$ & $7.2\;\mathrm{kB}$ & $6.5\;\mathrm{MB}$ & $5.8\;\mathrm{GB}$ \\
$50$ & $20\;\mathrm{kB}$ & $50\;\mathrm{MB}$ & $125\;\mathrm{GB}$ \\
\bottomrule
\end{tabular}
\end{center}

\paragraph{Sparse storage.}
The number of non-zero elements per row is bounded by a small constant $c_d$
($\leq 9$ in 1D, $\leq 21$ in 2D, $\leq 37$ in 3D), so the total non-zero
count scales only linearly with basis size:
\begin{equation}\label{eq:nnz}
    N_\mathrm{nz} \leq c_d \times N_b^d,
\end{equation}
far better than the $N_b^{2d}$ entries of the dense matrix.  The sparsity
ratio
\begin{equation}\label{eq:sparsity}
    \text{sparsity} = \frac{c_d\, N_b^d}{N_b^{2d}} = \frac{c_d}{N_b^d}
\end{equation}
decreases rapidly: at $N_b = 30$ it is $30\%$ (1D), $0.26\%$ (2D), and
$0.0014\%$ (3D).  The matrices are \emph{extremely sparse} in higher
dimensions, making sparse storage and iterative eigensolvers essential.

\subsubsection{Eigenvalue computation}

\paragraph{Dense diagonalization.}
Full diagonalization (e.g., LAPACK \texttt{dsyev}) scales as
$\mathcal{O}(N_\mathrm{tot}^3)$.  At $N_b = 30$ this is
$\sim 2.7 \times 10^4$ operations in 1D,
$\sim 7.3 \times 10^8$ in 2D, and
$\sim 2.0 \times 10^{13}$ in 3D --- the last rendering dense
diagonalization wholly impractical.

\paragraph{Iterative sparse eigensolver.}
If only the lowest $k$ eigenvalues are needed, Lanczos or Jacobi--Davidson
methods require
\begin{equation}
    \text{Cost} \sim k \cdot c_d \cdot N_b^d \cdot N_\mathrm{iter}
    \ll N_b^{3d}.
\end{equation}
For $N_b = 30$ in 3D with $k=10$ and $N_\mathrm{iter} = 500$:
$\text{Cost} \approx 5 \times 10^9$ --- roughly $4000\times$ faster than
dense diagonalization.

% ============================================================
\subsection{Symmetry Reduction}\label{sec:symmetry_md}
% ============================================================

\subsubsection{Parity symmetry in the separable 2D case}

The separable Hamiltonian commutes with the parity operators along each axis
independently: $\hat{\Pi}_x\colon x \to -x$ with eigenvalue
$\pi_x = (-1)^{n_x}$, and analogously for $\hat{\Pi}_y$.  The Hamiltonian
block-diagonalizes into \textbf{four symmetry sectors}:

\begin{center}
\renewcommand{\arraystretch}{1.3}
\begin{tabular}{cccc}
\toprule
Sector & $\pi_x$ & $\pi_y$ & Basis states \\
\midrule
$(+,+)$ & $+1$ & $+1$ & $n_x$ even, $n_y$ even \\
$(+,-)$ & $+1$ & $-1$ & $n_x$ even, $n_y$ odd  \\
$(-,+)$ & $-1$ & $+1$ & $n_x$ odd,  $n_y$ even \\
$(-,-)$ & $-1$ & $-1$ & $n_x$ odd,  $n_y$ odd  \\
\bottomrule
\end{tabular}
\end{center}

Each sector contains $\approx N_b^2/4$ basis states.  The tunnel splitting
is the gap between the ground state in sector $(+,+)$ and the first
antisymmetric state in sector $(-,+)$ (for ground transverse mode $n_y=0$).

\subsubsection{Additional symmetry in the radial 2D case}

The isotropic radial potential satisfies $V(x,y)=V(y,x)$, so the Hamiltonian
also commutes with the exchange operator
$\hat{P}_{xy}\colon (x,y)\to(y,x)$, enabling further sub-blocking.

Beyond discrete symmetries, the full $\mathrm{SO}(2)$ rotational symmetry
decomposes the Hilbert space into sectors of definite angular momentum $m$
($\hat{L}_z\ket{N,m}=m\hb\ket{N,m}$).  Within the shell
$N = n_x + n_y$, the allowed values are $m = -N, -N+2, \dots, N$.
One may either work in the Cartesian basis exploiting parity + exchange
(four blocks, simpler to implement) or transform to the $(N,m)$ basis and
exploit the full SO(2) symmetry.

\subsubsection{SO(3) symmetry in the 3D radial case}

The spherically symmetric 3D potential commutes with $\hat{\mathbf{L}}^2$
and $\hat{L}_z$.  Writing
$\psi(\mathbf{r}) = R_l(r)\,Y_l^{m_l}(\theta,\varphi)$, the radial
equation
\begin{equation}\label{eq:radial_3d_reduced}
    -\frac{\hb^2}{2m}\dv[2]{u_l}{r}
    + \left[\frac{\hb^2 l(l+1)}{2m r^2} - a\,r^2 + b\,r^4\right]u_l = E\,u_l
\end{equation}
(with $R_l = u_l/r$, $r\geq 0$, $u_l(0)=0$) is independent for each $l$,
and eigenvalues carry $(2l+1)$-fold degeneracy from $m_l$.

In the Cartesian HO basis the full discrete symmetry group is
$(\mathbb{Z}_2)^3 \rtimes S_3$ (hyperoctahedral group, order~48), giving
a reduction factor of up to~48 in sector size.

% ============================================================
\subsection{Analytical Tractability Assessment}\label{sec:tractability}
% ============================================================

\subsubsection{Classification of cases}

\begin{center}
\renewcommand{\arraystretch}{1.4}
\begin{tabular}{p{4.2cm}ccp{5.2cm}}
\toprule
Case & Difficulty & ME analytic? & Key feature \\
\midrule
Separable 2D/3D & Easy & Yes &
  Reduces to 1D; eigenvalues are sums of 1D eigenvalues \\
Radially symmetric 2D & Medium & Yes &
  Cross term $x^2y^2$ factorizes; SO(2) or parity+exchange blocks \\
Radially symmetric 3D & Medium & Yes &
  All $q_i^2 q_j^2$ cross terms factorize; SO(3) or discrete symmetry \\
General asymmetric polynomial & Hard & Yes &
  No symmetry blocks; matrix remains sparse; all ME analytic \\
\bottomrule
\end{tabular}
\end{center}

(ME = matrix elements.)

\subsubsection{Proof that all matrix elements are analytic for polynomial
potentials}

\begin{result}[Analyticity of matrix elements]\label{res:analytic_me}
Let $V(x_1, \dots, x_d) = \sum_{\mathbf{k}} c_\mathbf{k}\,
x_1^{k_1} \cdots x_d^{k_d}$ be any polynomial potential.
Then every matrix element $\mel{\mathbf{n}'}{V}{\mathbf{n}}$ in the
tensor-product HO basis is a finite, algebraically closed-form expression.
\end{result}

\begin{proof}
By the tensor product structure,
\begin{equation}\label{eq:poly_factorize}
    \mel{\mathbf{n}'}{x_1^{k_1}\cdots x_d^{k_d}}{\mathbf{n}}
    = \prod_{j=1}^{d} \mel{n'_j}{\hat{x}_j^{k_j}}{n_j}.
\end{equation}
Each 1D factor is computed by writing
$\hat{x}^k = (2\alpha)^{-k/2}(\hat{a} + \hat{a}^\dagger)^k$
and expanding via the binomial theorem.  After normal ordering, every term
yields a factor $\sqrt{n!/(n-s)!}\,\delta_{n',n+j-s}$ for appropriate
$j,s$ --- a finite expression involving only square roots of integers.
The selection rule is $|\Delta n_j| \leq k_j$ per coordinate.
\end{proof}

\begin{remark}
For the quartic double-well potential in any dimension $d$, all Hamiltonian
matrix elements reduce to products of the 1D formulas in
Sections~\ref{sec:HO_matrix_elements} and~\ref{sec:HO_kinetic}.  No
numerical quadrature is ever required.
\end{remark}

% ============================================================
\subsection{Summary of Formulas for Multi-Dimensional
Implementation}\label{sec:summary_multidim}
% ============================================================

\paragraph{1D building blocks} (per coordinate $q$ with width $\alpha_q$, from
Sections~\ref{sec:HO_matrix_elements}--\ref{sec:HO_kinetic}):
\begin{align}
    \mel{n'}{\hat{T}_q}{n} &= \frac{\hb^2\alpha_q}{4m}
    \Bigl[(2n+1)\,\delta_{n',n} - \sqrt{(n+1)(n+2)}\,\delta_{n',n+2}
    - \sqrt{n(n-1)}\,\delta_{n',n-2}\Bigr],
    \label{eq:summary_T_md} \\[4pt]
    \mel{n'}{\hat{q}^2}{n} &= \frac{1}{2\alpha_q}
    \Bigl[(2n+1)\,\delta_{n',n} + \sqrt{(n+1)(n+2)}\,\delta_{n',n+2}
    + \sqrt{n(n-1)}\,\delta_{n',n-2}\Bigr],
    \label{eq:summary_q2_md} \\[4pt]
    \mel{n'}{\hat{q}^4}{n} &= \frac{1}{4\alpha_q^2}
    \Bigl[(6n^2+6n+3)\,\delta_{n',n}
    + (4n+6)\sqrt{(n+1)(n+2)}\,\delta_{n',n+2}
    \notag\\
    &\quad+ (4n-2)\sqrt{n(n-1)}\,\delta_{n',n-2}
    + \sqrt{(n+1)(n+2)(n+3)(n+4)}\,\delta_{n',n+4}
    \notag\\
    &\quad+ \sqrt{n(n-1)(n-2)(n-3)}\,\delta_{n',n-4}
    \Bigr].
    \label{eq:summary_q4_md}
\end{align}

\paragraph{2D separable} [eq.~\eqref{eq:H_kron}]:
\begin{equation*}
    \mathbf{H}_\mathrm{2D} = \mathbf{H}_x \otimes \mathbf{I}_{N_y}
    + \mathbf{I}_{N_x} \otimes \mathbf{H}_y.
\end{equation*}

\paragraph{2D radial} [eq.~\eqref{eq:H_2d_radial_kron}]:
\begin{equation*}
    \mathbf{H}_\mathrm{2D}^\mathrm{rad}
    = \mathbf{H}_x^\mathrm{DW} \otimes \mathbf{I}
    + \mathbf{I} \otimes \mathbf{H}_y^\mathrm{DW}
    + 2b\,\mathbf{X}^2 \otimes \mathbf{Y}^2.
\end{equation*}

\paragraph{3D radial} [eq.~\eqref{eq:H_3d_kron}]:
three 1D DW Hamiltonians (one per axis) plus three pairwise cross terms
$2b\,\mathbf{Q}_i^2\otimes\mathbf{Q}_j^2$.

\medskip
\noindent All matrix elements above are analytic, exact, and involve only
elementary arithmetic operations on integers and their square roots.  No
numerical integration is required at any stage.

% ============================================================

\section{General Polynomial Potentials}\label{sec:general}
% ============================================================
% ============================================================

The previous sections treated the quartic double-well $V(x)=ax^4-bx^2$ as the
canonical example. This section develops the complete analytical machinery for
an \emph{arbitrary} polynomial potential of degree $N$, an arbitrary number of
basis functions $M$, and user-defined initial wavepackets. The derivations here
underpins the \texttt{QuTu} general-case solver.

% ============================================================
\subsection{General polynomial potential}\label{sec:gen_pot}
% ============================================================

A general real-valued polynomial potential of even degree $N$ (with $N \geq 2$)
is written
\begin{equation}\label{eq:gen_V}
  V(x) = \sum_{k=0}^{N} c_k\,x^k, \qquad c_N > 0,
\end{equation}
with real coefficients $\{c_k\}$.  The requirement $c_N > 0$ ensures
confinement ($V(x)\to+\infty$ as $|x|\to\infty$).  The quartic double-well
of Section~\ref{sec:potential} corresponds to $N=4$, $c_4=a$, $c_2=-b$, and
all other coefficients zero.

\begin{remark}
Odd-degree leading terms are excluded because $V(x)\to-\infty$ along one
direction, yielding an unbounded-below Hamiltonian with no ground state.
Odd \emph{interior} coefficients ($c_1,c_3,\ldots$) are perfectly allowed;
they break the left--right symmetry of the potential.
\end{remark}

% ============================================================
\subsection{Hamiltonian matrix in the harmonic-oscillator basis}\label{sec:gen_H}
% ============================================================

Using the harmonic oscillator (HO) basis $\{\phi_n\}_{n=0}^{M-1}$ introduced
in Section~\ref{sec:HO}, the Hamiltonian matrix elements decompose as
\begin{equation}\label{eq:gen_H_elem}
  H_{mn} = T_{mn} + V_{mn},
\end{equation}
where $T_{mn}=\bra{m}\hat T\ket{n}$ is the kinetic energy matrix element
(derived in Section~\ref{sec:HO}) and
\begin{equation}\label{eq:gen_V_elem}
  V_{mn} = \bra{m}V(\hat x)\ket{n}
         = \sum_{k=0}^{N} c_k\, X^{(k)}_{mn},
  \qquad
  X^{(k)}_{mn} \equiv \bra{m}\hat x^k\ket{n}.
\end{equation}
The problem therefore reduces to computing all power-matrix elements
$X^{(k)}_{mn}$ for $k=0,1,\ldots,N$.

% ============================================================
\subsection{Power-matrix elements via the master recursion}\label{sec:master_rec}
% ============================================================

\subsubsection{Ladder-operator form of \texorpdfstring{$\hat x$}{x}}

Introduce the characteristic length
\begin{equation}\label{eq:ell}
  \ell = \frac{1}{\sqrt{2\mu\alpha}},
\end{equation}
where $\mu$ is the mass and $\alpha$ the HO width parameter.  Then
\begin{equation}\label{eq:x_ladder}
  \hat x = \ell\,(\hat a + \hat a^\dagger),
\end{equation}
with the standard actions $\hat a\ket{n}=\sqrt{n}\,\ket{n-1}$ and
$\hat a^\dagger\ket{n}=\sqrt{n+1}\,\ket{n+1}$.

\subsubsection{Master recursion theorem}

\begin{result}[Master recursion for $X^{(k)}_{mn}$]\label{res:master_rec}
  For $k \geq 1$ and arbitrary non-negative integers $m,n$,
  \begin{equation}\label{eq:master_rec}
    X^{(k)}_{mn}
    = \ell\Bigl[\sqrt{n}\,X^{(k-1)}_{m,n-1}
               + \sqrt{n+1}\,X^{(k-1)}_{m,n+1}\Bigr],
  \end{equation}
  with the initial condition $X^{(0)}_{mn} = \delta_{mn}$.
\end{result}

\begin{proof}
Insert $\hat x = \ell(\hat a+\hat a^\dagger)$ and use linearity:
\[
  X^{(k)}_{mn}
  = \bra{m}\hat x^k\ket{n}
  = \bra{m}\hat x^{k-1}\hat x\ket{n}
  = \ell\,\bra{m}\hat x^{k-1}\bigl(\hat a+\hat a^\dagger\bigr)\ket{n}.
\]
Acting with the ladder operators on $\ket{n}$ gives
\[
  X^{(k)}_{mn}
  = \ell\Bigl[\sqrt{n}\,\bra{m}\hat x^{k-1}\ket{n-1}
             + \sqrt{n+1}\,\bra{m}\hat x^{k-1}\ket{n+1}\Bigr]
  = \ell\Bigl[\sqrt{n}\,X^{(k-1)}_{m,n-1}
             + \sqrt{n+1}\,X^{(k-1)}_{m,n+1}\Bigr]. \qed
\]
\end{proof}

Equation~\eqref{eq:master_rec} allows one to compute the entire matrix
$X^{(k)}$ for any $k$ by simple matrix--vector operations; no additional
integrals are required.

% ============================================================
\subsection{Selection rules and bandwidth}\label{sec:selection}
% ============================================================

\begin{result}[Selection rule]\label{res:selection}
  $X^{(k)}_{mn} = 0$ unless the following two conditions both hold:
  \begin{enumerate}
    \item \textbf{Parity:} $(m-n) \equiv k \pmod{2}$.
    \item \textbf{Reach:} $|m-n| \leq k$.
  \end{enumerate}
\end{result}

\begin{proof}
By induction on $k$.  The base case $k=0$ gives $\delta_{mn}$, satisfying
both rules ($(m-n) \equiv 0 \pmod 2$ and $|m-n|=0\leq 0$).

Assume the rules hold for $k-1$.  A non-zero contribution from
$X^{(k-1)}_{m,n\pm 1}$ requires $(m-(n\pm1))\equiv k-1\pmod 2$ and
$|m-(n\pm 1)|\leq k-1$.  The shift $n\to n\pm 1$ changes the parity of
$m-n$ by one unit, so $m-n\equiv k\pmod 2$; and $|m-n|\leq|m-(n\pm1)|+1\leq
k$.  Hence the rules propagate. \qed
\end{proof}

\begin{result}[Bandwidth]\label{res:bandwidth}
  The Hamiltonian matrix $\mathbf{H}$ for the potential~\eqref{eq:gen_V} is
  \emph{banded} with half-bandwidth $N$:
  \begin{equation}
    H_{mn} = 0 \quad \text{if } |m-n| > N.
  \end{equation}
\end{result}

The bandwidth is determined solely by the polynomial degree; the kinetic
energy matrix $\mathbf{T}$ has half-bandwidth 2, which is dominated by $N$
for $N>2$.  For a degree-4 potential, $\mathbf{H}$ has half-bandwidth 4,
consistent with the explicit formulas in Section~\ref{sec:HO}.

% ============================================================
\subsection{Explicit matrix elements for low powers}\label{sec:explicit_Xk}
% ============================================================

Applying the master recursion~\eqref{eq:master_rec} iteratively yields closed
forms for each $k$.  Below, $\ell=(2\mu\alpha)^{-1/2}$ and
$\delta_{mn}\equiv\delta_{m,n}$.

\subsubsection{\texorpdfstring{$k=1$}{k=1}: position matrix}

\begin{equation}\label{eq:X1}
  X^{(1)}_{mn}
  = \ell\bigl(\sqrt{n}\,\delta_{m,n-1} + \sqrt{n+1}\,\delta_{m,n+1}\bigr).
\end{equation}

\subsubsection{\texorpdfstring{$k=2$}{k=2}: quadratic}

\begin{equation}\label{eq:X2}
  X^{(2)}_{mn}
  = \ell^2\Bigl[(2n+1)\,\delta_{mn}
    + \sqrt{n(n-1)}\,\delta_{m,n-2}
    + \sqrt{(n+1)(n+2)}\,\delta_{m,n+2}\Bigr].
\end{equation}

\subsubsection{\texorpdfstring{$k=3$}{k=3}: cubic}

\begin{align}\label{eq:X3}
  X^{(3)}_{mn}
  &= \ell^3\Bigl[3n^{3/2}\,\delta_{m,n-1}
               + 3(n+1)^{3/2}\,\delta_{m,n+1} \nonumber\\
  &\quad +\sqrt{n(n-1)(n-2)}\,\delta_{m,n-3}
    +\sqrt{(n+1)(n+2)(n+3)}\,\delta_{m,n+3}\Bigr].
\end{align}

\subsubsection{\texorpdfstring{$k=4$}{k=4}: quartic}

\begin{align}\label{eq:X4}
  X^{(4)}_{mn}
  &= \ell^4\Bigl[(6n^2+6n+3)\,\delta_{mn}
    + (4n-2)\sqrt{n(n-1)}\,\delta_{m,n-2}
    + (4n+6)\sqrt{(n+1)(n+2)}\,\delta_{m,n+2} \nonumber\\
  &\quad +\sqrt{n(n-1)(n-2)(n-3)}\,\delta_{m,n-4}
    +\sqrt{(n+1)(n+2)(n+3)(n+4)}\,\delta_{m,n+4}\Bigr].
\end{align}

\subsubsection{\texorpdfstring{$k=5$}{k=5}: quintic}

\begin{align}\label{eq:X5}
  X^{(5)}_{mn}
  &= \ell^5\Bigl[(10n^2+10n+5)n^{1/2}\,\delta_{m,n-1}
              + (10n^2+10n+5)(n+1)^{1/2}\,\delta_{m,n+1} \nonumber\\
  &\quad + 5(2n-1)\sqrt{n(n-1)(n-2)}\,\delta_{m,n-3}
    + 5(2n+3)\sqrt{(n+1)(n+2)(n+3)}\,\delta_{m,n+3} \nonumber\\
  &\quad + \sqrt{n(n-1)(n-2)(n-3)(n-4)}\,\delta_{m,n-5} \nonumber\\
  &\quad + \sqrt{(n+1)(n+2)(n+3)(n+4)(n+5)}\,\delta_{m,n+5}\Bigr].
\end{align}

\subsubsection{\texorpdfstring{$k=6$}{k=6}: sextic}

\begin{align}\label{eq:X6}
  X^{(6)}_{mn}
  &= \ell^6\Bigl[(20n^3+30n^2+40n+15)\,\delta_{mn} \nonumber\\
  &\quad + 3(10n^2-10n+4)\sqrt{n(n-1)}\,\delta_{m,n-2}
    + 3(10n^2+10n+4)\sqrt{(n+1)(n+2)}\,\delta_{m,n+2} \nonumber\\
  &\quad + (15n-15)\sqrt{n(n-1)(n-2)(n-3)}\,\delta_{m,n-4} \nonumber\\
  &\quad + (15n+15)\sqrt{(n+1)(n+2)(n+3)(n+4)}\,\delta_{m,n+4} \nonumber\\
  &\quad + \sqrt{n(n-1)(n-2)(n-3)(n-4)(n-5)}\,\delta_{m,n-6} \nonumber\\
  &\quad + \sqrt{(n+1)(n+2)(n+3)(n+4)(n+5)(n+6)}\,\delta_{m,n+6}\Bigr].
\end{align}

\subsubsection{Summary table}

Table~\ref{tab:Xk_bandwidth} summarises the non-zero diagonal offsets and
leading coefficient for each $k$.

\begin{table}[htbp]
  \centering
  \caption{%
    Non-zero off-diagonal levels and leading prefactors for
    $X^{(k)}_{mn}$ ($k=1,\ldots,6$).
    $\ell=(2\mu\alpha)^{-1/2}$.
    ``Offsets'' lists the values of $m-n$ for which $X^{(k)}_{mn}\neq 0$.%
  }
  \label{tab:Xk_bandwidth}
  \ra{1.3}
  \begin{tabular}{ccc}
    \toprule
    $k$ & Non-zero offsets $m-n$ & Prefactor \\
    \midrule
    1 & $\pm 1$ & $\ell$ \\
    2 & $0, \pm 2$ & $\ell^2$ \\
    3 & $\pm 1, \pm 3$ & $\ell^3$ \\
    4 & $0, \pm 2, \pm 4$ & $\ell^4$ \\
    5 & $\pm 1, \pm 3, \pm 5$ & $\ell^5$ \\
    6 & $0, \pm 2, \pm 4, \pm 6$ & $\ell^6$ \\
    \bottomrule
  \end{tabular}
\end{table}

% ============================================================
\subsection{Optimal width parameter \texorpdfstring{$\alpha$}{alpha}}\label{sec:opt_alpha}
% ============================================================

The convergence of the variational expansion depends critically on the choice
of the HO width parameter $\alpha$.  The stationarity condition
$\partial E_0/\partial\alpha = 0$ via the Hellmann--Feynman theorem gives the
implicit equation
\begin{equation}\label{eq:opt_alpha_general}
  \langle\hat T\rangle_\alpha = \frac{1}{2}\sum_{k=0}^{N} k\,c_k\,
  \langle\hat x^k\rangle_\alpha,
\end{equation}
where all expectation values are taken in the ground eigenstate at width
$\alpha$.  This must be solved self-consistently.

\subsubsection{Analytical estimate for a pure monomial}

For the pure monomial $V(x)=c_N x^N$ ($N$ even), the virial theorem gives
a closed-form estimate.  Setting the kinetic and potential energy scales
equal yields
\begin{equation}\label{eq:opt_alpha_mono}
  \alpha_{\mathrm{opt}}^{(\mathrm{mono})}
  = \left[\frac{N c_N\,(N-1)!!}{2\left(2\mu\right)^{N/2}}\right]^{2/(N+2)},
\end{equation}
where $(N-1)!! = 1\cdot 3\cdot 5\cdots(N-1)$ is the double factorial.
For $N=4$ (quartic double-well at the barrier) this recovers the standard
result $\alpha_{\mathrm{opt}}=(3c_4/\mu)^{1/3}$ up to numerical factors.

\subsubsection{Practical recommendation}

For production runs, fix $\alpha$ by the local harmonic frequency at the
minimum of $V(x)$.  If the potential has minima at $x_{\pm}$, set
\begin{equation}\label{eq:alpha_harm}
  \alpha = \mu\,\omega_{\min},
  \qquad
  \omega_{\min} = \sqrt{\frac{V''(x_+)}{\mu}}.
\end{equation}
For the symmetric double-well this reduces to $\alpha=\mu\omega_0$ with
$\omega_0 = \sqrt{4b}$ (see Section~\ref{sec:potential}).

% ============================================================
\subsection{Parity decomposition for symmetric potentials}\label{sec:parity_gen}
% ============================================================

When all odd-index coefficients vanish ($c_1=c_3=\cdots=0$) the potential is
even: $V(-x)=V(x)$.  The Hamiltonian then commutes with the parity operator
$\hat\Pi$, and the eigenstates are either even ($\phi_n$ with $n$ even) or
odd ($\phi_n$ with $n$ odd).  The Hamiltonian matrix block-diagonalises into
two independent $(M/2)\times(M/2)$ sub-matrices:
\begin{equation}
  \mathbf{H}
  = \mathbf{H}^{(+)} \oplus \mathbf{H}^{(-)},
\end{equation}
where $\mathbf{H}^{(+)}$ acts on even-index basis functions and
$\mathbf{H}^{(-)}$ on odd-index ones.  Exploiting this structure reduces the
diagonalisation cost by a factor of $\sim 4$.

% ============================================================
\subsection{Number of basis functions and convergence}\label{sec:conv_gen}
% ============================================================

The number of basis functions $M$ required for a given accuracy depends on
two competing factors:
\begin{enumerate}
  \item \textbf{Energy range:} All HO basis functions with energies
        $E_n^{(\mathrm{HO})}\lesssim E_{\mathrm{max}}$ should be included,
        where $E_{\mathrm{max}}$ is the highest eigenvalue of interest.
  \item \textbf{Coupling range:} The bandwidth $N$ means each function couples
        to its $N$-th nearest neighbours; including too few functions truncates
        off-diagonal couplings.
\end{enumerate}

A practical criterion: choose $M$ such that
\begin{equation}\label{eq:M_criterion}
  E_{M-1}^{(\mathrm{HO})} = \alpha\hb(M-\tfrac{1}{2}) \gg V_{\mathrm{max}},
\end{equation}
where $V_{\mathrm{max}}$ is the energy of the highest eigenstate to be
resolved.  Convergence should always be verified by increasing $M$ until the
eigenvalues of interest change by less than the desired tolerance.

% ============================================================
\subsection{General wavepacket construction}\label{sec:gen_wp}
% ============================================================

The general initial state $\ket{\Psi(0)}$ is specified as a superposition of
the first $M$ eigenstates $\{\ket{\psi_j}\}_{j=0}^{M-1}$:
\begin{equation}\label{eq:gen_wp}
  \ket{\Psi(0)} = \sum_{j=0}^{M-1} c_j\,\ket{\psi_j},
  \qquad
  \sum_{j=0}^{M-1}|c_j|^2 = 1.
\end{equation}
Three natural construction modes are supported by \texttt{QuTu}.

\subsubsection{Mode~1: localized Gaussian}\label{sec:wp_gaussian}

Centered at position $x_0$ with momentum $p_0$ and width $\sigma$:
\begin{equation}\label{eq:wp_gauss}
  \Psi(x,0) = \mathcal{N}\exp\!\left(-\frac{(x-x_0)^2}{4\sigma^2}
              + i p_0 x\right), \qquad
  \mathcal{N} = (2\pi\sigma^2)^{-1/4}.
\end{equation}
The expansion coefficients are
\begin{equation}\label{eq:cj_gauss}
  c_j = \braket{\psi_j}{\Psi(0)}
      = \int_{-\infty}^{\infty} \psi_j^*(x)\,\Psi(x,0)\,dx,
\end{equation}
evaluated numerically on a grid fine enough to resolve both $\psi_j$ and
$\Psi(x,0)$.

\subsubsection{Mode~2: energy-filtered state}\label{sec:wp_energy}

Populate all eigenstates in an energy window
$[E_{\mathrm{low}},E_{\mathrm{high}}]$ with equal weight:
\begin{equation}\label{eq:wp_energy}
  c_j =
  \begin{cases}
    \mathcal{N}^{-1} & \text{if } E_{\mathrm{low}} \leq E_j \leq E_{\mathrm{high}},\\
    0 & \text{otherwise,}
  \end{cases}
  \quad
  \mathcal{N}^{2} = \#\{j : E_{\mathrm{low}}\leq E_j\leq E_{\mathrm{high}}\}.
\end{equation}
This creates a \emph{quasi-classical} packet that samples a controlled energy
range without presupposing spatial localisation.

\subsubsection{Mode~3: user-defined amplitudes}\label{sec:wp_user}

The user specifies the complex amplitudes $\{c_j\}$ directly.  The code
renormalises: $c_j \leftarrow c_j/\|\mathbf{c}\|$.  This mode is useful for
initialising the two-state packet $c_0=c_1=1/\sqrt{2}$ (symmetric
tunneling) or the four-state packet of Section~\ref{sec:wavepackets}.

% ============================================================
\subsection{Time evolution and observables}\label{sec:gen_time_evo}
% ============================================================

Given the general expansion~\eqref{eq:gen_wp}, the time-evolved state is
\begin{equation}\label{eq:gen_psi_t}
  \ket{\Psi(t)} = \sum_{j=0}^{M-1} c_j\,e^{-iE_j t/\hb}\,\ket{\psi_j}.
\end{equation}
All observables of Section~\ref{sec:observables} carry over without
modification: survival probability, position expectation value, tunneling
splitting, and turning points are each determined by the same formulas with
$\{E_j,\psi_j,c_j\}$ now drawn from the general diagonalisation.

\begin{remark}
  The tunneling splitting $\Delta E=E_1-E_0$ and half-period
  $T_{\mathrm{tun}}=h/\Delta E$ remain meaningful for any symmetric double-well
  regardless of the polynomial degree $N$, as long as $E_0$ and $E_1$ are
  below the barrier.
\end{remark}

% ============================================================
\subsection{Implementation summary}\label{sec:gen_impl}
% ============================================================

Table~\ref{tab:gen_impl} summarises the algorithmic flow for the general
polynomial case as implemented in \texttt{QuTu}.

\begin{table}[htbp]
  \centering
  \caption{%
    Step-by-step algorithm for the general polynomial potential solver.
    All matrix operations are $O(M^2)$ in memory; the dense \texttt{dsyev}
    eigensolver is $O(M^3)$.  A banded \texttt{dsbev} path is planned.%
  }
  \label{tab:gen_impl}
  \ra{1.3}
  \begin{tabular}{clp{7cm}}
    \toprule
    Step & Operation & Notes \\
    \midrule
    1 & Read $\{c_k\}$, $N$, $M$, $\mu$, $\alpha$ &
        User input; $\alpha$ from Eq.~\eqref{eq:alpha_harm} if not given \\
    2 & Compute $X^{(k)}$ for $k=0,\ldots,N$ &
        Master recursion~\eqref{eq:master_rec}, $O(kM^2)$ per $k$ \\
    3 & Assemble $V_{mn} = \sum_k c_k X^{(k)}_{mn}$ &
        Linear combination, $O(NM^2)$ \\
    4 & Add $T_{mn}$; form $H_{mn}$ &
        Analytic kinetic matrix, Section~\ref{sec:HO} \\
    5 & Diagonalise $\mathbf{H}$ &
        LAPACK \texttt{dsyev} (symmetric dense); banded \texttt{dsbev} planned \\
    6 & Construct initial state $\ket{\Psi(0)}$ &
        One of the three modes, Section~\ref{sec:gen_wp} \\
    7 & Propagate: $c_j(t)=c_j e^{-iE_j t/\hb}$ &
        Analytic; $O(M)$ per time step \\
    8 & Compute observables &
        Section~\ref{sec:observables} and~\ref{sec:additional_obs} \\
    \bottomrule
  \end{tabular}
\end{table}


% ============================================================
%                       References
% ============================================================
% \clearpage
% \printbibliography[title={References}]

\end{document}
%%%%%%%%%%%%%%%%%%%%%%%%%%%%%%%%%%%%%%%%%%%%%%%%%%%%%%%%%%%%%%%%%%%%%%%%%%%%%%%
